\documentclass[a4paper,12pt]{article}
\usepackage[utf8]{inputenc}
\usepackage{titlesec}
\usepackage{stmaryrd}
 
\titleformat
{\chapter} % command
[display] % shape
{\bfseries\Large\itshape} % format
{Chapter \ \thechapter} % label
{0.5ex} % sep
{
    \rule{\textwidth}{1pt}
    \vspace{1ex}
    \centering
} % before-code
[
\vspace{-0.5ex}%
\rule{\textwidth}{0.3pt}
] 
%%%%%%%%%%%%%%%%%%%%%%%%%%%%%%%%%%%%%%%%%%%%%%%%%%%%%%
%AGGIUNGO TEOREMA LEMMA PROPOSIZIONE COROLLARIO PROVA%
%%%%%%%%%%%%%%%%%%%%%%%%%%%%%%%%%%%%%%%%%%%%%%%%%%%%%%

\newtheorem{theorem}{Theorem}[section]
\newtheorem{lemma}[theorem]{Lemma}
\newtheorem{proposition}[theorem]{Proposition}
\newtheorem{corollary}[theorem]{Corollary}
\newtheorem{definition}[theorem]{Definition}

\newenvironment{proof}[1][Proof]{\begin{trivlist}
\item[\hskip \labelsep {\bfseries #1}]}{\end{trivlist}}
\newenvironment{example}[1][Example]{\begin{trivlist}
\item[\hskip \labelsep {\bfseries #1}]}{\end{trivlist}}
\newenvironment{remark}[1][Remark]{\begin{trivlist}
\item[\hskip \labelsep {\bfseries #1}]}{\end{trivlist}}




\bibliographystyle{plain}% the recommended bibstyle


\usepackage{
amssymb,
color,
amsmath,
proof
}
\usepackage{graphicx}

%%%%%%%%%%%%%%%%%%%%%%%%%%%%%%%%%%%%%%%%%%%%%%%%
% Author macros 
%%%%%%%%%%%%%%%%%%%%%%%%%%%%%%%%%%%%%%%%%%%%%%%%

\newcommand{\setA}              { {\mathcal A} }
\newcommand{\setB}              { {\mathcal B} }
\newcommand{\setC}              { {\mathcal C} }
\newcommand{\setD}              { {\mathcal D} }
\newcommand{\setE}              { {\mathcal E} }
\newcommand{\setF}              { {\mathcal F} }
\newcommand{\setG}              { {\mathcal G} }
\newcommand{\setH}              { {\mathcal H} }
\newcommand{\setI}              { {\mathcal I} }
\newcommand{\setJ}              { {\mathcal J} }
\newcommand{\setK}              { {\mathcal K} }
\newcommand{\setL}              { {\mathcal L} }
\newcommand{\setM}              { {\mathcal M} }
\newcommand{\setN}              { {\mathcal N} }
\newcommand{\setO}              { {\mathcal O} }
\newcommand{\setP}              { {\mathcal P} }
\newcommand{\setR}              { {\mathcal R} }
\newcommand{\setS}              { {\mathcal S} }
\newcommand{\setT}              { {\mathcal T} }
\newcommand{\setU}              { {\mathcal U} }
\newcommand{\setW}              { {\mathcal W} }
\newcommand{\setX}              { {\mathcal X} }
\newcommand{\setY}              { {\mathcal Y} }
\newcommand{\setZ}              { {\mathcal Z} }

\newcommand{\Nat}               { {\tt N}  }
\newcommand{\nth}               { {\tt nth}  }

\newcommand{\Power}             { {\setP}  }
\newcommand{\Pfin}              { {\setP_{\mbox{\tiny fin}}}  }

\newcommand{\Prop}              { {\tt Prop}  }
\newcommand{\LA}                { {\tt LA} }
\newcommand{\HA}                { {\tt HA}}
\newcommand{\PA}                { {\tt PA}}

\newcommand{\ord}               { {\tt ord} }
\newcommand{\false}             { {\tt false} }
\newcommand{\False}             { {\tt False} }
\newcommand{\true }             { {\tt true } }
\newcommand{\True }             { {\tt True } }

\newcommand{\universe}  [1]     { {|#1|} }
\newcommand{\Universe}  [1]     { {|\!|#1|\!|} }
\newcommand{\interpret} [1]     { {[\![#1]\!]} }
\newcommand{\ON}                { {\tt ON} }
\newcommand{\WF}                { {\tt WF} }
\newcommand{\Dil}               { {\tt Dil} }

\newcommand{\Nicla}             { {\mathcal N} }

\newcommand{\Def}               { {\tt Def} }
\newcommand{\Par}               { {\tt Par} }
\newcommand{\Gen}               { {\tt Gen} }
\newcommand{\supremum}          { {\bigvee} }
\newcommand{\Analytical}        { {\setA} }
\newcommand{\Rel}               { {\tt Rel} }
\newcommand{\length}            { {\tt l} }
\newcommand{\norm}              { {\tt norm} }
\newcommand{\NF}                { {\tt NF} }
\newcommand{\Main}              { {\tt Main} }
\newcommand{\subexpr}           { {\prec} }
\newcommand{\subtree}           { {\prec} }
\newcommand{\subexpreq}         { {\preceq} }
\newcommand{\subtreeeq}         { {\preceq} }
\newcommand{\Acc}               { {\tt Acc} }
%\newcommand{\AccRec}            { {\tt AccRec} }
\newcommand{\Rec}               { {\tt R} }
\newcommand{\rec}               { {\tt r} }
%\newcommand{\constr}            { {\implies_{\!\!\!\!\!\theta}} }
\newcommand{\equal}             { {\sim} }
\newcommand{\notequal}         { {\not \sim} }
\newcommand{\identical}       { {=} }
\newcommand{\notidentical}    { {\not =} }
\newcommand{\CUT}             { {\tt CUT} }
\newcommand{\cut}             { {\tt cut} }
\newcommand{\Exists}          { {\bigvee} }
\newcommand{\Forall}          { {\bigwedge} }
\newcommand{\inac}    [2]     { {{#1}^{+#2}} }
\newcommand{\subinac} [2]     { {{#1}^{-#2}} }
\newcommand{\join}    [3]     { {{\tt j}^{#1}_{#2} {#3} } }
\newcommand{\bigjoin}    [3]     { {{\tt J}^{#1}_{#2} {#3} } }
\newcommand{\Clos}            { {\tt Clos} }
\newcommand{\Backward}         { {\tt B} }
\newcommand{\Forward}         { {\tt F} }
\newcommand{\List}            { {\tt List} }
\newcommand{\FV}              { {\tt FV} }
%\newcommand{\Lift}     [2]    { {(#1)^{*#2}} }
\newcommand{\Lift}            { {\tt lift} }
\newcommand{\Deg}             { {\tt Deg} }
\newcommand{\fdeg}            { {\tt fdeg} }
\newcommand{\Occ}             { {\tt Occ} }
\newcommand{\Proof}           { {\tt Pr} }
\newcommand{\id}              { {\tt id} }
\newcommand{\ID}              { {\tt ID} }
%\newcommand{\comp}            { {{}^{\mbox{\tiny $\circ$}}} }
\newcommand{\comp}            { {\circ} }
\newcommand{\Br}              { {\tt Br} }
%\newcommand{\Win}             { {\tt val} }
\newcommand{\dom}             { {\tt dom} }
\newcommand{\cod}           { {\tt cod} }
\newcommand{\Cod}           { {\tt Cod} }
\newcommand{\rel}             { {\tt rel} }
\newcommand{\cons}            { {\tt cons} }
\newcommand{\Cons}            { {\tt Cons} }
%\newcommand{\cons}            { {\Forarrow} }
\newcommand{\Cone}            { {\tt Cone} }
\newcommand{\Pow}             { {\tt Pow} }
\newcommand{\Fix}             { {\tt Fix} }
\newcommand{\incomplete}      { {\tt incomplete} }
\newcommand{\ordinals}        { {\tt ord} }
\newcommand{\Form}            { {\tt Form} }
\newcommand{\Ext}             { {\tt Ext} }
\newcommand{\conc}            { {*} }
\newcommand{\leftsymb}        { {\tt l} }
\newcommand{\rightsymb}       { {\tt r} }
\newcommand{\Valid}           { {\tt V} }
\newcommand{\Winning}         { {\tt W} }
\newcommand{\Win}        [1]  { {\models_{#1}} }
\newcommand{\numeral}    [1]  { {\underline{#1}} }
\newcommand{\var}             { {\tt var} }
\newcommand{\Var}             { {\tt Var} }
\newcommand{\Const}           { {\tt Ax} }
\newcommand{\eql}             { {=_\leftsymb} }
\newcommand{\eqr}             { {=_\rightsymb} }
\newcommand{\Cl}              { {\Const_\leftsymb} }
\newcommand{\Cr}              { {\Const_\rightsymb} }
\newcommand{\contr}           { {{\tt contr}} }
\newcommand{\contrl}          { {{\tt c}_\leftsymb} }
\newcommand{\contrr}          { {{\tt c}_\rightsymb} }
\newcommand{\Vl}              { {\Var_\leftsymb} }
\newcommand{\Vr}              { {\Var_\rightsymb} }
\newcommand{\impliesl}        { {\implies_{\!\!\!\!\!\leftsymb}} }
\newcommand{\impliesr}        { {\implies_{\!\!\!\!\!\rightsymb}} }
\newcommand{\foralll}         { {\forall_\leftsymb} }
\newcommand{\forallr}         { {\forall_\rightsymb} }
\newcommand{\Tree}            { {\tt Tree} }
\newcommand{\FTree}           { {\tt FTree} }
\newcommand{\Set}             { {\tt Set} }
\newcommand{\Poly}            { {\tt P} }
\newcommand{\ap}              { {\tt ap} }
\newcommand{\poly}            { {\tt poly} }
\newcommand{\Skeleton}        { {\tt Ske} }
\newcommand{\Label}             { {\tt l} }
\newcommand{\head }           { {\tt head} }
\newcommand{\tail }           { {\tt tail} }
\newcommand{\Tail }           { {\tt Tail} }
\newcommand{\Eq}              { {\tt Eq} }
\newcommand{\Hom}             { {\tt Hom} }
\newcommand{\inj}             { {\tt j} }
\newcommand{\EM}              { {\tt EM} }
%\newcommand{\restr}           { {\lceil} }
\newcommand{\Ind}            { {\tt Ind} }
\newcommand{\ZF}              { {\tt ZF} }
\newcommand{\Conv}            { {\tt Conv} }
\newcommand{\eq}              { {\tt eq} }
\newcommand{\casedef} [3]     { { {\tt if} \ (#1)
                                \ {\tt then} \ #2 \ {\tt else} \ #3} }
\newcommand{\Unit}            { {\tt Unit} }
\newcommand{\unit}            { {\tt unit} }
\newcommand{\Bool}            { {\tt Bool} }
\newcommand{\lexicographic}   { {  <_{\mbox{\tiny lex}}}  }
\newcommand{\lexicographiceq} { {\le_{\mbox{\tiny lex}}}  }
\newcommand{\KleeneBrouwer}   { {  <_{\mbox{\tiny KB}}} }
\newcommand{\KleeneBrouwereq} { {\le_{\mbox{\tiny KB}}} }
\newcommand{\HypothesisLine}  { {\tt HL} }
\newcommand{\NestedHypothesisLine}  { {\tt NHL} }
\newcommand{\Bck}             { {\tt Bck} }
\newcommand{\hypothesisline}  { {\tt hl} }
\newcommand{\JUST}            { {\tt J} }
\newcommand{\changes}         { {\tt ch} }
\newcommand{\FConv}           { {\tt FConv} }
\newcommand{\RejectionLine}   { {\tt RL} }
\newcommand{\conv}            { {\Downarrow} }
\newcommand{\stationary}      { {>} }
%\newcommand{\stationaryeq}    { {\ge} }
\newcommand{\strict}          { {>} }
\newcommand{\Incr}            { {\tt Incr} }
\newcommand{\Segm}            { {\tt Sg} }
\newcommand{\Val}             { {\tt V} }
\newcommand{\val}             { {\tt val} }
\newcommand{\canonical}       { {\tt j} }
\newcommand{\constr}          { {\tt cons} }
\newcommand{\Max}             { {\tt Max} }
\newcommand{\Fun}             { {\tt Fun} }
\newcommand{\reduces}         { {\leadsto} }
%\newcommand{\reduces}         { {\rightarrowtail  \!\!\!\!\!\! \rightarrowtail } }
\newcommand{\trunk}           { {\tt tr} }
\newcommand{\Real}            { {\tt \bf R} }
\newcommand{\Hyper}           { {\setH} }
\newcommand{\Ctxt}            { {\tt Ctxt} }
\newcommand{\Auxiliary}       { {\setA} }
\newcommand{\case}            { {\tt case} }
\newcommand{\segment}  [2]    { {\universe{#1}_{#2}} }
\newcommand{\Twice}           { {\tt Tw} }
\newcommand{\Red}             { {\tt Red} }
\newcommand{\Term}            { {\tt Term} }
\newcommand{\Player}          { {\tt P} }
\newcommand{\Opponent}        { {\tt O} }
\newcommand{\turno}           { {\tt t} }
\newcommand{\vince}           { {\tt w} }
\newcommand{\Foglie}          { {\tt Lv} }
\newcommand{\nil}             { {\tt nil} }
\newcommand{\sub}             { {\tt sub} }
\newcommand{\STOP}            { {\tt STOP} }
\newcommand{\UNFOLD}          { {\tt UNFOLD} }
\newcommand{\DROP}            { {\tt DROP} }
\newcommand{\Drop}            { {\tt D} }
\newcommand{\END}             { {\tt END} }
\newcommand{\EXH}             { {\tt EXH} }
\newcommand{\exh}             { {\tt exh} }
\newcommand{\REPEAT}          { {\tt repeat} }
\newcommand{\ignore}          { {\tt ignore} }
\newcommand{\tryall}          { {\tt tryall} }
\newcommand{\ch}              { {\tt ch} }
\newcommand{\Ch}              { {\tt Ch} }
\newcommand{\Atom}            { {\tt At} }
\newcommand{\range}           { {\tt range} }
\newcommand{\Mosse}           { {\tt Mv} }
\newcommand{\EXISTS}          { {\exists \!\!\! \exists} }
\newcommand{\FORALL}          { {\forall \!\!\! \forall} }
\newcommand{\CA}              { {\tt CA} }
\newcommand{\Wit}             { {\tt Wit} }
\newcommand{\EXISTSGAME}      { {\bigvee} }
\newcommand{\FORALLGAME}      { {\bigwedge} }
\newcommand{\Root}            { {\tt root} }
\newcommand{\cof}             { {\tt cof} }
\newcommand{\const}           { {\tt c} }
\newcommand{\Back}            { {\tt B} }
\newcommand{\Forth}           { {\tt F} }
\newcommand{\forth}           { {\tt f} }
\newcommand{\Zeta}            { {\tt Z} }
%\newcommand{\merge}           { {\tt m} }
\newcommand{\Leaf}            { {\tt 0} }
\newcommand{\reduction}  [2]  { {({#1}|{#2})} }
\newcommand{\Delay}           { {\tt D} }
\newcommand{\Type}            { {\tt Tp} }
\newcommand{\dummy}           { {\tt dummy} }
\newcommand{\Tot}           { {\tt T} }
\newcommand{\Bot}             { {\bot \!\!\!\! \bot} }
\newcommand{\Data}            { {\tt D} }
\newcommand{\restr}           { {\preceq} }
\newcommand{\Th}              { {\tt Th} }
\newcommand{\axiom}           { {\tt ax} }
\newcommand{\Exch}            { {\tt Exch} }
\newcommand{\Cand}            { {\tt Cand} }
\newcommand{\Index}           { {\tt I} }
\newcommand{\realizes}        { {\models} }
\newcommand{\uniform}         { {\tt u} }
\newcommand{\future}          { {\tt f} }
\newcommand{\CMod}            { {\tt CMod} }
\newcommand{\Mod}             { {\tt Mod} }
\newcommand{\smallest}      { {\tt mod} }
\newcommand{\Iset}            { {\tt I} }
\newcommand{\Pos}            { {\tt Pos} }
\newcommand{\Proofs}        { {\tt P} }
\newcommand{\PER}           { {\tt \sim} }
\newcommand{\Nu}             { {\mathcal V} }
\newcommand{\SET}           { {\tt SET} }
\newcommand{\HC}             { {{\setH}_1} }
\newcommand{\Id}               { {\tt Id} }
\newcommand{\Dom}               { {\tt Dom} }
\newcommand{\Range}               { {\tt Range} }
\newcommand{\Supp}          { {\tt Supp}}
\newcommand{\Inv}          { {\tt Inv}}
\newcommand{\Comp}         {  {\odot} } % {  {\otimes} }
\newcommand{\PreSupp}    { {\tt S} }
\newcommand{\eval}    { {\tt eval} }
\newcommand{\fun}    { {\tt fun} }
\newcommand{\hy}        { {\operatorname{-}}}
\newcommand{\eqL}     {= \! \!  \!  L}
\newcommand{\eqR}    {= \! \!  \!  R}
\newcommand{\Ack}   {{\tt A}}
\newcommand{\LKID} { {\tt LKID} }
\newcommand{\CLKID} { {\tt CLKID} }
\newcommand{\Weak} { {\tt Weak} }

%%%%%%%%%%%%%%%%%%%%%%%%%%%%%
%
%
%         FINE MACRO
%
%
%%%%%%%%%%%%%%%%%%%%%%%%%%%%%



\begin{document}

\title{Negative and Positive Results about Cut-elimination for Cyclic Proofs (DRAFT)}

\author{S. Berardi, M. Tatsuta}


\maketitle

{\bf Abstract}.
In this paper we  give four examples of cyclic proofs having no cut-free form, for different reasons. We prove that cut-free  cyclic proofs define algorithms requiring polynomial time and constant space, and therefore only exist for simple theorems.  We found a common syntactical pattern to most failures of cut-elimination: the cuts include what we call \emph{local} cuts. On the positive side, we prove cut-elimination for cyclic proofs having only non-local cuts. 



\section{Introduction}
We consider $\LKID$,  Martin-L\"{o}f theory of inductive definitions plus classical logic (\cite{FOSSACS2017}), and an alternative proof system for the same language, cyclic proofs (\cite{Brotherston2006}), with no explicit induction rule. 

In this paper, we prove that cut-elimination for a formula $A$ in a cyclic proof may fail for many reasons: if $A$ includes quantifiers; if $A$ includes some inductive predicate $p$ whose productions discard some free variable in the conclusion; if our goal is proving the existence of an individual whose size is exponential in the parameters of the problem; and if our goal requires an algorithm running in non-constant space. This is the case of proving that two definitions of transitive closure for a binary relation are equivalent. 

In all cases (except perhaps the first one) the cuts we could not eliminate are what we call \emph{local} cuts. We give an example of a local cut in the case of the inductive predicate $N$, ``being a natural number''. Assume there is some cut between $\Gamma \vdash \Delta, N(Su)$ and $N(Su), \Gamma \vdash \Delta$. Assume we proved the first $N(Su)$ from $N(u)$ and the second $N(Su)$ from $Su=Sx, N(x)$ for some eigenvariable $x$ of the proof. If $x \in \FV(u)$ we call the cut  a local cut. Here ``local'' means that the value $u$ assigned by the cut to the eigenvariable $x$ depends from $x$ itself. If this is never the case we call the cut a global cut.

Local and global cuts may be generalized to cuts on any inductive predicate, on quantifier-free formulas, and to sets of cuts. Our positive result will be: if the set of all cuts of a proof is global then we may remove all cuts.

We explain informally why often we cannot eliminate a local cut. A local cut may destroy regularity, the fact that a cyclic proof has only finitely many subproofs and formulas. Indeed, the normalization process applies infinitely many times the substitution $[u/x]$, and if $x \in \FV(u)$ then the substitutions $[u/x]$, $[u/x][u/x]$, $[u/x][u/x][u/x]$, \ldots produce different formulas, hence infinitely many formulas, hence a non-cyclic proof. 

Global cuts are is simple, yet are not trivial. For instance, for any polynomial $P$ there is a canonical proof that $\Nat$ is closed under $P$. The set of cuts of this proof is global, hence all these cuts may be removed. For comparison, we prove that there is a cyclic proof that $\Nat$ is closed under $2^x$, but there is no cut-free cyclic proof of the same result.

%11:34 13/09/2018
%An induction rule may be translated by a cut-free proof when it uses no left-rule on the assumption $F[x]$ of the proof of the inductive step $\Gamma, F[x] \vdash F[Sx], \Delta$.



%We call \emph{infinitary binary proof-tree}, an infinitary proof for short, any proof tree made with rules of arity $0, 1, 2$ and possibly infinite. Finite proofs (for instance, all first order logical proofs) are included as a trivial case. A cyclic proof $\Pi$ is an infinitary binary proof-tree which is regular, that is, is a possibly infinite tree with finitely many different subtrees, and which satisfies the global trace condition (taking the place of the induction rule).  The global trace condition says that in every infinite branch in the proof tree there is some infinite sequence of inductively defined predicates $p(\vec{t})$, $p'(\vec{t'})$, $p''(\vec{t''})$,\ldots, with $p'(\vec{t})$ assumption of $p(\vec{t'})$, $p''(\vec{t''})$ assumption of $p'(\vec{t'})$ and so forth. The global trace condition may be used to show that if all proof-rules are correct then the conclusion of the infinitary proof is true. 

\section{Cut-elimination fails for cyclic proofs and quantified formulas}
In this section we prove that cut-elimination fails for cyclic proofs if the cut include at least one quantifier. This result has the advantage of having a short proof, but it leaves open the possibility of cut-elimination for quantifier-free formulas. In the next sections, using a fine analysis of the properties of cut-free proofs, we will show that even for quantifier-free formulas we have cut-elimination only in particular cases.

We first introduce some notations. We call $\CLKID^\omega$ the set of cyclic theorems. $\CLKID^\omega$ is a non-conservative extensions of $\LKID$ (\cite{FOSSACS2017}). If we add some injectivity axioms for constructors to both theories they become equal (\cite{Simpson2017}, \cite{LICS2017}). From now on, let $I\hy \Sigma^0_n$, $I\hy \Pi^0_n$ be the fragments of arithmetic with $\Sigma^0_n\hy $induction and $C\hy \Sigma^0_n$, $C\hy \Pi^0_n$ be the fragments of cyclic proofs whose formulas are all $\Sigma^0_n$, $\Pi^0_n$. The sub-formula property for a theory is: all formulas occurring in a proof are subformula of the conclusion or of some assumption.

This paper is about cut-elimination and the sub-formula property. Arithmetical proofs have cut-elimination, yet a cut-free arithmetical proof has no sub-formula property when an induction axiom is applied to an open term. Indeed, let $\Cons_n$ be the consistency statement for $I\hy \Sigma^0_n$, expressed as a $\Delta^0_0$-open statement. Then $\Cons_n$ is provable in arithmetic, yet it has no proof with the subformula property, otherwise this would be a proof in $I\hy \Sigma^0_0 \subseteq I\hy \Sigma^0_n$, contradicting G\"{o}del's Incompleteness Theorem. 

Cut-free cyclic $\Sigma^0_n$-proof have a weak form of sub-formula property. Indeed, since we have a left-rule for elimination of equality, whenever we have formulas $A[a,b]$ and $a=b$ in the proof we may also have $A[b,a]$. Due to the case rule, if we have a formula $\Nat t$ in the proof, then we may have $t = Sx$ in the same proof. Due to this weak sub-formula property, the argument against cut-elimination for arithmetical proofs may be adapted to show that cyclic proofs with quantifiers and whose only inductive predicate is $\Nat$ have no cut-elimination in general. For cyclic proofs, the proof of this result is more involved than for arithmetical proofs. 
%The reason is that cyclic proofs have no explicit induction axiom, hence their lack of cut-elimination is less evident than for arithmetic. 

\begin{theorem}
Cut-elimination fails for $C\hy \Sigma^0_{1}$ (for cyclic proofs with sequents made of $\Sigma^0_1$-formulas).
\end{theorem}

\begin{proof}
We use the following result (\cite{Anupam}): for all $n \in \Nat$ the theories $I\hy \Sigma^0_{n+1}$ and $C\hy \Sigma^0_{n}$ prove the same $\Pi^0_{n+1}$-theorems. Let $F = \Cons_1 \in \Delta^0_1$ be the consistency statement for  $I\hy \Sigma^0_1$, expressed as an open formula. $\Cons_1$ is provable in $I\hy \Pi^0_2$.  Indeed, truth for a $\Sigma^0_1\hy $open formula is a $\Pi^0_2\hy $statement. Thus, we may prove by $\Pi^0_2\hy $induction that all proofs using only $\Sigma^0_1\hy $sequents are true. Then we use free-cut elimination, which reduces all formulas to $\Sigma^0_1\hy $formula, in order to prove that all $I\hy \Sigma^0_1\hy $proofs with a $\Sigma^0_2\hy $conclusion are true, hence are not a proof of false.  From $I\hy \Pi^0_2 \vdash F$ we obtain $I\hy \Sigma^0_2 \vdash F$ by \cite{Buss}. By \cite{Anupam} for $n=1$ and $F \in \Delta^0_1$ we have $C\hy \Sigma^0_1 \vdash F$, with a proof $\Pi$ possibly using cuts. We prove that $\Pi$ has no cut-free form. For contradiction, take any cut-free proof of $C\hy \Sigma^0_1 \vdash G$. By the subformula property and $F \in \Delta^0_1$ this proof is quantifier-free, hence it is a proof in $C\hy \Sigma^0_0 \vdash F$. By  \cite{Anupam} for $n=0$ and $F \in \Pi^0_1$ we have $I\hy \Sigma^0_1 \vdash F$. By G\"{o}del's Incompleteness Theorem we conclude $I\hy \Sigma^0_1 \not \vdash F$, which is the required contradiction. 
$\Box$ \end{proof}

%16:26 04/07.2018

%This means that if we are looking for a fragment of cyclic proofs with cut-elimination we should consider quantifier-free cuts. But this is not enough: in this draft we give several examples of connective-free (hence quantifier-free) proofs without cut-elimination. In order to find a fragment with cut-elimination we will restrict ourselves to a proper subset of of cyclic proofs with quantifier-free cuts corresponing to reasoning on maps with a polynomial growth. 

As we said, cut-elimination for quantifiers is ruled out by the previous theorem. If we have all inductive definitions then inductive predicates are exactly the $\Sigma^0_1$-predicates (see next section), and quantifier-free formulas are a subset of $\Delta^0_2$-formulas. Thus, cut-elimination restricted to quantifier-free formulas is not trivial, and in general it fails. However, we could define a set of quantifier-free cuts which we called \emph{global cuts} and which may be eliminated. 

% APPARENTLY IT IS NO MORE REQUIRED 12:10 13/09/2018
%We call global inductive definitions all inductive definitions introducing no fresh variable: we require that for each production the free variables of all premises are included in the free variables on the conclusion. 

%We give one more detail about our cut-elimination procedure for cyclic proofs.




%Assume $A, B$ are quantifier-free formulas, and write $A \sim B$ when $A, B$ are equal if we replace all terms in $A, B$ with a dummy symbol: ``$\_$''. We denote with $A=B$ the list of equational assumptions required to prove that terms in $A, B$ are ordinately equal, provided $A \sim B$. We skip identities. For instance, $p(a)=p(b)$ is $a=b$ and $p(a)=p(a)$ is empty. If $p \not = q$, then $p(a) = q(b)$ is not defined because $p(\_), q(\_)$ are different. We define a generalized cut-rule $\nu$ for $A \sim B \sim C$: from $\Gamma_1 \vdash \Delta_1,A$ and $B, \Gamma_2 \vdash \Delta_2$ we deduce  $A=C, B=C, \Gamma_1, \Gamma_2 \vdash \Delta_1, \Delta_2$. Generalized cut-rule is derivable from cut-rule. In the case $A, B, C$ are the same we will deduce $A=A, A=A, \Gamma_1, \Gamma_2 \vdash \Delta_1, \Delta_2$. $A = A$ is empty, therefore the conclusion is  $\Gamma_1, \Gamma_2 \vdash \Delta_1, \Delta_2$: in this case $\nu$ is an ordinary cut.
%will be a list of identities and will be removed from the proof, without introducing new cuts. 

% 10:29 06/07/2018 ADESSO CERCO DI ELIMINARE ANCHE QUESTI TAGLI
% ALTRIMENTI LA PROCEDURA NON FUNZIONA
%Our cut-elimination procedure does not remove some ``trivial'' cuts. Indeed, it inserts one extra rule which we call the fresh variable rule: \emph{``if $x=u, \Gamma \vdash \Delta$ and $x \not \in \FV(u, \Gamma, \Delta)$ then $\Gamma \vdash \Delta$''}. This rule is equivalent to a cut with the canonical proof of $\vdash \exists x.x=u$. This means that our cut-elimination procedure will insert (the equivalent of) new cuts, even if a of a trivial sort. These new cuts are not eliminated.

%21:56 06/06/2018
%07/05/2018

\section{Failure of cut-elimination for quantifier-free proofs}
\label{section-failure}
In this section we prove that cut-elimination fails for atomic cuts and variable-discarding inductive definitions. 

An inductive definition is variable-discarding if it has some production having some free variable in the assumption not occurring in the conclusion of the production. An example is the inductive predicate $p$ with the production: if $q(x,y)$ then $p(x)$. A variable-discarding inductive definition hydes an existential quantifier inside the inductive predicate: in the previous example, $p(x)$ is equivalent to $\exists y.q(x,y)$. 

If we use a variable-discarding inductive definitions then cut-elimination fails even in the simplest cases. We formulate in a variable-discarding way totality the ``twice'' map $f(x)=2x$, a very weak arithmetical result. We prove that this statement has no cut-free proof. For comparison, we show that if we add cut to cyclic proofs we may prove the totality of $f(x)=2x$, and even the totality of the Ackermann function. This latter is a strong arithmetical result with no proof in $I \hy \Sigma^0_1$, yet it has a cyclic proof with no propositional connectives nor quantifiers.

\newcommand{\quotationMarks}  [1] {{``#1''}}

\subsection{The \quotationMarks{Twice} Map $f(x)=2x$}
In this sub-section we define totality for the map in $f(x)=2x$ using a variable-discarding inductive definition. If we state totality of $f$ in this way, we prove that totality of $f(x)$ has a cyclic proof with local cuts, but no cut-free total proof. 

We consider the language $L = \{0, S\}$. We first define the graph $a(x,y)$ of the map $f(x)=2x$ by: $a(0,0)$, and if $a(x,y)$ then $a(Sx,SSy)$. The totality of the map $f(x)=2x$ may be expressed by the predicate $\Tot(x)$ and the inductive rule: if $a(x,y)$ then $\Tot(x)$. This inductive definition discards the variable $y$: $y$ occurs in $a(x,y)$ bu does not occur in $\Tot(x)$. The interpreation of $\Tot(x)$ is $\exists y.a(x.y)$: local inductive definition allow to express $\exists$ inside an inductive predicate. 

There is a connective-free cyclic proof with cuts on $\Tot(y)$ of $\Nat x \vdash \Tot(x)$ in the language $0,S$, and no cut-free cyclic proof in $\CLKID^\omega$ of the same result. 
%Instead, there is a cut-free proof of $\forall x.x+0=x, \forall x,y. x+Sy=S(x+y), \Nat x \vdash \Tot(x)$: if we assume that we have a map implementing addition we may prove cut-free that $\ldots, \Nat x \vdash a(x,x+x)$, hence $\ldots, \Nat x \vdash \Tot(x)$. 

\begin{lemma}[The twice map is total]
\label{lemma-twice-total}
Let $t,u$ be any terms of the language $0,S$. 
There is a a connective-free cyclic proof with \emph{local} cuts on $\Tot(x)$ of $\Nat t \vdash \Tot(t)$ in the same language.
\end{lemma}


\begin{proof}
\begin{enumerate}
\item
We first define a finite proof $\Pi_0$ of $\Nat t \vdash \Tot(t)$ from the open assumption $(1) \Nat x \vdash \Tot(x)$ for any $x \not \in \FV(t)$. This is not yet the regular proof of $\Nat t \vdash \Tot(t)$.

We prove $\Nat t \vdash \Tot(t)$ by $\Nat L$ on $\Nat t$, from two premises: $t=0 \vdash \Tot(t)$ and $\Nat x, t=Sx, \Nat x \vdash \Tot(t)$. For each premise we provide a proof. 

\begin{enumerate}
\item
We prove $t=0 \vdash \Tot(t)$ by $\eqL$ from $\vdash \Tot(0)$, then by $\Tot R$-rule from: $\vdash a(0,0)$, which is an $a R$-axiom.

\item
Proof of $\Nat x, t=Sx, \Nat x \vdash \Tot(t)$. We use $\eqL$ with premise $\Nat x \vdash \Tot(Sx)$. Then we use a cut with $\Nat x \vdash \Tot(x)$ and $\Tot(x) \vdash \Tot(Sx)$.  $\Nat x \vdash \Tot(x)$ is an open assumption. We prove $\Tot(x) \vdash \Tot(Sx)$ by $\Tot L$-rule and $x=x', a(x',y') \vdash \Tot(Sx)$ (for $x',y'$ fresh), then by $\eqL$-rule and $a(x,y') \vdash \Tot(Sx)$.  We use $\Tot R$-rule with premise $a(x,y') \vdash a(Sx, SSy')$, then $a R$-rule with premise $a(x,y') \vdash a(x, y')$, and eventually $\id$.
\end{enumerate}
\item
We define a regular infinite proof $\Pi$ from $\Pi_0$ by repeating cyclically from  $\Nat x \vdash \Tot(x)$. Next time we use as eigenvariable some $y \not \equiv x$, then $x$ again and so forth. Thus, we may define a regular proof $\Pi$ by unfolding $\Pi_0$.
\item
We check that $\Nat x \vdash \Tot(x)$ and $\Tot(x) \vdash \Tot(Sx)$ is a local cut in $\Pi$. Among the ancestors of $T(x)$ (right-hand side) we have an occurrence of $T(Sx)$, hence of $a(Sx, SSy')$. Among the ancestors of  $T(x)$ (left-hand side) we have an occurrence of $a(x', y')$, which is the conclusion of a rule with eigenvariable $y'$.  
\item
We prove the global trace condition for $\Pi$. Take any infinite path $\pi$ of the proof $\Pi$: then this path passes infinitely many times through  $\Nat x \vdash \Tot(x)$ and  $\Nat y \vdash \Tot(y)$, progressing each time.
\end{enumerate}
$\Box$ \end{proof}
%13:46 13/09/2018
We may now prove that there is no cut-free proof of $\Nat x \vdash \Tot x$. We prove that if there is some cut-free proof then for some $k \in \Nat$ and all $x \in \Nat$ we would have $2x \le x+k$, that is, $x \le k$, which is impossible for $x=k+1$.

Given any cut-free $\Pi$ of $\Gamma \vdash \Delta$ we define some canonical operations on $\Pi$, which we use to extract information from $\Pi$. We assume that all formulas in $\Gamma$ are atomic or quantified equations, and all formulas in $\Delta$ are atomic.


For any node $\nu \in \Pi$ we denote with $\Gamma_\nu \vdash \Delta_\nu$ the sequent decorating $\nu \in \Pi$. We define $\mu \le \nu$ the ordering of $\Pi$ having the root as bottom element. We call induction on $\nu$ the induction on the distance of $\nu$ from the root. By induction on $\nu \in \Pi$ we may prove that  all formulas in $\Gamma_\nu$ are atomic or quantified equations, and all formulas in $\Delta_\nu$ are atomic and the rule with conclusion $\nu$ in $\Pi$ is $\id, \eqL, \eqR, \Weak, \forall L$ or is $p L, p R$ for some inductive predicate $p$ or is a substitution. 

We define the eigenvariables of a rule $p L$ are all the fresh variables of the rule. The eigenvariables of a substitution rule are all free variables in the assumption of the rule.
By a renaming we may assume that no eigenvariable of a rule with conclusion $\nu$ occurs in any $\mu \le \nu$. 

Assume we fixed a model $\setM$ of the language $L$. Assume we have an assignment $\sigma:\setM \rightarrow \FV(\Gamma_\nu \vdash \Delta_\nu)$, such that $\setM, \sigma \models\Gamma$ in the sense of Tarski. We call a production proof of an inductive predicate in $\setM, \sigma$ a proof $\Theta$. made of productions for the predicate and starting from atomic logical predicates true in $\setM, \sigma$.  We suppose that there is at least one value $\dummy \in \range(\sigma)$ and that for any inductive predicate $p_1(\vec{t}_1), p_k(\vec{t}_k) \in \Gamma$ we have some production proofs $\Theta_1, \ldots, \Theta_k$ for them in $\setM, \sigma$. 

We define a canonical extension $\sigma_\nu:\bigcup\{\FV(\Gamma_\nu \vdash \Delta_\mu) | \mu \le \nu\}\rightarrow \setM$ of $\sigma$, for each inductive predicate $p(\vec{t}) \in \Gamma_\nu$ a production proof for $p(\vec{t})$ in $\setM, \sigma_\nu$. 
We define a pruning $\Pi_1$ of $\Pi$ such that for all $\nu \in \Pi_1$ we have $\setM, \sigma_\nu \models \Gamma_\nu$.

If $\nu = $ the root of $\Pi$ then $\sigma_\nu = \sigma$ and production proofs for $p(\vec{t}) \in \Gamma_\nu$ are already given. Assume $\nu$ is the conclusion of a rule $r$ whose premise number $i$ is $\nu_i$. If $r = \id$ we are adone. Assume $r = p L$. Then $r$ has main formula some inductive predicate $p(\vec{t}) \in \Gamma_\nu$. We choose in $\Pi_1$ the premise $\mu = \nu_i$ corresponding to the production in the conclusion of $\Theta$, and we assign the eigenvariables of $r$ in such a way that the auxiliary formulas in $\nu_i$ are equal to the assumption of the production. For each auxiliary formula which is inductive we have an assumption $\Theta'$ of $\Theta$ proving it (hence the formula is true in $\setM, \sigma_\mu$). For each auxiliary formula $R(\vec{u})$ which is logical we have $\setM, \sigma \models R(\vec{u})$ by assumption of $\Theta$. All auxiliary formulas which are equations are true in $\setM, \sigma_{\nu_i}$ by definition of $\nu_i$. If $r$ is a substitution $\theta$ then we set $\sigma_\mu = \sigma_\nu \cup \sigma_\nu \theta$.  Assume $r \not = p L$ and $r$ is not a substitution.  Then we for all premises $\mu$ of $\nu$ we set $\mu \in \Pi_1$, and we assign any fresh variable in $\FV(\Gamma_\mu \vdash \Delta_\mu)$ to $\dummy \in \range(\sigma)$. 

We defined a pruning $\Pi_1 \subseteq \Pi$. By induction on $\nu \in \Pi_1$ we may prove that $\setM, \sigma_\nu \models \Gamma_\nu$. We already checked the case $r = \id, \eqR, p L$. If $r$ is a substitution $\theta$ then we set $\sigma_\mu = \sigma_\nu \cup \sigma_\nu \theta$. The thesis follows.  If $r = \forall L$ then $\Gamma_\mu = \Gamma_\nu, a[\vec{t}/\vec{x}]=b[\vec{t}/\vec{x}]$ for some $\forall \vec{x}.a=b \in \Gamma_\nu$, and $\sigma_\mu = \sigma_\nu$, hence $\setM, \sigma_\mu \models \Gamma_\mu$. If $r = \Weak, p R$ then $\Gamma_\mu \subseteq \Gamma_\nu$ and $\sigma_\mu = \sigma_\nu$, hence $\setM, \sigma_\mu \models \Gamma_\mu$. If $r = \eqL$ then each formula in $\Gamma_\nu$ is replaced by a formula with equal arguments and $\sigma_\mu = \sigma_\nu$, hence $\setM, \sigma_\mu \models \Gamma_\mu$. 

Every infinite path in $\Pi_1$ would have an infinite progressing trace, which cannot be because all inductive predicates in the left-hand side are true, therefore no predicate may progress forever. Thus, all paths of $\Pi_1$ are finite and $\Pi_1$ is finite. 

Assume $\nu \in \Pi_1$ and $A \in \Delta_\nu$. We say that $A$ has a production proof in $\Pi$ if there is a trace $\tau$ of $A$ in $\Pi_1$ ending in the main formula of $\id$ or $\eqR$, such that all $B$ in $\tau$ which are main formulas have all auxiliary formulas with a production proof in $\Pi$. 

By induction on the height of $\nu \in \Pi_1$ (existing because $\Pi_1$ is finite) we may prove that all $A \in \Delta_\nu$ with a production proof in $\Pi_1$ are true in $\setM, \nu$. 

By induction on the height of $\nu \in \Pi_1$ we may prove that some $A \in \Delta_\nu$ has a production proof in $\Pi_1$.

We apply these results to any cut-free $\infty$-proof $\Pi$ of $\Nat x \vdash \Tot x$ and we prove that $\Pi$ is not regular. If $\Pi$ were regular, we would have some maximum $k$ such that $S^k(y)$ or $S^k(0)$ occurs in $\Pi$. We choose $\setM = \Nat$, $\sigma(x) = k+1$ and $\Theta = $ the unique production proof of $\Nat n$. By induction on $\nu$ we prove that $\range(\sigma_\nu) \subseteq [0,k+1]$. We deduce that there is a production proof of $\Tot x$. There is no $\Tot t$ occurring in any $\Gamma_\nu$, therefore some ancestor $\Tot t \in \sigma_\nu$ of $\Tot x$ in $\Pi_1$ is proved from some $a(t,u)$. By induction on $\nu$ we may prove that $\sigma_nu(t) = \sigma(x) = k+1$. By definition of $k$ we have $u \equiv S^h(y)$ or $u \equiv S^h(0)$, hence $\sigma(u) \le \sigma_\nu(y) + h \le (k+1) + k < 2(k+1)$. Instead, from $\setM, \sigma_\nu \models a(t,u)$ we deduce $\sigma_\nu(t) = 2\sigma_\nu(u) = 2(k+1)$. This is a contradiction, therefore $\Pi$ is not regular. 

This rules out the possibility of proving $2n$ total. The same argument holds for any other map $f(\vec{n})$ which is definitively $\ge 2 \max(\vec{n})$, with the following proviso. We have to define the graph $a$ of $f$ inductively, then the totality $\Tot(\vec{x})$ of $f$ in $x$ by: $a(\vec{x},y)$ implies $\Tot(\vec{x})$. In the next subsection we provide an example for $f$: the Ackermann function.

%21:00 15/07/2018



\newpage

\subsection{The Ackermann function}
The Ackermann–Peter version of the Ackermann function is a map $\Ack:\Nat^2 \rightarrow 
\Nat$ defined by the following one-step rewriting rules, denoted by  $\reduces_1$. 
For all $x, y$ if $\Nat x, \Nat y$ we ask that:
\begin{enumerate}
\item
$\Ack(0,y) \reduces_1 Sy$ 
\item
$\Ack(Sx,0) \reduces_1 \Ack(x,1)$
\item
$\Ack(Sx,Sy) \reduces_1 \Ack(x,\Ack(Sx,y))$
\end{enumerate}
We write $\reduces_t$ for reduction $t$ times. 

We define the totality statement for Ackermann as a map $\Nat^2 \rightarrow 
\Nat$ on a pair $x, y$ as 
$\Tot(x,y) = \exists z, t. (\Nat z) \wedge (\Nat t) \wedge (\Ack(x,y) \reduces_t z)$, 
requiring that the image $z$ of $(x,y)$ is in $\Nat$.
The Ackermann binary function defines a family $\{Ack(x,\cdot)| x \in \Nat \}$ of unary
maps. The totality of the map $\Ack(x,\cdot)$ from $\Nat$ to $\Nat$ 
is formulated as the $\Pi^0_2$-predicate 
$T(x) = \forall y. \Nat y \rightarrow \Tot(x,y)$. The totality of $\Ack$ 
from $\Nat \times \Nat$ to $\Nat$ is stated with the
$\Pi^0_2$-formula $\forall x. \Nat x \rightarrow T(x)$.
 
\subsubsection{Totality of $\Ack:\Nat^2 \rightarrow \Nat$ (informal proof)}
By $\Pi^0_2$-induction on $x$ such that $\Nat x$ we may prove 

\begin{lemma}[Ackermann is Total (informal proof)]
\label{lemma-ackermann-informal}
$\Ack:\Nat^2 \rightarrow \Nat$ is a total map. 
\end{lemma}

\begin{proof}
What we in fact prove by induction on $x$ such that $\Nat x$ is the totality of 
$\Ack(x,\cdot)$ from $\Nat$ to $\Nat$, as stated by $T(x)$. This is the informal proof.

\begin{enumerate}
\item
{\bf Base Case: $x=0$}. $\Ack(0,\cdot)$ is total. Indeed, $\Ack(0,y) \reduces Sy$, therefore
$\Ack(0,\cdot)$ is the successor map, which is total from from $\Nat$ to $\Nat$.

\item
{\bf Inductive Case: from $x$ to $Sx$}. 
Assume that $\Nat x$ and that $\Ack(x,\cdot)$ is total in order to prove that
$\Ack(Sx,\cdot)$ is total. We prove by induction on $y$ such that $\Nat y$ that 
$\Ack(Sx,y) \reduces_t z$ for some $t,z$ such that $\Nat t$, $\Nat z$. 

\begin{enumerate}
\item
{\bf Base Case of the secondary induction: $y=0$}.
We prove that $\Ack(Sx,0) \reduces_t z$ for some $t,z$ such that $\Nat t$, $\Nat z$.
By principal induction hypothesis \emph{instanciated on $y=1$} 
we have $\Ack(x,1) \reduces_u z$ for some $u,z$ such that $\Nat u$, $\Nat z$, 
The thesis follows from $\Ack(Sx,0) \reduces_1 \Ack(x,1)$ and $\Ack(x,1) \reduces_u z$.

\item
{\bf Inductive Case of the secondary induction: from $y$ to $Sy$}.
We prove that $\Ack(Sx,Sy) \reduces_t z$ 
for some $t,z$ such that $\Nat t$, $\Nat z$. 
By principal induction hypothesis we have $\Ack(Sx,y) \reduces_u z$ for some $u,z$ such that 
$\Nat u, \Nat z$. By secondary induction hypothesis \emph{istanciated on $y=z$} 
we have $\Ack(x,z) \reduces_v w$. The thesis follows from 
$\Ack(Sx,Sy) \reduces_1 \Ack(x,z)$ and $\Ack(x,z) \reduces_v w$.

\end{enumerate}
\end{enumerate} 
\end{proof}
%10:34 09/11/2025

\subsection{Ackermann as an inductive predicate in $\LKID$}
We replace the statement $\exists t. (\Nat  t) \wedge \Ack(x,y) \reduces_t z$ in $\LKID$ 
as an \emph{inductive predicate} $\Ack(x,y,z)$. We replace the correctness 
statement $\exists z. \Nat z \wedge \Ack(x,y,z)$ with an \emph{inductive predicate}  
 $\Tot(x,y)$. We insert two simplifications:
 \begin{enumerate}
\item 
In the definition of $\Ack(x,y,z)$
we skip the index $t$ such that $\Nat$ counting the number of reduction steps.
\item
In the definition of $\Tot(x,y)$ we skip the requirement $\Nat z$.
The reason is that $\Nat z$ can be proved from 
$\Nat y, \Ack(x,y,z)$, by induction on the definition of $\Ack(x,y,z)$.
\end{enumerate}

\begin{definition}[The $R$-rules for $\Ack$ and $\Tot$]

\begin{enumerate}
\item
Assume $x,y,z,t$ are any variables and $\Gamma \vdash \Delta$ any sequent.
Then the $\Ack R$-rules for the inductive definition of $\Ack$ in $\LKID$ are: 

\begin{enumerate}
\item
$\infer{\Gamma \vdash \Delta, \Ack(0,y,Sy)}{}$
\\

\item
$\infer{\Gamma \vdash \Delta, \Ack(Sx,0,z)}
{\Gamma \vdash \Delta,\Ack(x,1,z)}$
\\

\item 
$\infer{\Gamma \vdash \Delta, \Ack(Sx,Sy,t)}
{\Gamma \vdash \Delta,\Ack(Sx,y,z) \ \ \ \Gamma \vdash \Delta, \Ack(x,z,t)}$
\end{enumerate}

\item
The only $\Tot R$-rule is: \ \ \ 
$\infer{\Gamma \vdash \Delta, \Ack(x,y,z)}{\Gamma \vdash \Delta,\Tot(x,y)}$ 
\end{enumerate}
\end{definition}

The $\Ack L$, $\Tot L$-rules are defined in $\LKID$ as a consequence, using \emph{fresh}
variables. 

$L$-rules are required to prove that the image of $(x,y)$ such that $\Nat x$,
$\Nat y$ is in $\Nat$ (see later).
$L$-rules are also required to prove that the image of $(x,y)$ such that $\Nat x$,
$\Nat y$ is, but we do not include this proof here.

Assume $x,y,z,t \not \in \FV(t,u,v,\Gamma,\Delta)$.

\begin{enumerate}
\item
The$L$-rule with conclusion $\Ack(t,u,v),\Gamma \vdash \Delta$ 
has $3$ premises, one for each $R$-rule for the same predicate:
\begin{enumerate}
\item
Premise 1 is: 
$t=0, u=y, v=Sy,\Gamma \vdash \Delta$
\item
Premise 2 is:  
$t=Sx, u=0, v=z, \Ack(x,1,z),\Gamma \vdash \Delta$
\item 
Premise 3 is:  
$t=Sx, u=Sy, v=z, \Ack(Sx,y,z), \Ack(x,z,t),\Gamma \vdash \Delta$
\end{enumerate}

\begin{center}
{\scriptsize          %in order to make the proof-tree smaller
  \hspace{-3cm} 
 $\infer
  {\Ack(t,u,v),\Gamma \vdash \Delta}
  {   \deduce{t\!=\!0, u\!=\!y, v\!=\!Sy,\Gamma \vdash \Delta}{(1)} & 
    	 \deduce{t\!=\!Sx, u\!=\!0, v\!=\!z, \Ack(x,1,z),\Gamma \vdash \Delta}{(2)} & 
    	 \deduce{t\!=\!Sx, u\!=\!Sy, v\!=\!z, \Ack(Sx,y,z), \Ack(x,z,t),\Gamma \vdash \Delta}{(3)}
  }
 $
}
\end{center}


\item
The$L$-rule with conclusion $\Tot(t,u), \Gamma \vdash \Delta$ has a unique premise:
$$t=x, u=y, \Ack(x,y,z),\Gamma \vdash \Delta$$  
\end{enumerate}
%17:33 02/07/2018

As an immediate consequence we have the following finite proofs, which are parts of 
the totality proof. 
We could prove that $\Tot(Sx,y), \forall z. \Nat z \rightarrow \Tot(x,z) \vdash \Tot(Sx,Sy)$,
but we only use the weaker property $\Ack(Sx,y,z), \Tot(x,z) \vdash \Tot(Sx,Sy)$.
%19:51 09/11/2025

\begin{lemma}
\label{lemma-ackermann-base}
There are finite \emph{cut-free and connective-free} proofs of:
\begin{enumerate}
\item
\label{lemma-ackermann-base-1}
$\Ack(x,1,z) \vdash \Ack(Sx,0,z)$

\item 
\label{lemma-ackermann-base-2}
$\Ack(Sx,y,z), \Ack(x,z,t) \vdash \Ack(Sx,Sy,t)$

\item
\label{lemma-ackermann-base-3}
$\vdash \Tot(0,y)$

\item 
\label{lemma-ackermann-base-4}
$\Tot(x,1) \vdash \Tot(Sx,0)$

\item 
\label{lemma-ackermann-base-5}
$\Ack(Sx,y,z), \Tot(x,z) \vdash \Tot(Sx,Sy)$
\end{enumerate}
The proofs use $L$-rule for $\Tot$ but \emph{no $L$-rule for $\Ack$}:
no \emph{induction on the definition of $\Ack$} is used.
\end{lemma}

\begin{proof}
\begin{enumerate}
\item
%\label{lemma-ackermann-base-1}
$\Ack(x,1,z) \vdash \Ack(Sx,0,z)$ is derived by $\Ack R$-rule with premise 
$\Ack(x,1,z) \vdash \Ack(x,1,z)$. This latter is derived by $\id$-rule.

\item 
%\label{lemma-ackermann-base-2}
$\Ack(Sx,y,z), \Ack(x,z,t) \vdash \Ack(Sx,Sy,t)$ is derived by $\Ack R$ with premise
$\Ack(Sx,y,z), \Ack(x,z,w) \vdash \Ack(Sx,y,z)$ and $\Ack(Sx,y,z), \Ack(x,z,w) \vdash  \Ack(x,z,w)$. These latter are derived by $\id$-rules.

\item
%\label{lemma-ackermann-base-3}
$\vdash \Tot(0,y)$ is derived by $\Tot R$-rule with premise $\vdash \Tot(0,y,Sy)$.
This latter is an $\Ack R$-rule.

\item 
%\label{lemma-ackermann-base-4}
$\Tot(x,1) \vdash \Tot(Sx,0)$ is derived by $\Tot L$-rule with premise 
$x=x', \Ack(x',1,z) \vdash \Tot(Sx,0)$ with $x',z$ fresh. This latter is derived
by $\Tot R$-rule with premise $x=x', \Ack(x',1,z) \vdash \Ack(Sx,0,z)$,
which is derived by $\eq$-rule with premise $\Ack(x,1,z) \vdash \Ack(Sx,0,z)$.
This is point \ref{lemma-ackermann-base-1} above.

\item 
%\label{lemma-ackermann-base-5}
$\Ack(Sx,y,z), \Tot(x,z) \vdash \Tot(Sx,Sy)$ is derived by $\Tot L$-rule with premise 
$x=x', z=z', \Ack(Sx,y,z), \Ack(x',z',t) \vdash \Tot(Sx,0)$ with $x',z',t$ fresh.
This latter is derived by $\Tot R$-rule with premise $x=x', z=z', \Ack(Sx,y,z), \Ack(x',z',t) \vdash  \Ack(Sx,Sy,z)$,
which is derived by $\eq$-rule with premise $\Ack(Sx,y,z), \Ack(x,z,t) \vdash \Ack(Sx,Sy,t)$.
This is point \ref{lemma-ackermann-base-2} above.
\end{enumerate}
\end{proof}

\subsubsection{Ackermann is a Partial Function on $\Nat$}
In $\CLKID^\omega$ we may bypass $\Pi^0_2$-induction and we provide a connective-free 
(hence quantifier-free) proof with cuts of the totality of $\Ack$. We prove first that all $\Ack(t)$ 
are a \emph{partial} map $:\Nat \rightarrow \Nat$: that is, if $\Nat t, \Nat u, \Ack(t,u,v)$ then
$\Nat v$. In fact we prove a little more, namely: if $\Nat u, \Ack(t,u,v)$ then $\Nat v$. 

%17:30 04/07/2018

\begin{lemma}[Ackermann is a partial function on $\Nat$]
\label{lemma-ackermann-partial}
Let $t,u$ be any terms of the language $0,S$. There is a connective-free cyclic proof 
{\bf with cuts} of $\Nat u, \Ack(t,u,v) \vdash \Nat v$ in the same language.
\end{lemma}


\begin{proof}
Let $x,y,z,w\not \in \FV(t,u,v)$ be pairwise distinct variables. We define a connective-free 
cyclic proof $\Pi_0$ of $\Nat u, \Ack(t,u,v) \vdash \Nat v$ using $\Ack L$ in the root, with 
$3$ open assumptions: 
\\

%\begin{enumerate}
%\item
%$ \Nat 1, \Ack(x,1,z) \vdash \Nat z$
%\item
%$ \Nat y, \Ack(Sx,y,z) \vdash \Nat z$
%\item
%$\Nat z, \Ack(x,z,w) \vdash \Nat w$. 
%\end{enumerate}

$\infer*[\Ack L] %having 3 vertical dots instear of an horizontal line
 {\Nat u, \Ack(t,u,v) \vdash \Nat v}
 {\deduce{\Nat 1, \Ack(x,1,z)   \vdash \Nat z} {(1)} & 
   \deduce{\Nat y, \Ack(Sx,y,z) \vdash \Nat z} {(2)} &
   \deduce{\Nat z, \Ack(x,z,w)  \vdash \Nat w}{(3)}
 }
$
\\

We define the $3$ premises of the rule $\Ack L$. Let $x,y,z,w$ be \emph{fresh}
variables.
\\

{\scriptsize
$
\infer[\Ack L] %having 3 vertical dots instear of an horizontal line
 {\Nat u, \Ack(t,u,v) \vdash \Nat v}
 {\infer*{\Nat u, t\!=\!0, u\!=\!y, v\!=\!Sy \vdash \Nat v} {} 
 & 
   \infer{\Nat u, t\!=\!x, u\!=\!1, v\!=\!z, \Ack(x,1,z) \vdash \Nat v} 
            {\deduce{\Nat 1, \Ack(x,1,z) \vdash \Nat z}{}}
 &
   \infer
   {\Nat u, t\!=\!Sx, u\!=\!y, v\!=\!w, \Ack(Sx,y,z), \Ack(x,z,w) \vdash \Nat v}
   {\infer[\tt cut]
     {\Nat y,  \Ack(Sx,y,z), \Ack(x,z,w) \vdash \Nat w}
     {\deduce{\Nat y, \Ack(Sx,y,z) \vdash \Nat z} {(2)} 
     &
       \deduce{\Nat z, \Ack(x,z,w)  \vdash \Nat w}{(3)}
     }
   }
 }
$
}

\begin{enumerate}
\item
\emph{Premise $1$} is $\mathbf{\Nat u, t=0, u=y, v=Sy \vdash \Nat v}$. 
It is proved by $\eqL$ with premise $\Nat y \vdash \Nat Sy$, 
which is proved by $\Nat R$ with premise 
$\Nat y \vdash \Nat y$. This latter is proved by the rule $\id$. 

\item
\emph{Premise $2$} is $\mathbf{\Nat u, t=x, u=1, v=z, \Ack(x,1,z) \vdash \Nat v}$. 
It is proved by $\eqL$ with premise $\Nat 1, \Ack(x,1,z) \vdash \Nat z$. This is the 
\emph{open assumption $1$}.

\item
\emph{Premise $3$} is 
$\mathbf{\Nat u, t=Sx, u=y, v=w, \Ack(Sx,y,z), \Ack(x,z,w) \vdash \Nat v}$. 
It is proved by $\eqL$ with premise $\Nat y,  \Ack(Sx,y,z), \Ack(x,z,w) \vdash \Nat w$. 
This latter is proved by a cut with two premises $\Nat y,  \Ack(Sx,y,z) \vdash \Nat z$ and 
$\Nat z, \Ack(x,z,w) \vdash \Nat w$. These  two premises are 
\emph{open assumptions $2$ and $3$}
\end{enumerate}

We define a regular infinitary cut-free proof tree $\Pi$ out of $\Pi_0$. 

We just defined $\Pi_0$ with $3$ open assumptions, 
for some variables $x, y, z, w\not \in \FV(t,u,v)$: 

\begin{enumerate}
\item
$\Nat 1, \Ack(x,1,z) \vdash \Nat z$
\item
$\Nat y, \Ack(Sx,y,z) \vdash \Nat z$
\item
$\Nat z, \Ack(x,z,w) \vdash \Nat w$
\end{enumerate} 

We can repeat the same operations we did on $\Nat u, \Ack(t,u,v) \vdash \Nat v$ 
on each of the $3$ open assumption. From each open assumption we obtain $3$ new open assumptions $\Nat 1, \Ack(x',1,z') \vdash \Nat z$ and $\Nat y', \Ack(Sx',y',z') \vdash \Nat z'$ and $\Nat z', \Ack(x',z',w') \vdash \Nat w'$, for some $x', y', z', w' \not \in \FV(x,y,z,t,u,v)$. We rename and we re-obtain the $3$ open assumptions: $\Nat 1, \Ack(x,1,z) \vdash \Nat z$ and $\Nat y, \Ack(Sx,y,z) \vdash \Nat z$ and $\Nat z, \Ack(x,z,w) \vdash \Nat w$, for $x, y, z, w\not \in \FV(t,u,v,x',y',z')$.  We defined a regular tree $\Pi$.

We prove the global trace condition for $\Pi$. In $\Pi$ whenever we prove $\Nat u, \Ack(t,u,v) \vdash \Nat v$ we derive it from $3$ sequents of the form $\Nat u', \Ack(t',u',v') \vdash \Nat v'$, and in each case there is a progressing trace from $\Ack(t,u,v) $ to $\Ack(t',u',v') $. Every infinite path of $\Pi$ has to pass infinitely many times from  $\Nat u, \Ack(t,u,v) \vdash \Nat v$ to $\Nat u', \Ack(t',u',v') \vdash \Nat v'$, therefore it includes an infinitely progressing trace from the unique atom $\Ack(t,u,v) $ in the root.
$\Box$ \end{proof}

From the fact that $\Ack$ is the graph of a partial map on $\Nat$ we derive that $\Ack$ is the graph of a total map on $\Nat$, using cuts but no logical connective.

\begin{lemma}[Ackermann is a total function]
\label{lemma-ackermann-total}
Let $t,u$ be any terms of the language $0,S$. 
There is a a connective-free cyclic proof with cuts of $\Nat t, \Nat u \vdash \Tot(t,u)$ in the same language.
\end{lemma}

%15:39 07/11/2025

\begin{proof}

We define a finite proof $\Pi_0$ of $\Nat t, \Nat u \vdash \Tot(t,u)$ from $3$ open assumptions: $(1) \Nat x, \Nat 1 \vdash \Tot(x,1)$ and $(2) \Nat t, \Nat y \vdash \Tot(t,y)$ and $(3) \Nat x, \Nat z \vdash \Tot(x,z)$ for any variables $x,y,z \not \in \FV(t,u)$.

We prove $\Nat t, \Nat u \vdash \Tot(t,u)$ by $\Nat L$ on $\Nat u$, from two premises: $\Nat t, u=0 \vdash \Tot(t,u)$ and $\Nat t, u=Sy, \Nat y \vdash \Tot(t,u)$. For each premise we provide a proof. 

\begin{enumerate}
\item
We prove $\Nat t, u=0 \vdash \Tot(t,u)$ by $\eqL$ with premise $\Nat t \vdash \Tot(t,0)$, then by $\Nat L$-rule with assumptions: $t=0 \vdash \Tot(t,0)$ and $t=Sx, \Nat x \vdash \Tot(t,0)$. 

\begin{enumerate}
\item
Proof of $t=0 \vdash \Tot(t,0)$. We use $\eqL$ with premise $\vdash \Tot(0,0)$, 
then $\Tot R$-rule and $\Ack(0,0,1)$, which is an $\Ack R$-axiom. 
\item
We prove $t=Sx, \Nat x \vdash \Tot(t,0)$ by $\eqL$-rule with premise 
$\Nat x \vdash \Tot(Sx,0)$, and this latter by a \emph{cut} 
between $\Nat x \vdash \Tot(x,1)$ and 
$\Tot(x,1) \vdash \Tot(Sx,0)$. $\Nat x \vdash \Tot(x,1)$ is proved by a cut between $\Nat 1$ and $\Nat x, \Nat 1 \vdash \Tot(x,1)$: this is \emph{(open assumption 1)}. 
$\Tot(x,1) \vdash \Tot(Sx,0)$ is proved by $\Tot L$-rule from $x' = x, y'=1, \Ack(x',y',z) \vdash \Tot(Sx,0)$, then by $\eqL$ with assmption $\Ack(x,1,z) \vdash \Tot(Sx,0)$, by $\Tot R$-rule from $\Ack(x,1,z) \vdash \Ack(Sx,0,z)$, by $\Ack R$-rule from $\Ack(x,1,z) \vdash \Ack(x,1,z)$, eventually by $\id$.
\end{enumerate}

\item
Proof of $\Nat t, u=Sy, \Nat y \vdash \Tot(t,u)$. We use $\eqL$ with premise $\Nat t, \Nat y \vdash \Tot(t,Sy)$. Then we use a cut with $\Nat t, \Nat y \vdash \Tot(t,y)$ and $\Nat t, \Nat y,  \Tot(t,y) \vdash \Tot(t,Sy)$.  $\Nat t, \Nat y \vdash \Tot(t,y)$ is \emph{(open assumption 2)}. 
We prove $\Nat t, \Nat y,  \Tot(t,y) \vdash \Tot(t,Sy)$ by $\Nat L$-rule with premises
$t=0, \Nat y,  \Tot(t,y) \vdash \Tot(t,Sy)$ and $t=Sx, \Nat x, \Nat y,  \Tot(t,y) \vdash \Tot(t,Sy)$. 

\begin{enumerate}
\item
We prove the first premise $t=0, \Nat y,  \Tot(t,y) \vdash \Tot(t,Sy)$ by $\eqL$ and $W L$ from $\vdash \Tot(0,Sy)$.  By $\Tot R$-rule it is enough to prove  $\vdash \Ack(0,Sy,SSy)$, which is some $\Ack R$-axiom.
\item
We prove the second premise $t=Sx, \Nat x, \Nat y,  \Tot(t,y) \vdash \Tot(t,Sy)$, by $\eqL$ and $W L$ from $\Nat x, \Tot(Sx,y) \vdash \Tot(Sx,Sy)$. By $\Tot L$-rule we have to prove $\Nat x, Sx=x', y=y', \Ack(x',y',z) \vdash \Tot(Sx,Sy)$, by $\eqL$ we have to prove $\Nat x, \Ack(Sx,y,z) \vdash \Tot(Sx,Sy)$. We prove $\Nat x, \Ack(Sx,y,z) \vdash \Tot(Sx,Sy)$
by a cut between $\Nat x, \Ack(Sx,y,z) \vdash \Nat z$ (provable by Lemma \ref{lemma-ackermann-partial}) and $\Nat x, \Nat z, \Ack(Sx,y,z) \vdash \Tot(Sx,Sy)$.
We prove  $\Nat x, \Nat z, \Ack(Sx,y,z) \vdash \Tot(Sx,Sy)$ by a cut between $\Nat x, \Nat z \vdash \Tot(x,z)$, which is is \emph{(open assumption 3)}, and $\Ack(Sx,y,z), \Tot(x,z) \vdash \Tot(Sx,Sy)$. We prove $\Ack(Sx,y,z), \Tot(x,z) \vdash \Tot(Sx,Sy)$ by $\Tot L$-rule and $\Ack(Sx,y,z), x=x',z=z', \Ack(x',z',w) \vdash \Tot(Sx,Sy)$. By $\eqL$ we have to prove $\Ack(Sx,y,z), \Ack(x,z,w) \vdash \Tot(Sx,Sy)$. By $\Tot R$ it is enough to prove that $\Ack(Sx,y,z), \Ack(x,z,w) \vdash \Ack(Sx,Sy,w)$, by $\Ack R$ that 
$\Ack(Sx,y,z), \Ack(x,z,w) \vdash \Ack(Sx,y,z)$ and that $\Ack(Sx,y,z), \Ack(x,z,w) \vdash  \Ack(x,z,w)$. These sequents are both $\id$-rules.
\end{enumerate}
\end{enumerate}

We define a regular infinite proof $\Pi$ from $\Pi_0$. In $\Pi_0$ we proved  $\Nat t, \Nat u \vdash \Tot(t,u)$ from $3$ open assumptions: $\Nat x, \Nat 1 \vdash \Tot(x,1)$ and $\Nat t, \Nat y \vdash \Tot(t,y)$ and $\Nat x, \Nat z \vdash \Tot(x,z)$ for any variables $x,y,z \not \in \FV(t,u)$. We may prove these latter from $\Nat x', \Nat 1 \vdash \Tot(x,1)$ and $\Nat t, \Nat y' \vdash \Tot(t,y')$ and $\Nat x', \Nat z' \vdash \Tot(x',z')$ for any variables $x',y',z' \not \in \FV(t,u,x,y,z)$. BY renaming we obtain $\Nat x, \Nat 1 \vdash \Tot(x,1)$ and $\Nat t, \Nat y \vdash \Tot(t,y)$ and $\Nat x, \Nat z \vdash \Tot(x,z)$ because we have $x,y,z \not \in \FV(t,u,x',y',z')$. Thus, we may define a regular proof $\Pi$ by unfolding $\Pi_0$.

We prove the global trace condition for $\Pi$. Each $\Nat t, \Nat u \vdash \Tot(t,u)$ is proved by: $(1) \Nat x, \Nat 1 \vdash \Tot(x,1)$ or $(2) \Nat t, \Nat y \vdash \Tot(t,y)$ or $(3)\Nat x, \Nat z \vdash \Tot(x,z)$. In cases $(1)$ and $(3)$ there is a progressing trace from $\Nat t$ to $\Nat x$. In case $(2)$ there is a trace from $\Nat t$ to $\Nat t$ and a progressing trace from $\Nat u$ to $\Nat y$. Take any infinite path $\pi$ of the proof $\Pi$, in order to prove that there is some infinitely progressing trace. $\pi$ passes infinitely many times through $(1)$ or $(2)$ or $(3)$. Assume $\pi$ passes infinitely many times through $(1)$ or $(3)$ (and any number of times through $(2)$). Then there is an infinite trace $\tau$ from $\Nat t$ in the root, $\tau$ is infinitely progressing (whenever $\tau$ passes through $(1)$ or $(3)$) and we are done. Assume $\pi$ passes finitely many times through $(1)$ or $(3)$. Then from some node $\xi$ on, $\pi$ passes cyclically through $(2)$, progressing from $\Nat u$ to $\Nat x$. From $\xi$ on, the trace $\upsilon$ from the occurrence of $\Nat u$ in $\xi$ progresses in \emph{every} step. Also in this case there is some infinitely progressing trace in $\pi$.
$\Box$ \end{proof}

%18:33 03/07/2018

\newpage

We may now prove that there is some connective-free cyclic proof whose cuts cannot be eliminated.


\begin{theorem} 
There a connective-free cyclic proof of $\Nat x, \Nat y \vdash \Tot(x,y)$ with cuts in the language $0,S$. There is no cut-free proof cyclic proof of $\Nat x, \Nat y \vdash \Tot(x,y)$ in the same language. 
\end{theorem}

\begin{proof} As for the map $f(x) = 2x$. $\Box$ \end{proof}

%\begin{proof}
%Wehave a connective-free cyclic proof with cuts by Lemma \ref{lemma-ackermann-total}. Assume that we have some cut-free proof $\Pi$ of $\Nat x, \Nat y \vdash \Tot(x,y)$in the language $0,S$. We claim that for some $k \in \Nat$ we have $\Ack(n,m) \le \max(n,m) + k$ for all $n,m \in\Nat$, and this cannot be.
%
%%The cut-free proof may be used to compute some bound to the first $p$ such that $\Ack(n,m,p)$. 
%
%%Cyclic proofs of $\Tot(x,y)$ include no witness in general, but from a cut-free proof $\Pi$ of $\Nat x, \Nat y \vdash \Tot(x,y)$ we may compute an upper bound, as follows. 
%
%
%We compute a bound $ \max(n,m) + k$ for $\Ack(n,m)$ as follows. We define a substitution $\sigma$ on $\Pi$ with $\sigma(x)=n$ and $\sigma(y)=m$. Assume $z$ is any eigenvariable of a rule $\Nat L$ for $\Nat t$. If $\sigma(t)=Sn$ we set $\sigma(z)=n$, otherwise we set $\sigma(z)=0$. In the second case, the second premise of $\Nat L$ includes $t = Sz$ in the left-hand side, and $\sigma(t) = 0 \not = 1 = \sigma(Sz)$. We set $\sigma(w)=0$ for any any other variable $w$ of $\Pi$. By definition $\sigma(w) \le \max(n,m)$ for all variables $w$ of $\Pi$.
%
%We cut the infinitary proof $\Pi$ whenever we find some $t=u$ in the left-hand side such that $\sigma(t) \not = \sigma(u)$.  We call this new rule $\eq$-elimination. We call $\Pi_0$ the resulting proof, with $\eq$-elimination.
%
%We claim that $\Pi_0$ is finite. By the sub-formula property, in $\Pi_0$, we may have equalities, occurrences of some $\Nat t$ only in the left-hand side, occurrences $\Ack (a,b,c)$ and $\Tot(d,e)$ only in the right-hand side. The only possibility for an $\id$-rule is $\id$ for an equality $t=u$ such that $\sigma(t)=\sigma(u)$.  $\Pi$ only includes structural rules different from cut: $W$-rules, exchange, contraction, $\id$. $\Pi$ may include induction rules $\Nat L$ in the left and $\Ack R$, $\Tot R$ in the right; rules $\eqL$ and $\eqR$ and $\eq$-elimination. 
%Any infinite branch $\pi$ includes some infinitely progressing trace $\tau$ from some $\Nat t$. Assume $\sigma(t)=n$: then $\pi$ is stopped after $n+1$ progresses of $\tau$, because we reach some sequent including $u=Sx$ where $\sigma(u)=0$. Thus, $\Pi_0$ is binary and with all branches finite, hence $\Pi_0$ is finite by  K\"{o}nig's Lemma. 
%
%%All rules of $\Pi_0$ are intuitionistically correct. 
%Let $\Pi_1$ be any subtree of $\Pi_0$. By induction on $\Pi_1$, we may prove that there are $3$ possibilities: either the root of $\Pi_1$ has some false equality in the left, or some true equality or $\Ack(t,u,v)$ in the right, or some $\Tot(t,u)$ in the right and $\Pi_1$ includes some term which is a true witness for $\Tot(t,u)$. The witness is a term $S^k z$ or $S^k 0$ of $\Pi_0$, with $\sigma(z) \le \max(n,m)$. There are finitely many$k$ with $S^k x$ or $S^k 0$ occurring in $\Pi$ by regularity, hence the same holds for in $\Pi_0$. If we take the maximum $k$, then any witness in $\Pi_0$ is $\le \max(n,m) + k$. If we take the sequent $\Nat x, \Nat y \vdash \Tot(x,y)$ at the root of $\Pi_0$, then $\Nat x$, $\Nat y$ are true w.r.t. $\sigma$, hence we have a true witness for $\Tot(x,y)$ in $\Pi_0$ with $x,y$ assigned to $n,m$. We deduce the bound $\Ack(n,m) \le \max(n,m) + k$ for $k$ depending on $\Pi$ and for all $n,m \in \Nat$. This cannot be. We conclude that a cut-free proof $\Pi$ of $\Nat x, \Nat y \vdash \Tot(x,y)$ cannot exists in the language $0,S$. 
%$\Box$ \end{proof}

%As a corollary, there is no cut-free proof $\Pi$ of $\Nat t, \Nat u \vdash \Tot(t,u)$ in any language, because we may replace any subterm in $\Pi$ different from $0$, $St$ with $0$ while preserving the correctness of $\Pi$.

%A similar proof shows that there is a connective-free proof with cuts in the language $0,S$ that the twice map is total, and there is not cut-free proof of the same result in $\CLKID^\omega$ (with the full language).
%17:14 05/07/2018



\section{Upper and lower bound for the fragment of cyclic proofs with cut-elimination}
In \S \ref{section-failure} we proved that normalization fails for atomic formulas if we allow variable discarding productions in inductive definition. In this section we consider a single inductive definition, 
$\Nat$ (the simplest relevant example), discarding no variables. We prove that if we consider only cuts over formulas $\Nat t$  then normalization succeeds for totality of polynomials and fails for totality of exponential.

Fix any production for an inductive definition in $\LKID$. We recall that a production discards a variable if the variable occurs in the assumption and not in the conclusion. The definition of $\Nat$ discards no variables. The definition of the graph $a(x,y)$ of the map $f(x)=2x$ discards no variables. The inductive definition $\Tot(x)$ of totality of $f$ is: ``if $a(x,y)$ then $\Tot(x)$'' discards the variable $y$. Assume $A$ is quantifier-free. A cut occurrence between twos formulas $A$ is local if some ancestor of one $A$ is the main formula of a $p L$ rule with eigenvariable $x$, and some ancestors of the other $A$ has the same $x$ free.
%all inductive predicates in $A$ discard no variables and no $x \in \FV(A)$ is an eigenvariable of the proof. 

In this section we provide some cuts between two occurrences of $\Nat(t)$ with some $x \in \FV(t)$ which is an eigenvariable for the other occurrence. We show that some of these cuts cannot be eliminated.

The cyclic proofs we consider are again totality proofs for maps, but instead of using a local inductive definition $T(x)$ of totality as in \S \ref{section-failure} we express totality of $f$ in the form: $E_f, \Nat x \vdash \Nat f(x)$, with $E_f$ the equational axioms of $f$.  We prove that all polynomials with coefficients in $\Nat$ have cut-free proofs of totality, while the map $2^x$ has some totality proof whose cuts are all of the form $\Nat(2^y)$, but no totality proof which is cut-free. Our first steps are proving cut-free that sum and product are total.

Let $E_+ = \forall x.x+0=x, \forall x,y.x+Sy=S(x+y)$. We claim that there is a cut-free proof of the totality of $+$.

\begin{lemma}
\label{lemma-sum-total}
Assume $E_+ \subseteq \Gamma$ and $t$ is any term.
\begin{enumerate}
\item
There is a cut-free cyclic proof of $\Gamma, \Nat u \vdash \Nat (t+u)$ with free assumption $\Gamma \vdash \Nat t$.
\item
There is a cut-free cyclic proof of $E_+, \Nat t, \Nat u \vdash \Nat (t+u)$. 
\item
There is a cut-free cyclic proof of $E_+, \Nat t \vdash \Nat (t+t)$. 
\end{enumerate}
\end{lemma}

\begin{proof}
\begin{enumerate}
\item
We prove $\Gamma, \Nat u \vdash \Nat (t+u)$ by $\Nat L$-rule for $\Nat u$. We have to prove: \emph{(1)} $\Gamma, u=0 \vdash \Nat (t+u)$ and \emph{(2)} $\Gamma, \Nat z, u=Sz \vdash \Nat (t+u)$ for $z \not \in \FV(\Gamma, t,u)$. 
\begin{enumerate}
\item
Proof of $\Gamma, u=0 \vdash \Nat (t+u)$. We use $\eqL$ and $\Gamma \vdash \Nat (t+0)$, then $\forall L$ for $\forall x.(x+0)=x \in \Gamma$ and $\Gamma,  (t+0)=t \vdash \Nat (t+0)$, then $\eqL$ and  $\Gamma \vdash \Nat t$. This latter is a free assumption.
\item
Proof of $\Gamma, \Nat z, u=Sz \vdash \Nat (t+u)$. We use $\eqL$ and $\Gamma, \Nat z \vdash \Nat (t+Sz)$, then $\forall L$ for $\forall x.(x+Sy)=S(x+y) \in \Gamma$ and $\Gamma,  (t+Sz)=S(t+z), \Nat z \vdash \Nat (t+Sz)$, then $\eqL$ and  $\Gamma, \Nat z \vdash \Nat S(t+z)$, eventually $\Nat R$ and $\Gamma, \Nat z \vdash \Nat (t+z)$. Here we repeat cyclically: next time we choose $z' \not \in \FV(\Gamma, t,u,z)$ as eigenvariable, next time we choose $z$ again. In this way we have a regular proof.
 \end{enumerate}
In order to have a cyclic proof we have to check the global trace condition. The trace of $\Nat u$ progresses whenever we pass a $\Nat L$-rule, otherwise it is stationary. Every infinite path has infinitely many $\Nat L$-rule, therefore the trace of $\Nat u$ is infinitely progressing. 
\item
By point $1$ above there is a cut-free cyclic proof of $E_+, \Nat t, \Nat u \vdash \Nat (t+u)$ with free assumption $E_+, \Nat t, \Nat u \vdash \Nat t$. We prove the free assumption with $\id$. Regularity is preserved because we add no sub-proofs, and every infinite path in the new proof is in the old proof, therefore is includes some infinitely progressing trace. Thus, we defined a cut-free cyclic proof (without assumptions).
\item
By point $2$ above if we choose $t \equiv u$.
\end{enumerate}
$\Box$ \end{proof}

Let $E_* = E_+, \forall x.x*0=0, \forall x,y.x*Sy=x*y+x$. We claim that there is a cut-free proof of the totality of $*$.

\begin{lemma}
There is a cut-free proof of $E_*, \Nat x, \Nat y \vdash \Nat (x*y)$. 
\end{lemma}

\begin{proof}
We prove $E_*, \Nat x, \Nat y \vdash \Nat (x*y)$ by $\Nat L$-rule for $\Nat *$. We have to prove \emph{(1)} $E_*, \Nat x, y=0 \vdash \Nat (x*y)$ and \emph{(2)} $E_*, \Nat x, \Nat z, y=Sz \vdash \Nat (x*y)$. 
\begin{enumerate}
\item
Proof of $E_*, \Nat x, y=0 \vdash \Nat (x+y)$. We use $\eqL$ and $E_*, \Nat x \vdash \Nat (x*0)$, then $\forall L$ and $E_*,  (x*0)=0, \Nat x \vdash \Nat (x*0)$, then $\eqL$ and  $E_*, \Nat x \vdash \Nat 0$, eventually $\Nat R$.
\item
Proof of $E_*, \Nat x, \Nat z, y=Sz \vdash \Nat (x*y)$. We use $\eqL$ and $E_*, \Nat x, \Nat z \vdash \Nat (x*Sz)$, then $\forall L$ and $E_*,  (x*Sz)=x*z+x, \Nat x, \Nat z \vdash \Nat (x*Sz)$, then $\eqL$ and  $E_*, \Nat x, \Nat z \vdash \Nat x*z+x$. By Lemma \ref{lemma-sum-total} there is a cut-free cyclic proof $\Pi$ of $E_*, \Nat x, \Nat z \vdash \Nat x*z+x$ with free assumption $E_*, \Nat x, \Nat z \vdash \Nat x*z$. Here we repeat cyclically, next time we choose $y$ as eigenvariable, in order to obtain a regular proof.
 \end{enumerate}
We check the global  trace condition. Every infinite path $\pi$ either passes infinitely many times through a $\Nat L$-rule inferring $\Nat y$, with the trace of $\Nat y$ progressing each time, or from some node on is inside some inserted proof $\Pi$ of $E_*, \Nat x, \Nat z \vdash \Nat x*z+x$. In the first case the global trace condition for $\pi$ is satisfied. In the second case, the suffix of $\pi$ which is in $\Pi$ includes from some node on some infinite progressing trace, the same holds for $\pi$ and the global trace condition for $\pi$ is satisfied.
$\Box$ \end{proof}

If global cuts are removable, then from $t$ total and $u$ total we deduce that $t[u/x]$ is total. If $x \not \in \FV(t)$ this is immediate. Assume $x \in \FV(t)$. Let $E_t$ be the union of the lists of axioms for function symbols in $t$, and $E_u$ for function symbols in $u$. From $x \in \FV(t)$ we deduce that the function symbols in $[t[u/x]$ are those in $t$ or in $u$, hence $E_t, E_u = E_{t[u/x]}$. Let $\vec{x} = \FV(t,u)$. From $E_t, \Nat \vec{x}, \Nat x \vdash \Nat t$ we deduce $\Nat \vec{x}, \Nat u \vdash \Nat t[u/x]$ by substitution. From $E_u, \Nat \vec{x} \vdash \Nat u$ and cut we conclude $E_t, E_u, \Nat \vec{x} \vdash \Nat t[u/x]$, that is, $E_{t[u/x]}, \Nat \vec{x} \vdash \Nat t[u/x]$. The cut is global because it is the unique cut in the proof, hence by renaming eigenvariables we may may prevent an eigenvariable in one side of the proof to occur in the other side of the proof. If we may remove global cuts, then we may obtain a cut-free proof of $E_{t[u/x]}, \Nat \vec{x} \vdash \Nat t[u/x]$. 

By repeating this reasoning we obtain that all polynomials, being compositions of $0$, $S$, $x+y$ and $x*y$ are provably total with a cut-free proof, provided global cuts are removable. 

With the same reasoning, if global cuts are removable, and there is a cut-free proof of totality for $2^x$, then there is a cut-free proof of totality for all poly-exponential maps, obtained composing $0$, $S$, $x+y$, $x*y$ and $2^x$. We will prove that this cannot be: at least one exponential map has no cut-free proof of totality.

%09:53 16/07/2018


\begin{lemma}[Failure of cut-elimination for some exponential map]
\label{lemma-exp-failure}
There is some exponential map $f(x) = 2^{(x+3)^2}+ x$ with no cut-free totality proof. If global cuts may be eliminated, then the same holds for $2^x$.
\end{lemma}

\begin{proof}
Assume there is some cut-free proof $\Pi$ of totality of $f(x)$. Let us consider the set of traces from some inductive predicate in the left-hand side of a node $a$ to some inductive predicate in the left-hand side of a node $b$ in the graph representation of the cyclic proof. Assume there is a path from $a$ to $b$. Then set of traces $i$ from $a$ to $b$ and the set of traces $j$ from $b$ to $c$ may be composed by composing each trace with every compatible trace. The result $ji$ is the set of traces from $a $ to $c$. We add an empty trace $0$ and we set $ji = 0$ if the last node of $i$ is different from the first node of $j$. Assume $\pi$ is any infinite path in the infinite unfolding of the cyclic proof. Then the nodes of $\pi$ define a complete graph whose edges are colored on $M$, with coloring compatible with composition, and with all colors $\not = 0$.  

Any homogeneous subset of $\pi$ with at least $3$ elements $a,b,c$ has all edges colored with some idempotent $i \in M$, $i \not 0$. Indeed, if $i$ colors $(a,b)$ and $(b,c)$ and $(a,c)$, then by compatibility with composition $(a,c)$ has also color $ii$. Thus, $ii = i$. 

By Ramsey Theorem there is some infinite homogeneous subset, a set with all edges in the same color $i$, with $i \not = 0$, $i$ idempotent as we just proved. If it were $a \not = b$ then $ii = 0$, hence $i=0$, contradiction. Thus, $i$ connects $a$ with itself. There is an infinite subsequence $a,a,a,\ldots$ of nodes of $\pi$, with each connection from one $a$ to another one with color $i$.

By the global trace condition there is some infinitely progressing trace $\tau$ in $\pi$. $\tau$ passes  infinitely many times through $a$, and $a$ has finitely many inductive predicate in the left-hand side. Thus, $\tau$ passes infinitely many times through the same inductive formula $p(\vec{t})$. Take the first occurrence of $p(\vec{t})$ in $\tau$, the first progress point of $\tau$ after it, then the first occurrence of $p(\vec{t})$ in $\tau$ after the progress point. We deduce that there are two nodes $a$ such that $p(\vec{t})$ in the first $a$ progresses in $\tau$ to $p(\vec{t})$ in the second $a$. The color between the two $a$'s is $i$, therefore $p(\vec{t})$ progress to $p(\vec{t})$ in the color $i$. Thus, in any homogeneous set with $h+1$ elements the inductive predicate $p(\vec{t})$ progresses at least $h$ times. 

In $\Pi$, the formula $p(\vec{t})$ is some ancestor of the formula $\Nat x$ in the conclusion of $\Pi$, hence some ancestor of $\Nat x$ progresses at least $h$ times in any homogeneous set with $h+1$ elements. 

Let $R(m, \ldots, m)$ be the diagonal Ramsey number for $m$: is the minimum size of a complete graph with  $k$ colors having some homogeneous sets of cardinality $m$. A folk-lore result says that $R(m, \ldots, m) \le k^{km}$. We claim that every poly-exponential map $f(n)$ with a cut-free proof of totality has upper bound $r(n) = R(n+3,\ldots,n+3)$ for all $n \in \Nat$. 

In order to prove it, fix any assignment $\sigma(x) = n$ for $x$ and we extend it to all sequents in $\Pi$. We call the canonical extension of $\sigma$ the extension assigning $\max(0,p-1)$ to the eigenvariable of the case rule with main formula $\Nat u$ and $\sigma(u) = p$. There is exactly one path $\pi$ of $\Pi$ with all formulas in the left-hand side true with respect to the canonical extension of $\sigma$. $\pi$ has length $< r(n)$, because a path with length $r(n)$ would include an homogeneous set of $n+3$ elements, hence $n+2$ progresses of $\Nat x$, and this is impossible because $\sigma(x) = n$. $\pi$ includes $\ge f(n) - n$ rules $\Nat R$, then either the rule $\id$ or the rule $\Nat R$ with main formula $\Nat 0$. $r(n)$ cannot be a bound for $f(n) - n$ for all poly-exponential maps $f$. For instance, we have $f(z) = 2^{(z+3)^2}+ z > k^{k(z+3)} + z \ge r(z) + z \ge f(z)-z+z= f(z)$ for all $z > k\log_2(k)-3$, because in this case we have  $2^{(z+3)^2} > 2^{k\log_2(k)(z+3)} = (2^{\log_2(k)})^{k(z+3)} = k^{k(z+3)} = r(z)$.

$\Box$ \end{proof}

%Therefore if we prove that every poly-exponential map with a cut-free proof of totality has upper bound the diagonal Ramsey number $R(x,\ldots,x)$, then, assuming we may eliminate global cuts, we derive that $2^x$ has no cut-free proof of totality. 

We prove that if global cuts are eliminable, then $E_{2^x}, \Nat x \vdash \Nat 2^x$ has no cut-free cyclic proof. Instead, $E_{2^x}, \Nat x \vdash \Nat 2^x$ has a cyclic proof of totality whose cuts are all local cuts.
All cuts of this proof are local.

%13:56 15/07/2018


\begin{lemma}[Totality proof with cuts for $2^x$]
\label{lemma-exp-totality}
There is proof of $E_{2^x}, \Nat x \vdash \Nat 2^x$ with all cuts local.
\end{lemma}

\begin{proof}
We already checked that there is no cut-free proof of $E_{2^x}, \Nat x \vdash \Nat 2^x$: we provide a proof with cuts on $\Nat 2^y$. We use $\Nat L$ with main formula $\Nat x$ and premises $E_{2^x}, x=0 \vdash \Nat 2^x$ and $E_{2^x}, \Nat y, x = Sy\vdash \Nat 2^x$ for $y \not \equiv x$. 
\begin{enumerate}
\item
Proof of $E_{2^x}, x=0 \vdash \Nat 2^x$. We use $\eqL$ and $E_{2^x} \vdash \Nat 2^0$. Then we use $\eqL$ on $2^0=1 \in E_{2^x}$ with premise $E_{2^x} \vdash \Nat 1$, where $1 \equiv S0$. We prove $E_{2^x} \vdash \Nat S0$ from $\Nat R$ for $S$ and $E_{2^x} \vdash \Nat 0$, then by $\Nat R$ for $0$. 
\item
Proof of $E_{2^x}, \Nat y, x = Sy\vdash \Nat 2^x$. We use $\eqL$ and $E_{2^x}, \Nat y \vdash \Nat 2^{Sy}$. Then we use $\forall L$ on $\forall x.2^{Sx} = 2^x + 2^x$ and $E_{2^x}, 2^{Sy} = 2^y + 2^y, \Nat y\vdash \Nat 2^{Sy}$, then $\eqL$ and $E_{2^x}, \Nat y \vdash \Nat 2^y + 2^y$. We derive this latter by a cut between $E_{2^x}, \Nat y \vdash \Nat 2^y$ and $E_{2^x}, \Nat 2^y \vdash \Nat 2^y + 2^y$. In the first premise we repeat cyclically the proof, in the second one we apply Lemma \ref{lemma-sum-total}. The second time we choose $x$ as eigenvariable, in order to have a regular proof. 
\end{enumerate}
We check the global trace condition. Any infinite path in the proof either passes infinitely many times through $E_{2^x}, \Nat x \vdash \Nat 2^x$ and $E_{2^x}, \Nat y \vdash \Nat 2^y$, with the trace of $\Nat x$ progressing each time, or from some node on is in some proof of $E_{2^x}, \Nat 2^y \vdash \Nat 2^y + 2^y$. In this case there is a trace infinitely progressing from some point on. 

We check that the cut on $\Nat 2^y$ is local. The first premise of the cut on $\Nat 2^y$ has ancestor $\Nat 2^x$, therefore has ancestor $\Nat 2^y + 2^y$. The ancestors of  $\Nat 2^y + 2^y$ include the eigenvariable $z$ of the case rule of $E_{2^x}, \Nat 2^y \vdash \Nat 2^y + 2^y$. The second premise of the cut on $\Nat 2^y$ is the main formula inferred by the case rule of $E_{2^x}, \Nat 2^y \vdash \Nat 2^y + 2^y$, therefore has ancestor $\Nat z$ with $z$ eigenvariable of the case rule. Thus, the cut is local.

$\Box$ \end{proof}

%20:20 14/07/2018
%20:00 26/09/2018

%%%%%%%%%%%%%%%%%%%%%%%%%%%%%%%
% 10:06 30/09/2018 
%%%%%%%%%%%%%%%%%%%%%%%%%%%%%%%

\section{A last example: the transitive closure}
We prove that the equivalence of right-to-left and left-to-right definitions of the transitive closure has a cyclic proof but no cut-free cyclic proof. 

A word before we start: this example is rather surprising. The equivalence just unfolds the definition of right-to-left transitive closure, generating a list, then reverses the list in square time, and eventually proves left-to-right transitive closure. Our interpretation of the lack of cut-elimination is: this algorithm requires non-constant space, having to remember the list unfolding the right-to-left transitive closure. For this reason, the proof requires left-hand sides of unbounded size and it is not regular. Apparently, cut-elimination is only possible when the proof uses a polynomial-time algorithm requiring constant space.
\\

Equivalence of two transitive definitions is a simplified version of an example of a proof from Separation Logic. 
We consider a first order language $L = \{R\}$, with no constants and no function symbols, only one binary predicate $R(x,y)$ and variables. 
We extend $L$ with two equivalent definitions of the transitive closure of $R$: $\Backward(x,y)$ and $\Forward(x,y)$. 
\begin{enumerate}
\item
$\Backward(x,y)$ means:  ``we recursively define an $R$-chain from $x$ to $y$ by listing $R$-related elements right-to-left, from $y$ to $x$''. In this definition we first list $R(z',y)$, then $R(z'',z')$, \ldots, until we reach some $R(x,z)$. 
\item
$\Forward(x,y)$ means: ``we recursively define an $R$-chain from $x$ to $y$ by listing $R$-related elements left-to-right, from $x$ to $y$''.  In this definition we first list $R(x,z')$, then $R(z',z'')$, \ldots, until we reach some $R(z,y)$''. 
\end{enumerate}
$\Backward(x,y)$ and $\Forward(x,y)$ are equivalent because we may recursively reverse the order with which we list the $R$-chain from $x$ to $y$. We will inductively define $\Backward$ and $\Forward$, then prove that $\Backward(x,y) \vdash \Forward(x,y)$ has a cyclic-proof but no cut-free cyclic proof.

The productions of $\Backward$ are: if $R(x,y)$ then $\Backward(x,y)$; if $\Backward(x,y)$, $R(y,z)$ then $\Backward(x,z)$.  
The productions of $\Forward$ are: if $R(x,y)$ then $\Forward(x,y)$; if $R(x,y)$, $\Forward(y,z)$ then $\Forward(x,z)$.  These productions define the rules: $\Backward L$, $\Backward R$, $\Forward L$ and $\Forward R$.


As a Lemma, we first prove (cut-free) the following statement: ``if we take a left-to-right $R$-chain and we add one step to the right we obtain a left-to-right $R$-chain''. 

\begin{lemma}
\label{lemma-premise-cut}
$\Forward(x,z), R(z,y) \vdash \Forward(x,y)$ has some cut-free cyclic proof $\Pi$. 
\end{lemma}

\begin{proof}
\begin{enumerate}
\item
We first define a regular infinite proof $\Pi$. We choose any variable $t \not \equiv x,z,y$ and we infer the conclusion of $\Pi$ by $\Forward L$-rule.
\begin{enumerate}
\item
The base case is $x=x', z=z', R(x',z'), R(z,y) \vdash \Forward(x,y)$ for $x',z'$ fresh. We use $\eqL$ with premise $R(x,z), R(z,y) \vdash \Forward(x,y)$. We prove it by $\Forward R$-rule with premises:  $R(x,z), R(z,y) \vdash \Forward(x,z)$ and $R(x,z), R(z,y) \vdash R(z,y)$. The first premise is proved by $\Forward R$ with premise $R(x,z), R(z,y) \vdash R(x,z)$, then $\id$. The second premise is proved by $\id$.
\item
The induction case is $x=x', z=z', R(x',t), \Forward(t,z'), R(z,y) \vdash \Forward(x,y)$  for $x',z',t$ fresh. We use $\eqL$ twice and we have to prove $R(x,t), \Forward(t,z), R(z,y) \vdash \Forward(x,y)$. We use $\Forward R$ with two premises: $R(x,t), \Forward(t,z), R(z,y) \vdash R(x,t)$ and $R(x,t), \Forward(t,z), R(z,y) \vdash \Forward(t,y)$. The first premise is an $\id$-rule. We prove the second premise by $\Weak$ and $\Forward(t,z), R(z,y) \vdash \Forward(t,y)$. This latter is a renaming by a fresh variable $t$ of the goal, hence we are done.
\end{enumerate}
\item
We check the global trace condition for $\Pi$. The unique infinite path is the rightmost path $\pi$. $\pi$ moves from $\Forward(x,z), R(z,y) \vdash \Forward(x,y)$ to $\Forward(t,z), R(z,y) \vdash \Forward(t,y)$, and each time the unique assumption $\Forward(x,z)$ progresses. Thus, the trace connecting all ancestors of $\Forward(x,y)$  is infinitely progressing in $\pi$.
\end{enumerate}
$\Box$ \end{proof}

We prove $\Backward(x,y) \vdash \Forward(x,y)$ by a cut between the Lemma $\Forward(x,z), R(z,y) \vdash \Forward(x,y)$ and the very sequent $\Backward(x,y) \vdash \Forward(x,y)$ we are proving. This cut is of particular kind. It is a cut on a sequent decorating a descendant of the conclusion of the cut itself, and this is considered a restricted version of cut. We will also prove that this cut is a local cut in our terminology.  

\begin{lemma}
\label{lemma-cut-proof}
$\Backward(x,y) \vdash \Forward(x,y)$ has some cyclic proof $\Pi$ having a local cut. 
\end{lemma}
%11:54 20/09/2018
\begin{proof}
\begin{enumerate}

\item
We define a regular infinite proof $\Pi$. We infer the conclusion of $\Pi$ by $\Backward L$-rule.
%19:45 18/09/2018
\begin{enumerate}
\item
The base case of the $\Backward L$-rule is $x=x', y=y', R(x',y') \vdash \Forward(x,y)$ for $x',y'$ fresh. We apply $\eqL $ twice and we have to prove $R(x,y) \vdash \Forward(x,y)$. This is provable by $\Forward R$-rule with premise $R(x,y) \vdash R(x,y)$, which is an $\id$-rule.
\item
The inductive case of the $\Backward L$-rule is $x=x', y=y', \Backward(x',z), R(z,y') \vdash \Forward(x,y)$. We apply $\eqL $ twice and we have to prove $\Backward(x,z), R(z,y) \vdash \Forward(x,y)$. We prove it by a cut on $\Forward(x,z)$: the premise of the cut are $\Backward(x,z) \vdash \Forward(x,z)$ and $\Forward(x,z), R(z,y) \vdash \Forward(x,y)$. The first premise $\Backward(x,z) \vdash \Forward(x,z)$ of the cut is the original goal renamed with $z$ fresh. The second premise of the cut is Lemma \ref{lemma-premise-cut}.
\end{enumerate}

\item
We check that the cut on $\Forward(x,z)$ is local in $\Pi$. \emph{(i)}. $\Forward(x,z)$ in the first premise has ancestor $\Forward(x,y)$ in the goal, which has ancestor  $\Forward(x,y)$ in $\Forward(x,z), R(z,y) \vdash \Forward(x,y)$. This latter has ancestor $\Forward(t,y)$, which includes the eigenvariable $t$ of a $\Forward L$-rule. \emph{(ii)}. $\Forward(x,z)$ in the second premise is inferred from a $\Forward L$-rule and $\Forward(t,z')$, with $t$ eigenvariable of the rule. From \emph{(i)} and \emph{(ii)} we deduce that the same eigenvariable occurs in some ancestors of the left and right cut formula, hence the cut is local.

\item
We check the global trace condition for $\Pi$. Consider any infinite path $\pi$ in $\Pi$. If $\pi$ eventually moves to $\Forward(x,z), R(z,y) \vdash \Forward(x,y)$, proved by Lemma \ref{lemma-premise-cut}, then $\pi$ has an infinitely progressing trace from the unique assumption $\Forward(x,z)$. Otherwise $\pi$ always moves from $\Backward(x,y) \vdash \Forward(x,y)$ to $\Backward(x,z) \vdash \Forward(x,z)$, and each time the unique inductive assumption $\Backward(x,y)$ progresses. In both cases, the global trace condition is satisfied.
\end{enumerate}
$\Box$ \end{proof}

%%%%%%%%%%%%%%%%%%
%HARD TO EXPLAIN
%%%%%%%%%%%%%%%%%%
%Remark that in this case the local cut produces during cut-elimination a substitution of the form $[t/t]$. This substitution produces no new formulas. Apparently, repeating $[t/t]$ a potentially infinite number of times during cut-elimination does not destroy regularity. However, there is no cut-free proof of $\Backward(x,y) \vdash \Forward(x,y)$. The reason is that the existence substitution requires to rename $t$ in order to avoid that a rule with eigenvariable $t$ captures a variable $t$ bound by a different rule but  with the same eigenvariable $t$. This potentially infinite renaming destroys regularity, as we will see. 
%%%%%%%%%%%%%%%%%%

%15/09/2018 10:33
An $\infty$-proof is an infinitary proof with the global trace condition, hence it is a sound proof. It may be not regular. 

\begin{lemma} 
\label{lemma-cut-free-non-regular-proof}
There is a cut-free non-regular $\infty$-proof of $\Backward(x,y) \vdash \Forward(x,y)$
\end{lemma}

\begin{proof}
For all $m \in \Nat$, there  is a finite cut-free proof $\Phi_m$ of $R(x_0,x_1), \ldots, R(x_{m},x_{m+1}) \vdash \Forward(x_0,x_{m+1})$, using only $\Forward R$-rule and $\id$. The definition is by induction on $m$. From this family of proofs we define a cut-free not regular $\infty$-proof $\Psi_m$ of $\Backward(x_0,x_1),  R(x_1,x_2), \ldots, R(x_{m},x_{m+1}) \vdash \Forward(x_0,x_{m+1})$. We define $\Psi_m$ by $\Backward L$ and $\eqL $ from $R(x_0,x_1),  R(x_1,x_2), \ldots, R(x_{m},x_{m+1}) \vdash \Forward(x_0,x_{m+1})$ and $\Backward(x_0,y),  R(y,x_1), R(x_1,x_2), \ldots, R(x_{m},x_{m+1}) \vdash \Forward(x_0,x_{m+1})$ for some fresh $y$. The first premise is $\Phi_m$. Up to renaming, the second premise is $\Psi_{m+1}$. Now we set $\Psi = \Psi_0$: this is an infinitary proof of $\Backward(x_0,x_1)\vdash \Forward(x_0,x_{1})$. There is a unique infinite path $\pi$ in $\Psi$, and it is the rightmost. In $\pi$, the unique ancestor of the inductive formula $\Backward(x_0,x_1)$ progresses at each step. Thus, the global trace condition is satisfied in $\Psi$. For each $m$, $\Psi$ includes the sequent $R(x_0,x_1), \ldots, R(x_{m},x_{m+1}) \vdash \Forward(x_0,x_{m+1})$, hence $\Psi$ includes infinitely many sequents of different  size and it is not regular.
$\Box$ \end{proof}

We will check now that any cut-free $\infty$-proof $\Pi$ of $\Backward(x,y) \vdash \Forward(x,y)$ is similar to the proof defined in Lemma \ref{lemma-cut-free-non-regular-proof} above, and not regular. Indeed, we will show that $\Pi$ for each $m \in \Nat$ first unfolds the inductive definition of $\Backward(x,y)$ to some $R(x,x_1), \ldots, R(x_m,y)$. Then $\Pi$ proves $\Forward(x,y)$. This operation cannot be done in \emph{constant memory space}: for any $m \in \Nat$, all atoms $R(x_1,x_2), \ldots, R(x_m,y)$ should be together in some left-hand side of some sequent of $\Pi$. The reason is: if, say, we forget the atom $R(x_1,x_2)$ and we unfold again $\Backward(x,y)$, there is no reason why we should get the same $R$-chain we got before. Instead of obtaining $R(x_1,x_2)$, we obtain some $R(x'_1,x'_2)$, which cannot be combined with  $\Forward(x,x_1)$ to prove $\Forward(x,y)$.

Thus, $\Pi$ has sequents of size $\ge m$ for all $m \in \Nat$ and $\Pi$ is not regular. As a corollary, we will deduce that being cut-free and being regular are incompatible for any $\infty$-proof of $\Backward(x,y) \vdash \Forward(x,y)$. 
\\

The first step in our proof is the following one. For every $m \in \Nat$ we define a canonical model $\setM_m$ of the language $L = \{R, \Forward, \Backward\}$.
%, with the maximum $R$-chain having length $m+1$. 
We set $\setM_m = [0,m+1] \times \Nat$, and we define $R_{\setM_m}((p,q),(p',q')) \Leftrightarrow (p+1=p')$ .
%in $\setM_m$. 
The model $\setM_m$ consists of $m+2$ vertical lines, two elements are $R$-related if they are in consecutive lines. The longest $R$-chain in $\setM_m$ has $m+2$ elements and $m+1$ edges. Two elements in the same vertical line cannot be distinguished by $R_{\setM_m}$. $\Forward_{\setM_m}$ and $\Backward_{\setM_m}$ are defined inductively from $R_{\setM_m}$. As p consequence, we have $\Forward_{\setM_m}((p,q),(p',q'))$ if and only if $p<p'$ if and only if $\Backward_{\setM_m}((p,q),(p',q'))$.

We introduce some proof terminology.
Whenever a proof $\Pi$ is fixed, we write $\nu \in \Pi$ for: ``$\nu$ is a node of $\Pi$''. We write $\Gamma_\nu \vdash \Delta_\nu$ for the sequent decorating $\nu$ in $\Pi$ and $\nu:\Gamma \vdash \Delta$ for $\Gamma = \Gamma_\nu$ and $\Delta = \Delta_\nu$. 
We call ``induction on $\nu$'' the induction on the distance of $\nu$ from the root of $\Pi$. If $\mu \in \Pi$, we write $\nu \le \mu, \nu < \mu$ for $\nu$ descendant of $\mu$, proper descendant of $\mu$ respectively. The root of $\Pi$ is the the lowest node.
%We write $\nu \ge \mu, \nu > \mu$ for $\nu$ ancestor of $\mu$, proper ancestor of $\mu$ respectively.  

Assume $\Pi$ is a cut-free $\infty$-proof of $\Backward(x,y) \vdash \Forward(x,y)$. By possibly renaming $\Pi$, we may assume that all eigenvariables of any $\nu \in \Pi$ are different from all eigenvariables of any $\mu \in \Pi$ with $\mu < \nu$. By induction on $\nu$ we prove that if $r$ is the rule of conclusion $\nu \in \Pi$, then $r=\Weak, \eqL, \Backward L, \Forward R$, or $r=\id$, with $\id$ inferring some $R(x,y) \in \Gamma_\nu \cap \Delta_\nu$, or $r=$ a substitution rule; and that $\Gamma_\nu$ consists of instances of $=,\Backward,R$, while $\Delta_\nu$ consists of instances of $\Forward$, $R$.
%12:20 15/09/2018

For all $m \in \Nat$, $\nu \in \Pi$ we define by induction on $\nu$ the \emph{canonical $\setM_m$-assignment} $\sigma_\nu : \FV(\Gamma_\nu \vdash \Delta_\nu) \rightarrow \setM_m$, with $\dom(\sigma_\nu)$  the set of all variables occurring in some node  $\mu \le \nu$. Informally, we set $\sigma_\nu(x) = (i,j)$ if the only values for $x$ ``relevant for the proof'' are the elements in the vertical line number $i$ of $\setM_m$. When all values of $x$ are ``irrelevant for the proof'', we arbitrarily choose some previously selected value as value for $x$. When $x$ is an eigenvariable of $\Pi$, we use the index $j$ (the position in the vertical line in $\setM_m$) in order to distinguish $x$ from all previous values in $\range(\sigma_\nu)$, whenever this is possible. 

\begin{definition}[Canonical assignment $\sigma_\nu$]
Assume $\nu _0 \in \Pi$ is the root. We set $\sigma_{\nu_0}(x) = (0,0)$ and $\sigma_{\nu_0}(y) = (m+1,0)$. Assume $\nu \in \Pi$, $\nu:\Gamma_i \vdash \Delta_i$ and $\nu$ is the premise number $i$ of some rule $r$ of $\Pi$, with conclusion $\mu : \Gamma \vdash \Delta$. Let $x \in \dom(\sigma_\nu)$.
\begin{enumerate}
\item
Assume $x$ is an eigenvariable of $r$. If $r=$ some substitution rule $\theta$, we set $\sigma_\nu(x) = \sigma_\mu(\theta(x))$. If $r = \Backward L$, then  $r$ infers $\Backward(a,b),\Gamma' \vdash \Delta'$ from $\Backward(a',c'),a=a',b=b',R(c',b'),\Gamma' \vdash \Delta'$, and from $R(a',b'),a=a',b=b',\Gamma' \vdash \Delta'$. The possible choices for $x$ are $a',b',c'$. We set $\sigma_\nu(a') = \sigma_\mu(a)$ and $\sigma_\nu(b') = \sigma_\mu(b)$. If $\sigma_\mu(b) = (i,j)$ then: 
\begin{enumerate} 
\item 
in the case $i>0$ we set $\sigma_\nu(c') = (i-1,j')$, where $j'$ is the number of $\Backward L$-rules in the path between $\nu$ and the root.
%In this case we have $\sigma_\nu(c')  \not \in \range(\sigma_\mu)$. 
\item
in the case $i=0$ we set $\sigma_\nu(c') =$ some value in $\range(\sigma_\mu)$ (say, $(0,0)$). 
\end{enumerate}
\item
Assume $x$ is not an eigenvariable of $r$. If $x\in\FV(\Gamma \vdash \Delta)$ we set $\sigma_\nu(x) = \sigma_\mu(x)$, otherwise we set $\sigma_\nu(x) =$ some value in $\range(\sigma_\mu)$ (say, $(0,0)$).
\end{enumerate}
\end{definition}

By induction on $\nu$, we may prove that the definition is correct and that $\dom(\sigma_\nu)$ equal to the set of all variables occurring in some $\mu \le \nu$. In the case $x$ is some  eigenvariable, by our convention on eigenvariables, $x$ does not occur in any node $\rho \le \mu$, therefore $x \not \in \dom(\sigma_\mu)$ and we are free to choose any value for $\sigma_\nu(x)$.

By case analysis, we have $\range(\sigma_\nu) = \range(\sigma_\mu)$, except in the case $\nu$ is the second premise of a $\Backward L$-rule, in which case we have $\range(\sigma_\nu) \setminus \range(\sigma_\mu) = \{\sigma_\nu(c')\}$. 

%21:20 17/09/2018
As usual in Model Theory, we define $\sigma_\nu, \setM_m \models R(a,b)$ as $R_{\setM_m}(\sigma_\nu(a), \sigma_\nu(b))$. We define $\models$ in the same way for the formulas: $a=b, \Forward(a,b), \Backward(a,b)$. We define $\sigma_\nu, \setM_m \models \Gamma$ as $\sigma_\nu, \setM_m \models A$ for all $A \in \Gamma$. We define $\Pi_1 \subseteq \Pi$ as the ``pruning of $ \Pi$ whose nodes have all assumptions true''. Formally, we set $\nu \in \Pi_1$ if and only if $\setM_m, \sigma_\mu \models \Gamma_\mu$ for every $\mu \le \nu$. $\Pi_1$ is a subtree of the proof $\Pi$, but it is not a sound proof: some rules in $\Pi_1$ have missing premises.

%19:45 22/09/2018
%13:21 20/09/2018

\begin{lemma}[$\Pi_1$-Lemma]
\label{lemma-pi-1}
\begin{enumerate}
\item
$\Pi_1$ includes: \emph{(a)} the root of $\Pi$; \emph{(b)} exactly one premise for each $\Backward L$-rule; and \emph{(c)} all premises of all other rules. 
\item
$\Pi_1$ is finite.
\end{enumerate}
\end{lemma}

\begin{proof}
\begin{enumerate}
\item
\begin{enumerate}
\item 
Assume $\nu_0$ is the root of $\Pi$. Then $\Gamma_{\nu_0}=\{\Backward(x,y)\}$ and $\setM_m, \sigma_{\nu_0} \models \Backward(x,y)$, hence $\nu_0 \in \Pi_1$. 
\item
Assume $\nu \in \Pi_1$ and $r$ is a rule with conclusion $\nu$. Assume $r=\Backward L$. Then $\nu:\Backward(a,b),\Gamma' \vdash \Delta'$. $\nu$ has first premise $\alpha:a=a',b=b',\Backward(a',c'), R(c',b'), \Gamma' \vdash \Delta'$ and second premise $\beta:a=a',b=b',R(a',b'),\Gamma' \vdash \Delta'$. For $\mu =\alpha, \beta$ we have $\sigma_\mu(a')=\sigma_\nu(a)$ and $\sigma_\mu(b')=\sigma_\nu(b)$. By the assumption $\nu \in \Pi_1$ we have $\sigma_\nu, \setM_m \models \Backward(a,b)$: then $\sigma_\nu(a)$ and $\sigma_\nu(b)$ are connected by some $R$-chain of length $l \ge 1$, and the same holds for $\sigma_\nu(a')$ and $\sigma_\nu(b')$. If $l=1$ then $\Backward(a',c'), R(c',b')$ are false in $\sigma_\alpha, \setM_m$, independently from the value of $\sigma_\alpha(c')$, while $\sigma_\beta, \setM_m \models R(a',b')$. Thus, $\alpha \not \in \Pi_1$ and $\beta \in \Pi_1$. If $l>1$, then $\Backward(a',c'), R(c,b')$ are true because $\sigma_\alpha(c')$ is defined as some $R$-predecessor of $\sigma_\alpha(b')$, while $R(a',b')$ is false in $\sigma_\beta, \setM_m$. Thus, $\alpha \in \Pi_1$ and $\beta \not \in \Pi_1$. 
\item
Assume $r \not = \Backward L$, hence $r$ is $\eqL, \Forward R, \id$ or is the substitution rule. By case analysis we deduce that for all premises $\mu:\Gamma_i \vdash \Delta_i$ we have $\sigma_\mu, \setM_m \models \Gamma_i$. 
\end{enumerate}

%21:45 17/09/2018
\item
\emph{$\Pi_1$ is a finite tree}. Assume $\nu \in \Pi_1$ and $\Backward(a,b) \in \Gamma_\nu$ and $\sigma_\nu(b) = (i,h)$. Any ancestor $\Backward(a',b') \in \Gamma_\mu$ of $\Backward(a,b) \in \Gamma_\nu$ is true in $\setM_m, \sigma_\mu$. Therefore in any progress of $\Backward(a,b)  \in \Gamma_\nu$ the first component $i$ of $\sigma_\nu(b)$ is reduced by $1$. Since $i \in [0,m+1]$, this may happen at most $m+1$ times. Thus, no path in $\Pi_1 \subseteq \Pi$ has an infinite progressing trace: by the global trace condition for $\Pi$, all paths in $\Pi_1$ are finite. From $\Pi_1$ binary and K\"{o}nig's Lemma we deduce that $\Pi_1$ is a finite tree.
\end{enumerate}
$\Box$
\end{proof}

By Lemma \ref{lemma-pi-1}.1.(b), the only rule left with two premises in $\Pi_1$ is $\Forward R$. 

We need some more proof terminology. If $\nu, \mu \in \Pi$, $\nu > \mu$, and an occurrence of $A$ in $\Gamma_\nu \vdash \Delta_\nu$ is an ancestor of an occurrence $B$ in $\Gamma_\mu \vdash \Delta_\mu$, we say that $A, B$ have ancestor distance $0$ if there is no rule between $\nu, \mu$ whose main formula is some ancestor of $B$. $A,B$ could still be syntactically different, because we may have $\eqL$-rules or substitution rules between $\nu, \mu$. 

Ancestor of distance $0$ in $\Pi_1$ have equal arguments in $\setM_m$. Indeed, assume that $p=R, \Forward, \Backward$. If $\nu$ is an ancestor of $\mu$ in $\Pi_1$ and $p(a',b') \in \Gamma_\nu$ is an ancestor of $p(a,b) \in \Gamma_\mu$ (or $p(a',b') \in \Delta_\nu$ is an ancestor of $p(a,b) \in \Delta_\mu$)  and $p(a',b'), p(a,b)$ have distance $0$, then $\sigma_\nu(a')=\sigma_\mu(a)$ and $\sigma_\nu(b')=\sigma_\mu(b)$. The proof is by induction on $\nu \in \Pi_1$. We use the following fact. All $\eqL$-rules whose conclusion is some $\rho \in \Pi_1$ have main formula some $a=b \in \Gamma_\rho$. By definition of $\Pi_1$ we have $\sigma_\rho, \setM_m \models a=b$. Thus, $\eqL$ replaces an argument of $p$ by an equal argument in $\setM_m$.

Another feature of $\Pi_1$ is that for any $\nu \in \Pi_1$, any $\Backward(a,b) \in \Gamma_\nu$ we have $\sigma_\nu(a) = (0,0)$, and for any $\Forward(a,b) \in \Delta_\nu$ we have $\sigma_\nu(b) = (m+1,0)$. The proof is by induction on $\nu$.

The main feature of $\Pi_1$ is that for all right-hand side $\Delta_\nu$ of a sequent of $\Pi_1$ we may extract from $\Pi_1$ some proof of some $A \in \Delta_\nu$, starting from assumptions in  $\Pi_1$, then using zero or more $\Forward$-productions. All steps of this proof are sound in $\setM_m$. Formally, for any $\nu \in \Pi$ we define a notion of \emph{production proof} for $A \in \Delta_\nu$. 

We say that $R(a,b) \in \Delta_\nu$ has  a production proof if some ancestor $R(a',b')$ of $R(a,b)$ is a main formula of some $\id$-rule. 

We say that $\Forward(a,b) \in \Delta_\nu$ has a production proof if there is some trace of $\Forward(a,b)$ ending in some main formula of an $\id$, such that for all rules $r=\Forward R$ whose main formula is in the trace: the trace chooses the auxiliary formula $\Forward(c',b')$ of the $r$ if any, and the auxiliary formula $R(a',c')$ of $r$ has a production proof. 


\begin{lemma}
\label{lemma-production}
Assume $\nu \in \Pi_1$.
\begin{enumerate}
\item
Assume $R(a,b) \in \Delta_\nu$ has a production proof. Then: $R(a,b)$ has some ancestor $R(a',b') \in \Delta_\mu$, $\mu \ge \nu$ such that $R(a',b') \in \Gamma_\mu$, and: $\sigma_\nu, \setM_m \models R(a,b)$.
\item
If $A \in \Delta_\nu$ has a production proof then $\sigma_\nu, \setM_m \models A$. 
\item
There is some $A \in \Delta_\nu$ with a production proof.
\end{enumerate}
\end{lemma}

\begin{proof}
\begin{enumerate}
\item
Assume $A = R(a,b) \in \Delta_\nu$. Then some ancestor $R(a',b') \in \Delta_\mu$ of $R(a,b)$ in $\Pi_1$ is a main formula of an $\id$-rule. Then $R(a',b') \in \Gamma_\mu$, hence $\sigma_\mu, \setM_m \models R(a',b')$.  $R(a',b')$ is a $0$-ancestor of $R(a,b)$, because the only rule with main formula $R(a',b')$ is $\id$, and there are no $\id$-rules below $\mu$ because the $\id$-rule has no premises. Thus, $\sigma_\nu, \setM_m \models R(a,b)$.

\item
By induction on $\nu \in \Pi_1$. 

\begin{enumerate}
\item
By point $1$ above.
\item
Assume $A = \Forward(a,b) \in \Delta_\nu$. There is some trace $\tau$  of $\Forward(a,b)$ ending in some main formula of an $\id$, such that for all rules $r = \Forward R$ whose main formula is in $\tau$: $\tau$ chooses the auxiliary formula $\Forward(c',b')$ of the $r$ if any, and the auxiliary formula $R(a',c')$ of $r$ has a production proof. The $\id$-rule has main formula some $R(c,d)$, therefore $\tau$ includes some $0$-ancestor $\Forward(a',b')$ of  $\Forward(a,b)$, which is the main formula of some $\Forward R$-rule. By $\Forward(a',b')$ $0$-ancestor  of  $\Forward(a,b)$ it is enough to prove $\sigma_\mu, \setM_m \Forward(a,b)$. 

Assume $\Forward R$ has a unique premise $\mu$, and auxiliary formula $R(a'',b'')$, with $a'=a''$, $b'=b''$ true in $\sigma_\mu$. By induction on $\mu$ we have $\sigma_\mu, \setM_m \models R(a'',b'')$, hence $\sigma_\mu, \setM_m \models \Forward(a',b')$, as wished. 

Assume $\Forward R$ has two premises $\alpha, \beta$, and auxiliary formulas $R(a'',c'')$ and $\Forward(c'',b'')$, with $a'=a''$, $b'=b''$ true in $\sigma_\alpha$ and $\sigma_\beta$. By induction on $\alpha, \beta$ we have $\sigma_\alpha, \setM_m \models R(a'',c'')$ and $\sigma_\beta, \setM_m \models \Forward(c'',b'')$, hence $\sigma_\mu, \setM_m \models \Forward(a',b')$, as wished. 
\end{enumerate}

\item
By induction on the height of a node $\nu \in \Pi_1$: the height exists  because $\Pi_1$ is finite. We prove the thesis for all leaves of $\Pi$, then we assume the thesis for all premises of a rule of $\Pi_1$ and we prove it for the conclusion of the rule.
\begin{enumerate}
\item
Assume $\nu$ is a leaf of $\Pi_1$. Then $\nu$ is an $\id$-rule and we choose $A = $ the atom $R(a,b)$ inferred by the $\id$-rule. 
\item
Assume $\nu$ is the conclusion of some $\Forward R$-rule in $\Pi_1$. Then $\nu: \Gamma \vdash \Forward(a,b),\Delta$ has premises $\alpha: \Gamma \vdash R(a,c),\Delta$ and $\beta:\Gamma \vdash \Forward(c,b),\Delta$. If there is some $A \in \Delta$ in $\alpha$ we a production proof we choose $A$. If not, and there is some $B \in \Delta$ in $\beta$ with a production proof we choose $B$. If not, by induction hypothesis on the trees of root $\alpha$ and $\beta$, both $R(a,c)$ and  $\Forward(c,b)$ have a production proof. Thus, $\Forward(a,b)$ has a production proof.
\item
Assume $\nu$ is the conclusion of some rule $\not = \id, \Forward R$ in $\Pi_1$. Then $r$ has a single premise $\mu$ in $\Pi_1$. We apply induction hypothesis and we choose the descendant $A'$ of the formula $A$ we have for $\mu$. 
\end{enumerate}
\end{enumerate}
$\Box$
\end{proof}


%21:18 22/09/2018


Let $\nu \in \Pi_1$ and $k \in \Nat$. The last ingredient we need for proving the non-regularity of $\Pi$ is the notion of $R$-chain in the node. We say that $a_1, \ldots, a_{k+1} \in \setM_m$ is an $R$-chain in the node $\nu$ if $a_1 \in \range(\sigma_\nu)$ and there are $R(x'_1, z'_1), \ldots, R(x'_k, z'_k) \in \Gamma_\nu$ such that $\sigma_\nu(x'_i)=a_i, \sigma_\nu(z'_i)=a_{i+1}$ for $i=1, \ldots, k$. 
%We say that $a_1, \ldots, a_{k+1} \in \setM_m$ is an $R$-chain above the node $\nu$ if $a_1 \in \range(\sigma_\nu)$ and there are $R(x'_1, z'_1) \in \Gamma_{\nu_1}, \ldots, R(x'_k, z'_k) \in \Gamma_{\nu_k}$ such that $\nu \le \nu_i$ and $\sigma_\nu(x'_i)=a_i, \sigma_\nu(z'_i)=a_{i+1}$ for $i=1, \ldots, k$. 
If $k=0$ the definition of $R$-chain in $\nu$ becomes: $a_1 \in \range(\sigma_\nu)$. 

If there is an $R$-chain of $k+1$ elements then $\setM_m, \sigma_\nu \models R(x'_i,z'_i)$ by definition of $\Pi_1$, therefore $a_1=(i,h_1), \ldots, a_{k+1}=(i+k,h_{k+1})$ are in consecutive vertical lines in $\setM_m$, hence pairwise distinct. Thus, the $k$ atoms $R(x'_i,z'_i)$ are pairwise distinct because they are distinguished by $\sigma_\nu$. If we are able to prove that for all $m \in \Nat$ there is an $R$-chain of $m+1$ elements in some $\nu \in \Pi$ then $\Pi$ is not regular.

%20:20 23/09/2018

\begin{lemma}[Right-extension]
\label{lemma-right-extension}
Assume that $\nu \in \Pi_1$ and for some $\mu \ge \nu$ there is some $R(a, b) \in \Gamma_\mu$ such that $(0,0) \not =\sigma_\mu(a) \in \range(\sigma_\nu)$. Then $(0,0) \not = \sigma_\mu(b)\in \range(\sigma_\nu)$ and there is some $R(a',b') \in \Gamma_\nu$ such that $\sigma_\nu(a') = \sigma_\mu(a)$ and $\sigma_\nu(b') = \sigma_\mu(b)$.
\end{lemma} 

\begin{proof}
Either $R(a,b)$ is the $0$-ancestor of some $R(a',b') \in \Gamma_\nu$ or $R(a,b)$ is the $0$-ancestor of some $R(a',b')$ auxiliary formula of some premise $\mu \ge \rho > \nu$ of some $\Backward L$-rule.  

Assume the first case. Then $\sigma_\nu(a') = \sigma_\mu(a)$ and $\sigma_\nu(b') = \sigma_\mu(b)$, hence $\sigma_\mu(b) \in \range(\sigma_\nu)$. From $\sigma_\mu, \setM_m \models R(a,b)$ we deduce that $\sigma(b)$ cannot be in the first vertical line of $\setM_m$, hence $(0,0) \not = \sigma_\mu(b)$. 

Assume the second case: we derive a contradiction. We argue by cases on $\rho$.
\begin{enumerate}
\item
Assume $\rho$ is the first premise $a=a',b=b', R(a',b'), \Gamma' \vdash \Delta'$ of $\Backward(a,b), \Gamma' \vdash \Delta'$. Then $\sigma_\mu(a') = \sigma_\rho(a') = \sigma_\rho(a) = (0,0)$ (this is a property of $\Pi_1$), contradicting the hypothesis.
\item
Assume $\rho$ is the second premise $a=a',b=b', c=c', \Backward(a',c'), R(c',b'), \Gamma' \vdash \Delta'$ of $\Backward(a,b), \Gamma' \vdash \Delta'$. By definition of $\sigma_{\rho}$ we should have $\sigma_\mu(c') = \sigma_\rho(c') \not \in \range(\sigma_\nu)$, contradicting the hypothesis.
\end{enumerate}
$\Box$
\end{proof}


%\begin{lemma}[Right-extension of $R$-chains in a node]
%\label{lemma-right-extension}
%Assume that $\nu \in \Pi_1$ and $k \in \Nat$ and $a_1, \ldots, a_{k+1}$ is an $R$-chain in $\nu$ and $a_1 = (i,h)$ for some $i>0$, $h$. Assume that for some $\mu \ge \nu$ there is some $R(x''_{k+1}, z''_{k+1}) \in \Gamma_\mu$ with $\sigma_\mu(x'_{k+1})=a_{k+1}$ and $\sigma_\mu(z'_{k+1})=a_{k+2}$. Then we may extend the $R$-chain in $\nu$ to $a_1, \ldots, a_{k+2}$.
%\end{lemma} 
%
%\begin{proof}
%
%Either $R(x''_{k+1},z''_{k+1})$ is the ancestor of some $R(x'_{k+1},z'_{k+1}) \in \Gamma_\nu$ and we are done, or $R(x''_{k+1},z''_{k+1})$ is the ancestor of some $R$-atom  inserted by an $\Backward L$-rule $r$ between $\nu$ and $\mu$.  We prove that this cannot be. We argue by cases on $r$.
%\begin{enumerate}
%\item
%Assume $r$ is: $a=a',b=b', R(a',b'), \Gamma' \vdash \Delta'$ infers $\Backward(a,b), \Gamma' \vdash \Delta'$. Then $\sigma_\mu(a) = (0,0)$ (this is a property of $\Pi_1$), hence $a_{k+1} = \sigma_\mu(a') = \sigma_\mu(a) = (0,0)$. However, by assumption on the $R$-chain we have $a_{k+1} = (i+k,h_{k+1})$ hence $0 = i+k \ge i > 0$, contradiction.
%\item
%Assume $r$ is: $a=a',b=b', c=c', \Backward(a',c'), R(c',b'), \Gamma' \vdash \Delta'$ infers $\Backward(a,b), \Gamma' \vdash \Delta'$. By definition of $\sigma_{\mu}$ we  would have $\sigma_\mu(c') \not \in \range(\sigma_\nu)$. Instead, we have $a_{k+1} \in \range(\sigma_\nu)$. Indeed, if $k=0$ then $a_{k+1} = a_1 \in \range(\sigma_\nu))$ and if $k>0$ then $a_{k+1} = \sigma_\nu(z'_{k})$. Thus, $\sigma_\mu(c') \not = a_{k+1}$. This means that  $R(x''_{k+1},z''_{k+1})$ is no ancestor of $R(c',b')$, because the assignment $\sigma_\mu$ makes the first arguments of the two atoms different. 
%\end{enumerate}
%$\Box$
%\end{proof}

We prove that there is no cut-free cyclic proof $\Pi$ for $\Backward(x,y) \vdash \Forward(x,y)$.

\begin{theorem}
Assume $\Pi$ is a cut-free $\infty$-proof of $\Backward(x,y) \vdash \Forward(x,y)$. For all $m \in \Nat$, there is some $\nu \in \Pi$ and some $m$ atoms $R(x'_1,z'_1), \ldots, R(x'_{m},z'_{m}) \in \Gamma_\nu$ pairwise distinct. In particular, $\Pi$ is not regular.
\end{theorem}

\begin{proof}
$\Forward(x,y)$ is the unique formula in the right-hand side of the root of $\Pi_1$. Thus, by Lemma \ref{lemma-production}, $\Forward(x,y)$ has a production proof in $\Pi$. We unfold the definition of production proof. There are $m+1$ nodes $\nu_0 < \ldots < \nu_{m}$ in $\Pi_1$, all conclusions of $\Forward R$-rules. For all $0 \le k<m$, the node $\nu_k:\Gamma_k \vdash \Forward(x_k,y_k), \Delta'_k$ has first premise $\alpha_k:\Gamma_k \vdash R(x_k,z_k), \Delta'_k$ and second premise $\beta_k:\Gamma_k \vdash \Forward(z_k,y_k), \Delta'_k$. $\Forward(x_0,y_0)$ is a distance $0$ ancestor of $\Forward(x,y)$ and $\Forward(x_{k+1},y_{k+1})$ is a distance $0$ ancestor of $\Forward(x_k,y_k)$. The node $\nu_{m}:\Gamma_m \vdash \Forward(x_m,y_m)$ is the conclusion of a $\Forward L$-rule with a single premise $\rho = \alpha_m : \Gamma_m \vdash R(x_{m},y_{m}), \Delta'_k$. For all $k<m$, the formulas $R(x_k,z_k)$ and $\Forward(x_k,y_k)$ have a production proof, hence they are true in $\sigma_{\nu_k}$ and $\setM_m$. By the distance $0$ ancestor relation we have $(0,0) = \sigma_\rho(x) = \sigma_\rho(x_0)$ and $\sigma_\rho(z_k) = \sigma_\rho(x_{k+1})$ for all $k<m$ and $\sigma_\rho(y_k) = \sigma_\rho(y) = (m+1,0)$ for all $k \le m$. Thus there are $a_0, \ldots, a_m \in \setM_m$ which form an $R$-chain in $\setM_m$. Thus, $a_k=(k,h_k)$ for all $k\le m+1$ and $\sigma_\rho(x_k) = a_k$, $\sigma_\rho(y_k) = a_{m+1}$ for all $k \le m$. 

We claim that there is an $R$-chain in $\nu_0$. The proof is by induction on $k \in [0,m]$. Let $\tau = \sigma_{\nu_0}$ be the canonical assignment for the node $\nu_0: \Gamma_0 \vdash \Forward(x_0, y_0), \Delta'_0$. 
%We claim that that $a_1 \in \range(\tau)$ and if $k=1, \ldots, m$, there is some $R(x'_k,z'_k) \in \Gamma_0$ such that $\tau(x'_k)=a_k$, hence there are $m$ distinct formulas in $\Gamma_0$. This implies our thesis. 


\begin{enumerate}
\item
Let $k=0$: then $a_{1} = \tau(z_0) \in \range(\tau)$. We already checked that $a_1 = (1,h_1) \not = (0,0)$. 
\item
Assume the thesis for $0 \le k < m$, in order to prove it for $k+1$. By inductive assumption $(0,0) \not = a_{1} \in \range(\tau)$ and and there are $R(x'_1, z'_1), \ldots, R(x'_k, z'_k) \in \Gamma_{\nu_0}$ such that $\sigma_{\nu_0}(x'_i)=a_i, \sigma_{\nu_0}(z'_i)=a_{i+1}$ for $i=1, \ldots, k$. We already proved that there is some $R(x_{k+1},z_{k+1}) \in \Gamma_\mu$ in some $\mu \ge {\nu_0}$ such that $\sigma_{\mu}(x_{k+1})=a_{k+1}, \sigma_\mu(z_{k+1})=a_{k+2}$. We apply the Right Extension Lemma \ref{lemma-right-extension}and we deduce that there is some $R(x'_{k+1}, z'_{k+1}) \in \Gamma_{\nu_0}$ such that $\sigma_{\nu_0}(x'_{k+1})=a_{k+1}, \sigma_{\nu_0}(z'_{k+1})=a_{k+2} \not = (0,0)$.  
%Some ancestor $R(x''_{k+1},z''_{k+1})$  of the atom  $R(x_{k+1},z_{k+1})$ is deduced by some $\id$-rule above $\nu_{k+1}$. Thus, either $R(x''_{k+1},z''_{k+1})$ is the ancestor of some $R(x'_{k+1},z'_{k+1}) \in \Gamma_0$ or $R(x''_{k+1},z''_{k+1})$ is the ancestor of some atom $R(c',b')$ inserted by an $\Backward L$-rule between $\nu_0$ and $\nu_{k+1}$. 
%In this case, by definition of $\sigma_{\nu_{k+1}} \subseteq \sigma_\rho$ we  have $\sigma_\rho(c') \not \in \range(\tau)$, and in particular $\sigma_\rho(c') \not = a_{k+1} = \tau(x_{k+1})$. 
%This means that  $R(x''_{k+1},z''_{k+1})$ is no ancestor of $R(c',b')$, because there is a substitution making the first arguments of the two atoms different. 
%We conclude that there is some $R(x'_{k+1},z'_{k+1}) \in \Gamma_0$ such that $\tau(x'_{k+1})=a_{k+1}$ and $\tau(z'_{k+1})=a_{k+2}$, hence $a_{k+2} = \sigma_\rho(z'_{k+1}) \in \range(\tau)$. 
\end{enumerate}
We conclude that there are $m$ formulas in $\Gamma_0$ whose first arguments may be made pairwise distinct by the assignment $\sigma_{\nu_m}$. Thus, there are $m$ pairwise distinct formulas in $\Gamma_0$.
$\Box$ \end{proof}
%14:47 15/09/2018
 
%Let $\Gamma'_h \vdash \Delta'_h = \bigcup_{h \le l} \Gamma_h \vdash \Delta_h $ be the union of all sequent of $\pi$ up to the sequent number $k$. Assume there are $n$ $\Backward L$-rules from the root to the sequent number $k$. Then $\Gamma'_k$ consists of: one set of equations $a=b$ between $3n+2$ variables defining $n+2$ equivalence classes, and one set of atoms $R(a,b)$ with $b \not \equiv y$, and one set of atoms $\Backward(x,a)$. $\Delta'_k$ consists of a set atoms $\Forward(x,y)$ and $\Forward(c,y)$, with the variables $c$ forming a disjoint set. 



\section{Induction and failure cut}
%10:10 15/07/2018

In this section we prove that we cannot interpret (full) induction rule $\Ind$ without cuts, even for the induction rule for the predicate $\Nat t$, and we propose $\Ind R$, a restriction of induction which is conditionally derivable without cuts. 

\begin{lemma}[Induction and cut]
\label{lemma-induction-cut}
There is a cut-free proof of $E_{2^x}, \Nat x \vdash \Nat 2^x$ in cyclic proofs $+$ induction rule, and no cut-free proof in cyclic proofs.
\end{lemma}

\begin{proof}
Indeed, there is a cut-free proof of $E_{2^x} \vdash \Nat 2^0$: we use $\eqL$ on $2^0=1$ and $E_{2^x} \vdash \Nat 1$, with $1 \equiv S0$. We prove $E_{2^x} \vdash \Nat S0$ from $\Nat R$ for $S$ and $E_{2^x} \vdash \Nat 0$, then $\Nat R$ for $0$. There is a cut-free proof of $E_{2^x}, \Nat x \vdash \Nat 2^{Sx}$, we use $\forall L$ on $\forall x.2^{Sx}=2^x + 2^x$ and $E_{2^x}, 2^{Sx}=2^x + 2^x, \Nat x \vdash \Nat 2^{Sx}$, then $\eqL$ and $E_{2^x}, \Nat x, \Nat 2^x  \vdash \Nat (2^{x} + 2^x)$. This latter has a cut-free cyclic proof by Lemma \ref{lemma-sum-total}, point $3$. We already proved that there is no cut-free proof of $E_{2^x}, \Nat x \vdash \Nat 2^x$, this impliess that cut-free proofs do not conditionally derive $\Ind$-rule.
$\Box$ \end{proof}

We analyze our cyclic proof of $E_{2^x}, \Nat x \vdash \Nat 2^x$. Implicitly, the proof receives a value of $x$ and computes a value of $ 2^x$ in unary notation (with $0,S$ only): this is a basic example of non-polynomial computation. We assume that cuts are required by proofs hyding non-polynomial computations, and we study how to prevent this situation.

The proof of $E_{2^x}, \Nat x \vdash \Nat 2^x$ includes a proof $E_{2^x}, \Nat 2^x \vdash  \Nat (2^{x} + 2^x)$, in which we derive the assumption $ \Nat 2^x$ twice with some $\Nat L$-rule. This sub-proof defines a sub-recursion computing the value of $(2^{x} + 2^x)$ from the value of $2^x$. The point is that the value of $2^x$ is the very value being recursively computed by our proof. It is known that, in order to prevent non-polynomial computation, \emph{``the depth of sub-recursions \ldots cannot depend on the value being recursively computed''} (\cite{Bellantoni-Cook}, \S 2). This suggests to forbid any $\Nat L$-rule whose main formula is the induction hypothesis $\Nat 2^x$. We allow a $\Nat L$-rule inferring the assumption $\Nat x$. 
%The role of formulas $\Nat x$, $\Nat 2^x$ derived, respectively, with $\Nat L$ and $\Nat R$ are similar, respectively, to the role of a ``safe value'' and of a  ``normal value'' in \cite{Bellantoni-Cook}.

More in general, to prevent the need of cuts, we define a weaker rule $\Ind R$ in which the assumption $F[x]$ of the induction step cannot be the main formula of an $L$-rule. $\Ind R$ may be interpreted in cyclic proofs without adding new cuts. Our impression is that this is the way we argue by induction in a cut-free cyclic proof, in general. 

Formally, an $R$-proof of $\Gamma, \Nat x, F[x] \vdash F[Sx], \Delta$ for $F[x]$ is a cyclic proof in which we forbid the use of $L$-rules with main formula $F[x]$. 
%We consider $F[x]$ the conclusion of an $\eqL$-rule if $F[x,a,b]$ is inferred from $F[x,b,a]$ and $F[x,a,b] \not \equiv F[x,b,a]$. 
The only rule allowed to infer $F[x]$ is $\id$, then we may apply $R$-rules from the right occurrence of $F[x]$ in the $\id$-rule.

We define the restriction $\Ind R$ of induction rule  in which the second premise $\Gamma, \Nat x, F[x] \vdash F[Sx], \Delta$ in an $R$-rule for $F[x]$.  


\begin{lemma}[$\Ind R$ is conditionally derivable in cut-free proofs]
\label{lemma-ind-R}
\begin{enumerate}
\item
There is an $R$-cyclic proof $\Pi_1$ of  $\Gamma, \Nat x, F[x] \vdash F[Sx], \Delta$ $\Leftrightarrow$ there is a cyclic proof $\Pi$ of  $\Gamma, \Nat x \vdash F[Sx], \Delta$ with free assumption $\Gamma, \Nat x \vdash F[x]$ and $\Nat x$ in the free assumption ancestor of $\Nat x$ in the conclusion. 

Besides, $\Pi$ is cut-free $\Leftrightarrow$ $\Pi_1$ is cut-free. 
\item
Cut-free proofs conditionally derive the rule $\Ind R$.
\end{enumerate}
\end{lemma}

\begin{proof}
\begin{enumerate}
\item
$\Leftarrow$. Assume we have a cyclic proof $\Pi_1$ of  $\Gamma, \Nat x, F[x] \vdash F[Sx], \Delta$ in which the ancestors of $F[x]$ are infereed by no $L$-rule (including no $\eqL$-rule). We may define a proof $\Pi$ of of  $\Gamma, \Nat x \vdash F[x], \Delta$ with free assumption $\Gamma, \Nat x \vdash F[x]$, such that $\Nat x$ in the free assumption ancestor of $\Nat x$ in the conclusion, and $\Pi$ is cut-free if $\Pi_1$ is cut-free. 

We first rename all eigenvariables $y$ in $\Pi_1$ in order to have $y \not \in \FV(\Gamma, \Nat x \vdash \Delta)$, then we replace each node $\Gamma' \vdash \Delta'$ with $\Gamma, \Nat x, \Gamma' \vdash \Delta, \Delta'$, with $\Gamma, \Nat x, F[x] \vdash F[Sx], \Delta$ replaced by itself. All $\Nat x$ we add are an ancestor of $\Nat x$ in the conclusion.

%We replace each rule of $\Pi$ with a repeating rule, for instance we prove $a=b,\Theta[a,b] \vdash \Xi[a,b]$ from $a=b,\Theta[b,a], \Theta[a,b] \vdash  \Xi[a,b], \Xi[b,a]$. We skip all $W$-rules. In this way we add a finite set of formulas $I$ with $i$ elements, some sub-set $J \subseteq I$ for each sub-proof. Therefore the number of sub-proofs may increase by a factor $2^i$, but stays finite. 

We define in this way a proof $\Pi_2$ with at most the same number of subproofs, hence regular. For each infinite path $\pi_2$ of $\Pi_2$ there is some infinite path $\pi_1$ of $\Pi_1$. There is some infinitely progressing trace $\tau$ in $\pi_1$. $\tau$ is in $\pi_2$ and infinitely progressing, therefore $\Pi_2$ satisfies the global trace condition.

There are finitely many $\id$-rule for an ancestor of $F[x]$ in $\Pi_2$, say $\Gamma, \Nat x, \Gamma_i, F[x] \vdash F[x], \Delta, \Delta_i$ for $i=1, \ldots, n$: these are the only rules inferring some ancestor of $F[x]$. The ancestors of $F[x]$ have no progress because we forbid $L$-rules for $F[x]$. Therefore we may remove all ancestors of $F[x]$ and we remove no progressing trace and the global trace condition is still true. The number of sub-proof may duplicate because some $F[x]$ on the left-hand side are ancestors and removed and some are not ancestors and not removed. Thus, regularity is preserved. We defined a proof $\Pi_3$ with open assumptions $\Gamma, \Nat x, \Gamma_i \vdash F[x], \Delta, \Delta_i$ for $i=1, \ldots, n$, and with $\Nat x$ ancestor of $\Nat x$ in the conclusion. We prove each open assumption with $\Weak$ from new open assumptions $\Gamma, \Nat x \vdash F[x], \Delta$, while preserving regularity and global trace. This is a cyclic proof $\Pi = \Pi_4$ of $\Gamma \vdash F[x], \Delta$ with open assumption $\Gamma, \Nat x \vdash F[x], \Delta$, and $\Nat x$ ancestor of $\Nat x$ in the conclusion. In $\Pi$ we inserted some $\Weak$-rules but no cut-rule, therefore if $\Pi_1$ is cut-free $\Pi$ is cut-free.

$\Rightarrow$. Assume the converse. We rename all eigenvariables $y$ of $\Pi$ in order to have $y \not \in \FV(F[x])$, then we add $F[x]$ everywhere in the proof, and we replace open assumptions $\Gamma, \Nat x, F[x] \vdash F[x], \Delta$ with the $\id$-rule. 

\item
Assume we have a cyclic proof $\Pi_0$ of $\Gamma, \vdash F[0] \Delta$ and $\Pi_1$ of  $\Gamma, \Nat x, F[x] \vdash F[Sx], \Delta$ in which the ancestors of $F[x]$ use no $L$-rule. Let $\Pi$ as in the previous point. From $\Pi_0$ and $\Pi[x'/x]$ with $x' \not \in \FV(\Gamma, \Nat x \vdash F[x], \Delta)$ and the rule $\Nat L$ we may define a cyclic proof of $\Gamma, \Nat x \vdash F[x], \Delta$ with a rule $\Nat L$ in the root inferring $\Nat x$, and free assumption $\Gamma, \Nat x' \vdash F[x']$. 
$\Nat x$ in the root is the conclusion of a $\Nat L$-rule, hence there is a progress from $\Nat x$ to $\Nat x'$ in the premise. $\Nat x'$ is a descendant of $\Nat x'$ in each free assumption $\Gamma, \Nat x' \vdash F[x']$, therefore we have a progress from $\Nat x$ in the root to $\Nat x'$ in each free assumption. 

If we repeat the proof from the free assumption, choosing alternatively $\Pi$ and $\Pi[x'/x]$, we define a regular proof $\Pi'$. $\Pi'$ satisfies the global trace condition, because all infinite paths either are definitively in $\Pi_0$ or in some $\Pi$ or $\Pi[x'/x]$ and include an infinitely progressing trace, or they pass through infinitely many $\Pi$ and $\Pi[x'/x]$, and in this case $\Nat x$ in the root is infinitely progressing. Thus, $\Ind R$ is conditionally derivable without adding cuts, and cut-free proofs conditionally derive $\Ind R$.
\end{enumerate}
$\Box$ \end{proof}
 
 
\section{Cyclic proofs: substitution and renaming}
In order to prove a restricted version of cut-elimination, in this section we quickly recall some features of cyclic proofs: substitution and renaming rules.

A \emph{substitution} is a map from finitely many first order variables to first order terms. A substitution is a variable substitution if it maps variables into variables, and a \emph{renaming} if it is a bijection on variables. Substitutions are extended to terms/formulas/proof tree by compatibiliy, with renaming of bound variables to prevent variable capture. 

We call an \emph{infinitary binary tree} any possibly infinite tree in which each node has at most two children. We consider cyclic proofs as infinitary binary proof-trees for sequent calculus. The first group of rules are the usual rules for sequent calculus: identity rule, weakening, exchange, contraction, cut, left/right introduction for each logical connnective: $\vee, \wedge, \rightarrow, \forall, \exists$. As usual we often skip exchange and contraction rule, and we assume that left/right introduction and cut are merged with exchange and contraction: the inferred formulas and the assumptions formulas may be in any position in their side of the sequent, and inferring $A$ may or may not include contraction on $A$. 

Peano Arithmetic and Martin-L\"{o}f theory have more rules, induction rule, injectivity axioms for constructors and basic equality axioms, and in the case of Peano Arithmetic basic arithmetical axioms. The induction axiom for $\Nat$ and a formula $F$ is $F[0/x], \forall x. (\Nat x, F \rightarrow F[Sx/x] \vdash F)$ for any $F$. 

Cyclic proofs have different rules in their place:
\begin{enumerate}
\item 
\emph{Substitution rule}. From $\Gamma \vdash \Delta$ and $\sigma$ substitution we deduce $\sigma(\Gamma \vdash \Delta)$. 
\item
\emph{Renaming rule}.  From $\Gamma \vdash \Delta$ and $\sigma$ renaming we deduce $\sigma(\Gamma \vdash \Delta)$. (This is a restriction of Substitution rule).
\item
$(\eqL)$.  From $\Gamma[a,b] \vdash \Delta[a,b]$ deduce $a=b,\Gamma[b,a] \vdash \Delta[b,a]$. 
\item
$(\eqR)$. $\Gamma \vdash \Delta, t=t$.
\item 
$(\Nat L)$. From $\Gamma, t=0 \vdash \Delta$ and $\Gamma, t=Sx, \Nat x \vdash \Delta$ and $x$ fresh ($x \not \in \FV(\Gamma, t, \Delta)$) deduce $\Gamma, \Nat t \vdash \Delta$.
\item
$(\Nat R)$: $\Gamma \vdash \Nat 0$ and if $\Gamma \vdash \Delta, \Nat u$ then $\Gamma \vdash \Delta, \Nat S u$. 
\end{enumerate}
We add a left and right rule \emph{for each inductive definition}. The fresh variables in the rules $\Nat L$, $\exists L$, $\forall R$ are the eigenvariables in the proof.  Recall that induction is \emph{not} assumed in cyclic proofs. 

Assume $c$ is some atomic predicate. We distinguish three sets of atomic predicates: logical (given with the first order languague), equality ($=$), and inductively defined predicates. We assume identity rule for logical predicates and we consider it both as a left-rule and a right-rule. We consider the identity rule as the introduction rule for logical predicates. We may derive identity for equality and for inductively defined predicates, in a cut-free way. Here are two examples. 
\begin{enumerate}
\item
\emph{Identity for equality}.
from $\vdash t=t$ (rule $= \! \!  \!  R$) we derive $t=u \vdash t=u$ (rule $\eqL$).
\item
\emph{Identity for $\Nat$}. 
It is enough to prove that we may derive derive $\Nat x \vdash \Nat x$ from $\Nat y \vdash \Nat y$ for any $y \not \equiv x$. We reason by cases with the $\Nat L$ rule. First, we derive $\Nat x, x=0 \vdash \Nat x$ from $\Nat x \vdash \Nat 0$ (rule $\eqL$), then from rule $\Nat L$. Second, for any $y \not \equiv x$ we derive $\Nat y, x = Sy \vdash \Nat x$ from $\Nat y \vdash \Nat Sy$ (rule $\eqL$), then from $\Nat y \vdash \Nat y$ (rule $\Nat R$), as wished.
\end{enumerate}
We derive (cut-free) the identity rule for compound formulas from identity for atomic formulas.

%We consider one extra rule which we call the fresh variable rule: ``if $x=u, \Gamma \vdash \Delta$ and $x \not \in \FV(u, \Gamma, \Delta)$ then $\Gamma \vdash \Delta$.  If $x \not \in \FV(u)$ then $(x=u)[u/x] = (u=u)$, therefore $\vdash \exists x.x=u$ is provable from $\vdash u=u$, which is an axiom. Thus, the fresh variable rule is derivable  by deducing $\exists x. x=u, \Gamma \vdash \Delta$ from $x=u, \Gamma \vdash \Delta$ by $\exists L$, then $\Gamma \vdash \Delta$ with a cut with $\vdash \exists x.x=u$.

Regular infinitary binary trees have a finite graph representation. Assume $\Pi$ is regular, with at most $n$ subtrees. Thus, in each branch of $\Pi$ there is a first node, among the first $n+1$ nodes, which is the root of a subtree equal to the subtree of a previous node. If we include a back link from these nodes to the unique ancestor which is the root of an equal tree we obtain a finite graph representation of the infinite tree $\Pi$. There are graph representation with exactly $n$ nodes.

Assume $\Pi$ is some infinitary binary proof tree. Assume $A = p(\vec{t})$ has an inductive definition. A \emph{trace} of an occurence of $A$ in the left-hand side of some sequent of $\Pi$ is any possibly list $\tau$ of ancestors of $A$ in $\Pi$. A trace $\tau$ progresses whenever an element $p(\vec{t})$ of $\tau$ is the assumption formula of a case rule. The global trace condition for $\Pi$ is: for any infinite branch $\beta$ of $\Pi$ there is some infinitely progressing trace $\tau$ included in $\beta$. 

Assume $\Pi$ is a regular infinitary proof tree and $G$ is its graph representation. Then the global trace condition is decidable for $\Pi$: we sketch how.  We claim without proof that global trace condition is equivalent to:  every path $\pi$ from a node $\alpha$ to itself in $G$ has some power $\pi^n$ including a trace $\tau$ progressing from some $A$ in $\alpha$ to itself. There are infinitely many path from a node $\alpha \in G$ to itself when $G$ is cyclic, but if we quotient them according to the trace relationship they define among inductively defined formulas $A$ in the left-hand side of a sequent, they are reduced to finitely many. 
Thus, the global trace condition is decidable for regular trees.

%12:07 06/06/2018
Cyclic proofs are infinitary binary trees $\Pi$ which are regular and satisfy the global trace condition. Finitely many subtrees implies finitely many sequents and finitely many variables, including the eigenvariables of the proof.

From now on, we consider cyclic proofs without the substitution rule: they are enough to interpret the induction axiom, and the substitution rule is derivable from the other rules if we have cuts. The ordinary cut-elimination procedure applies substitutions on proofs. On infinite proofs, the cut-elimination procedure has to be repeated infinitely many times, cyclically, and if the proof substitutions compose, then they may define infinitely many different formulas and a non-regular proof-tree. We ignore whether this happens, but this could be the reason why cyclic-proofs have no cut-elimination, even if they lack the induction rule. In order to preserve regularity, we will  define a cut-elimination process involving a finite set of substitution $\setS = \{[u_1/x_1], \ldots, [u_n/x_n]\}$ on proofs with $\vec{x} \cap \FV(\vec{u}) = \emptyset$: in this way the closure by composition of $\setS$ is finite. 

The induction statement has cut-free cyclic proofs, both with and without renaming. With renaming we may obtain shorter proofs for the induction statement (and shorter proofs  in general).

%There are $\Sigma^0_1$ formulas $A$, having a cyclic proof including $\Sigma^0_1$-formulas only, for which there are no normal cyclic proofs. This rules out the cut-elimination for cyclic proofs, at least when we cut formulas with quantifiers.

Substitution rule may be replaced everywhere by a combination of the rules for $\rightarrow, \forall$ and cut, while preserving regularity and global trace condition. From now on we skip the substitution rule in cyclic proofs, and we only consider the renaming rule.


We may prove that cyclic proofs are closed under a substitution $[t/x]$ on proofs. This property is more difficult to prove than expected. Unlike the case of finite proofs, the number of different subtrees of the infinitary proof may double after substitution, because the substitution $[t/x]$ does not propagate to \emph{all} occurrences of $x$ in the proof. Indeed, $x$ may disappear by a weakening rule, then $x$ be inserted as a eigenvariable in the proof, and in this case $x$ cannot be substituted. Apparently, closure under substitution is not required for eliminating quantifier-free cuts, therefore we skip it.

We define a renaming procedure $\Pi_2 = \Theta(\Pi_1)$ of a cyclic proof $\Pi_1$. $\Theta$ applies one renaming substitution $\psi$ to each sequent of in $\Pi_1$, with two provisos. 
\begin{enumerate}
\item
In the case of a renaming rule we require commutation of renaming: if we deduce $\phi_1(\Gamma_1 \vdash \Delta_1)$ from $\Gamma_1 \vdash \Delta_1$ in $\Pi_1$ and $\phi_2(\Gamma_2 \vdash \Delta_2)$ from $\Gamma_2 \vdash \Delta_2$ in in $\Pi_2$,  then we require two renaming maps $\psi_1$ on $\phi_1(\Gamma_1 \vdash \Delta_1)$ and $\psi_2$ on $\Gamma_1 \vdash \Delta_1$ such that $\psi_1 \phi_1 = \phi_2 \psi_2$. 
\item
In any other case of two consecutive nodes of in $\Pi_1$ if we apply two renaming $\psi_1$, $\psi_2$ then $\psi_1$, $\psi_2$ coincide on the variables common to the two sequents. 
\end{enumerate}
Renaming preserves being an infinitary proof, regularity and the global trace condition. There is a canonical renaming for infinitary proofs with substitution rule restricted to renaming. Renaming may increase exponentially the number of different subtrees.

\begin{lemma}[Canonical Renaming ]
\label{lemma-renaming}
Assume $\Pi$ is an infinitary binary proof with conclusion $\Gamma \vdash \Delta$, and $V = \FV(\Gamma \vdash \Delta)$, and $\phi:V \rightarrow W$ is a bijection on variables. Assume $\Pi$ has substitution rule restricted to renaming. Then there is some renaming $\Theta_\psi$ such that $\Pi_0 =\Theta_\psi(\Pi)$ has conclusion $\phi(\Gamma \vdash \Delta)$ and substitution rule restricted to the identity. If $\Pi$ has the global trace condition (is regular) then $\Pi_0$ has the global trace condition (is regular).
\end{lemma}

\begin{proof}
We define $\Pi_0$ by applying $\phi$ through $\Pi$. Whenever we find some renaming rule: ``from $\Theta \vdash \Xi$ deduce $\psi(\Theta \vdash \Xi)$'', then we replace $\phi$ with $\phi\psi$ from the node $\Theta \vdash \Xi$, and we associate the renaming rule to $\id$. 
Whenever we find some rule $r$ with conclusion $\Theta \vdash \Xi$ and one assumption $\Theta' \vdash \Xi'$ with a variable $x$ bound in the proof, we replace $\phi$ with $\psi$. $\psi$ is obtained by restricting $\phi$ to $\FV(\Theta, \Xi)$, then adding $y/x$, with $y$ the first fresh variable, i.e., the first $\not \in \FV(\Theta, \Xi)$. 
Global trace condition is preserved because trace and progress are not affected by renaming. Assume $\Pi$ is regular. Then for some $n$ each sequent in $\Pi$ has at most $n$ free variables: $n$ is the maximum among the number of variables in the roots of the finitely many subtrees of $\Pi$. 
Since each node of $\Pi_0$ is some renaming of some node of $\Pi$, each sequent in $\Pi_0$ has at most $n$ free variables. The the index $i$ of a eigenvariable $x$ in $\Pi_0$ is at most $n-1$. Indeed, when we create a eigenvariable, we assign to it the first index not in use, and there are at most $n-1$ variables different from the eigenvariable, therefore the first first index not in use is some $i \le n-1$. 
Thus, the variables of $\Pi_0$ are among $\FV(\phi(\Gamma \vdash \Delta))$ or the initial segment $\{x_0, \ldots, x_{n-1}\}$ of $X$, hence are finitely many. 
Each subtree $\Pi_1$ of $\Pi_0$ is $\Theta_\psi(\Pi_2)$ for some subproof $\Pi_2$ of $\Pi$, and since we have finitely many $\Pi_2$ in $\Pi$, and finitely many variables in $\Psi_0$, then we have finitely many choices for $\psi$, hence finitely many $\Theta_\psi(\Pi_2)$ in $\Pi_0$. 
$\Box$ \end{proof}

In the last line of the proof above, the number of subtrees may raise by a factor $n!$, with $n$ the maximum number of free variable in a sequent. We say that two infinitary proofs are equivalent if they may be obtained one from the other by: renaming, adding/removing any number of identical renaming rules which are below a rule different from renaming. The previous lemma says that a cyclic proof is equivalent to a proof having identical renaming only. 


All identical renaming may be removed according to the definition of proof equivalent. Indeed, the only case in which this would not be possible is a path in which all rules from some point on are identical renaming. But this cannot be: by the global trace condition every renaming rule is below some progress point, hence some non-renaming rule. Thus, each cyclic proof with renaming is equivalent to some cyclic proof without renaming: this transformation, however, may raise exponentially the size of the graph representation. 

We consider an apparent weakening of the notion of regularity: a proof is cyclic-equivalent if it is equivalent to a cyclic proof.  A cyclic equivalent proof has the global trace condition and finitely many subtrees up to proof equivalence. Conversely, if we apply any canonical renaming, from a cyclic equivalent proof we obtain a cyclic proof without renaming, but the number of subtrees raised by $n!$, with $n$ the maximum number of free variable in a sequent. 

% TOLTO 11:15 06/07/2018
%A cut-eigenvariable of $\Pi$ is an eigen-variable in a rule of $\Pi$ inferring some ancestor of some cut of $\Pi$ or included in no cut of $\Pi$. A non-cut-eigenvariable is an eigenvariable which is not a cut-eigenvariable. Cut-elimination removes or recopies cut-eigenvariables and may duplicate then rename non-cut-eigenvariable in order to preserve the fresh variable condition. As a result the number of different non-cut-eigenvariable grows, while cut-eigenvariables disappear.

Assume $A$ is any quantifier-free formula occurrence in a cyclic proof $\Pi$. If $A$ is an inductive predicate, a progress point $B$ of $A$ in $\Pi$ is any progress point of a trace of $\Pi$ starting from $A$. By definition, $B$ is some ancestor of $A$ in $\Pi$.  If $A$ is any quantifier-free formula, a minimal inductive ancestor of $A$ in $\Pi$ is any ancestor of $A$ which is an inductive predicate and is the first found in some ancestor sequence from $A$ in $\Pi$. A progress point of $A$ in $\Pi$ is any progress point of any minimal inductive ancestor of $A$ in $\Pi$. 

% TOLTO 11:15 06/07/2018
%A formula occurrence $A$ in $\Pi$ is global if all eigenvariables of $\Pi$ occurring in $A$ are cut-eigenvariables (hence not renamed by cut-elimination). In particular, if $A$ occurs in the scope of no rule of $\Pi$ with eigenvariable, then $A$ is global in $\Pi$. 

\section{Cyclic proofs: cut-rule}
In this section we formulate the problem of cut-elimination for $\infty$-proofs and for cyclic proofs.

We plan to prove that every $\infty$-proof $\Pi_1$ has a normal form $\Pi_2$ which is an $\infty$-proof. Since some cyclic proofs have no normal form, it will follow that $\Pi_1$ may be regular and its normal form $\Pi_2$ not regular. An $A$-cut is global in $\Pi$ if some positive ancestor of $A$ has a free variable which is some eigenvariable of some negative ancestor of $A$. In particular, if a cut is the conclusion of $\Pi$, by a renaming we may assume that the cut is global. We will prove that if all cuts of $\Pi_1$ are global then the normal form of $\Pi_2$ is regular. This includes the case of a unique cut in the root of $\Pi$, hence it implies that the cut rule on quantifier-free formulas is conditionally derivable within cut-free proofs. 


%
%\begin{lemma}[Induction Rule with global cuts]
%The induction rule has a proof adding only global cuts, hence is it conditionally derivable using cut-free proofs if we may remove all global cuts. 
%\end{lemma}
%
%\begin{proof}
%Assume we have cut-free proofs of $\Gamma \vdash \Delta, F[0]$ and of $\Gamma, F[x] \vdash \Delta, F[Sx]$, with $x \not \in \FV(\Gamma, \Delta)$. By renaming, we have a proof of $\Gamma, \Nat y, F[y] \vdash \Delta, F[Sy]$ for any $y \not \in \FV(\Gamma, \Delta)$. For any $x, y \not \in \FV(\Gamma, \Delta)$ with $x \not \equiv y$ we define a proof of $\Gamma, \Nat x \vdash \Delta, F[x]$ from a proof of $\Gamma, \Nat y \vdash \Delta, F[y]$ using only global cuts. We prove $\Gamma, \Nat x \vdash \Delta, F[x]$ by $\Nat L$ from $\Gamma, x=0 \vdash \Delta, F[x]$ and $\Gamma, \Nat y, x=Sy \vdash \Delta, F[x]$. We prove $\Gamma, x=0 \vdash \Delta, F[x]$ by $\eqL$ from $\Gamma \vdash \Delta, F[0]$. We prove $\Gamma, \Nat y, x=Sy \vdash \Delta, F[x]$ by $\eqL$ from $\Gamma, \Nat y \vdash \Delta, F[Sy]$. We prove this latter by cut from $\Gamma \vdash \Delta, F[y]$ and $\Gamma, \Nat y, F[y] \vdash \Delta, F[Sy]$ (which is given). Then we prove  $\Gamma \vdash \Delta, F[y]$  from  $\Gamma \vdash \Delta, F[x]$ in the same way, cyclically. We defined a regular proof, whose subproofs are those of $\Gamma, \Nat y, F[y] \vdash \Delta, F[Sy]$ and of $\Gamma, \Nat y, F[x] \vdash \Delta, F[Sx]$ and finitely many others. In any infinite path either the assumption $\Nat y$ progresses infinitely many times, or from some node on the path is in the proof  of $\Gamma \vdash \Delta, F[0]$ or of $\Gamma, F[x] \vdash \Delta, F[Sx]$, hence it includes some infinitely progressing path. There were no cuts. All cuts we added have an eigen-variable introduced by some $\Nat L$-rule inferring some cutted formula $\Nat x$ or $\Nat y$, hence global, except the outermost cut. 
%
%%%%%%%%%%%%%%%%%%%%%%%%%%
%% ATTENZIONE
%%%%%%%%%%%%%%%%%%%%%%%%%%
%%11:12 06/07/2018
%%CORREGGERE: $\Nat x$, $\Nat y$ NON SONO TAGLIATE, E INOLTRE LA DEFINIZIONE DI
%%GLOBALE E' CAMBIATA
%%%%%%%%%%%%%%%%%%%%%%%%%%
%
%The outermost cut has eigen-variable $x$ introduced by a some $\Nat L$-rule included in no cut, hence global. In both cases all eigen-variables occurring in a cut-formula are global. Thus, all cuts we have in the proof of $\Gamma, \Nat x \vdash \Delta, F[x]$ are global cuts.
%$\Box$ \end{proof}


%The cut closure of a set $\setA$ of formulas occurrences in $\Pi$ consists all subformulas of some $A \in \setA$ and all progress points of some $A\in \setA$. 
%19:50 19/06/2018



We define an equivalence relation $A \sim B$ on quantifier-free formulas, meaning that $A, B$ are the same if we neglect terms. We define $p (\vec{t}) \sim p(\vec{u})$. We define $A_1 \rightarrow A_2 \sim B_1 \rightarrow B_2$ if $B_1 \sim A_1$ and $A_2 \sim B_2$. If $c \equiv \vee, \wedge$ we define $A_1 c A_2 \sim B_1c B_2$ if $A_1 \sim B_1$ and $A_2 \sim B_2$. In all other cases $A \sim B$ is false. 

When $A \sim B$ we define $A = B$ as the list of equations between terms of $A, B$ in the same position. We define $p(t_1, \ldots, t_n) = p(u_1,\ldots, u_n)$ as $t_1 = u_1, \ldots, t_n = u_n$. We define $A_1 \rightarrow A_2 = B_1 \rightarrow B_2$ as $B_1= A_1, A_2 = B_2$. If $c \equiv \vee, \wedge$ we define $A_1 c A_2 = B_1c B_2$ as $A_1 = B_1, A_2 = B_2$. 

For instance, $p(a)=p(b)$ is $a=b$ and $p(a)\rightarrow q(b) = p(c)\rightarrow q(d)$ is $c=a, b=d$. Remark that pairs of terms occuring positively in $A, B$ are written in the order: term of $A$/term of $B$, while pairs of terms occuring negatively in $A, B$ are written in the reverse order: term of $B$/term of $A$. This choice is intended to simplify the cut-elimination proof. Alternatively, we could use the rule $\eqL$ in order to replace $a=b$ with $b=a$ whenever this is required during the proof.

If $p \not = q$, then $p(a) = q(b)$ is not defined because $p(\_), q(\_)$ are different. 


We define a generalized cut-rule $\nu$ for $A \sim B$: from $\Gamma_1 \vdash \Delta_1,A$ and $B, \Gamma_2 \vdash \Delta_2$ we deduce  $A=B,\Gamma_1, \Gamma_2 \vdash \Delta_1, \Delta_2$. We call this rule an $A,B$-cut. 
%We have three trivial cases, with assumptions ``$\Gamma_1 \vdash \Delta_1$ and $B, \Gamma_2 \vdash \Delta_2$''  or ``$\Gamma_1 \vdash \Delta_1,A$ and $\Gamma_2 \vdash \Delta_2$'' or ``$\Gamma_1 \vdash \Delta_1$ and $\Gamma_2 \vdash \Delta_2$''.

Generalized cut-rule is derivable from cut-rule and $\eqL$ and weakening. In the case $A \equiv B$ we will deduce $A=A, \Gamma_1, \Gamma_2 \vdash \Delta_1, \Delta_2$. $A = A$ is is list of identities, and if we remove them with a cut the conclusion is  $\Gamma_1, \Gamma_2 \vdash \Delta_1, \Delta_2$. Thus, cut rule $\cut$ may be deduced from $\nu$ and cut rule for identities.  
%will be a list of identities and will be removed from the proof, without introducing new cuts. 



{\bf Remark on the restricted use of substitution in the cut elimination process.} We may normalize a proof with only global cuts by normalizing all maximal cuts in it. The result is a cyclic proof if and only if all normalizations of a maximal cut produce a cyclic proof. Our cut-elimination process involving a finite set of substitution $\setS$ of all pairs $[u_i/x_j$, where: $u_1, \ldots, u_n$ are all terms occurring in some ancestor of a global cut and $x_1, \ldots, x_n$ are the eigenvariables used to infer some ancestor of a global cut. Both terms and variables are finitely many because the proof is regular. Let $\setC$ be the closure under composition of $\setC$: we prove that $\setC$ is finite. Indeed, by definition $ \FV(\vec{u})$ includes no eigenvariables of a left rule whose main formula is an ancestor of a global cut. Thus, $\vec{x} \cap \FV(\vec{u}) = \emptyset$. In this way for all $[u/x], [v/y] \in \setS$ we have $[u/x][v/y]=[v/y]$ if $x \equiv y$ and $[u/x][v/y]=[v/y][u/x]$ if $x \not \equiv y$, therefore $\setC$ consists of all compositions $[u_1/x_1]\ldots[u_k/x_k]$ of elements of $\setS$ with $x_1, \ldots, x_k$ in increasing alphabetic order, hence $\setC$ is finite. During the cut-elimination process we will have arguments of the form $\sigma(\Pi)$, for $\sigma \in \setC$ and $\Pi$ subproof of the original proof. These proofs are finitely many, and this remark will be crucial while proving regularity of the normal form.

Thus, we assume that $\Pi$ has only global cuts and the root of $\Pi$ is a cut. Assume $\setA$ is the set of cut-formulas of $\Pi$ and $\setB$ the cut closure of $\setA$. 
%We restrict $\nu$ to cuts on $A \sim B \sim C$, for $B\in \setB$. 


The set $\setS$ of $\Pi$-states is the set of finite trees $S$ whose branches are made with the rule $\nu$, and with leaves some subtrees $\Pi_1, \ldots, \Pi_k$ of $\Pi$. We rename all eigenvariables in $\Pi_i$ which are not eigenvariables of an ancestor of a cut in order to make them not occurring in any $\Pi_j$ with $j \not = i$. We require that the cut-formulas of $S$ are exactly the formulas in the leaves which are ancestors of some cut-formula in $\Pi$. 
%Thus, each cut- (non-cut-) eigenvariable in $S$ is a cut- (non-cut-) eigenvariable in $\Pi$.

A $\Pi$-state $S$ is a tree made with the rule $\nu$ only, we do not counts the nodes of the subtrees $\Pi_1, \ldots, \Pi_k$ of $\Pi$ as nodes of $S$. 

%14:37 23/06/2018


%Assume $S \in \setS$ and $\alpha \in S$ be a node decorated by $\Gamma \vdash \Delta$. We define $\cut(\alpha) \subseteq \setB$ as the set of all $B \in \setB$ such that some formula of $\Gamma \vdash \Delta$ is removed by some $B$-cut. 

%We say that a state $S \in \setS$ is irredundant if for any two nodes $\alpha \ge \beta$ of $S$ with $\cut(\alpha) = \cut(\beta)$ we have $\alpha = \beta$. $S$ is redundant otherwise. Any redundant state $S$ may be turned into an irredundant state $S_0$ by choosing some $\alpha > \beta$ of $S$ such that $\cut(\alpha) = \cut(\beta)$ and replacing $\alpha$ with $\beta$. 
%For each $B \in \setB$ and each formula $C$ of $\beta$ removed by some $B$-cut and adding $C=B$ or $B=C$ to the conclusion there is some formula $D$ of $\alpha$ removed by some $B$-cut, and conversely.  
%All other formulas of $\beta$ are never cutted and included in $\alpha$ and in the root of $S$. Thus, $S_0$ is a correct state, with conclusion some subset of the conclusion of $S$.

If we fix a bound $b$ to the height of a tree-state $S \in \setS$, then up to renaming of eigenvariable not used for an ancestor of a cut there are finitely many trees with bound $b$. 
%We choose a bound $b=2^k$, with $k$ the number of formulas in $\setB$. Any tree of height $2^k$ is redundant, because the same set of cut formulas repeats in some branch of length $2^k$. 

The $\setP$-states are the states of the algorithm removing the cut in $\Pi$ with only global cuts, for instance removing the unique cut in $\Pi = \cut(\Pi_1, \Pi_2)$ for cut-free $\Pi_1, \Pi_2$.

%19:20 20/06/2018

%13:17 14/06/2018
%11:10 07/06/2018

\section{Trivial cut-elimination: removing identities}
%15:40 03/10/2018

%%%%%%%%%%%%%%%%%%%%%%%%%%%%%%%%%%%%%%%%%%%%
% REMOVED 07:29 21/06/2018. RE-INSERTED 16:56 23/06/2018
%%%%%%%%%%%%%%%%%%%%%%%%%%%%%%%%%%%%%%%%%%%%
We define a procedure eliminating any list of assumptions $t=t$ from left-hand side of a sequent. 


Eliminating $t=t$ from the left-hand side of a sequent is equivalent to  eliminate a cut with $\vdash t = t$. 
 
%16:37 12/06/2018

\begin{lemma}[Cut elimination for $t=t$]
\label{lemma-identity}
Assume $\Pi: \Gamma \vdash \Delta$ is a cyclic proof whose only cuts are with some $t=t$. Then there is some cut-free cyclic proof $\Pi_0:\Gamma \vdash \Delta$.
\end{lemma}

\begin{proof}
We define $\Pi_0$ by removing all ancestors of some cutted $t=t$. In the case $t=t$ is inferred by $\eqL$ in: ``from $\Theta \vdash \Xi$ deduce $t=t, \Theta \vdash \Xi$'', we replace $\eqL$ with its assumption. We cannot remove $\eqL$ for $t=t$ infinitely many times consecutively, otherwise $\Pi$ would contain an infinite branch $\beta$ definitely containing only $\eqL$, hence $\beta$ would include no infinitely progressing trace. Thus, we defined a correct infinitary binary tree $\Pi_0$, with the same conclusion as $\Pi$. We check that $\Pi_0$ is cyclic. 
\begin{enumerate}
\item 
\emph{$\Pi_0$ has the global trace condition}. Any infinite branch $\gamma$ of $\Pi_0$ is obtained from some branch $\beta$ of $\Pi$ with at least the same nodes, hence infinite. $\beta$ has some infinitely progressing trace $\tau$. We never remove progress points of $\tau$ while forming $\gamma$, hence $\tau$ is in $\gamma$. 
\item
\emph{$\Pi_0$ is regular}.  Each subtree of $\Pi_0$ is obtained from some subtree $\Pi_1$ of $\Pi$ crossing out finitely many identities of the form $t=t$. For each $\Pi_1$ in $\Pi$ there are finitely many possibilities of crossing out identities. Thus, the number of subtrees of $\Pi_0$ is finite.
\end{enumerate}
$\Box$ \end{proof}

%15:55 23/06/2018

%%%%%%%%%%%%%%%  
%      13:26 11/06/2018         % 
%%%%%%%%%%%%%%%

\begin{lemma}[Equation]
\label{lemma-equation}
Assume $A, B$ are quantifier-free formulas and $A \sim B$.
There are cut-free finite proofs (hence cyclic proofs) $\Pi_{A,B}:\Gamma, A=B, A  \vdash B, \Delta$ and $\Pi'_{A,B} : \Gamma, A=B, B  \vdash A, \Delta$.
\end{lemma}

\begin{proof}
We define $\Pi_{A,B}$ from the identity rule $\Gamma, A=B,A \vdash A, \Delta$. If $A=B$ is the sublist of non-identical $t_1=u_1, \ldots, t_n = u_n$, we use $\eqL$ for $t_1=u_1$ if $t_1 \not \equiv u_1$, then for $t_2 = u_2$ if $t_2 \not \equiv u_2$, and so forth. In all cases if $t_i \not \equiv u_i$ we switch $t_i$ and $u_i$ in the first $A$ in the right-hand side. The result is $\Gamma, A=B,A \vdash B, \Delta$. 
$\Box$ \end{proof}

%We say that $\Pi$ has context $n$ if $\Pi$ have at most $n$ variables and all the variables among $\{x_0, \ldots, x_{2n-1}\}$. All regular proofs have context some $n$. 
%
%We say that $\Pi_1, \Pi_2$ have fresh eigenvariables if the eigenvariables of the two proofs are disjoint each other and from the free variables of both proofs.
%
%We define the cut-elimination procedure $\nu(\Pi_1, \Pi_2)$ for $A \sim B$, $A, B$ without quantifiers and $\Pi_1, \Pi_2$ cut-free, with disjoint eigenvariables and context $n$. If $\Pi_1, \Pi_2$ have context $n$ we may independently rename the eigenvariables of $\Pi_1$ and those of $\Pi_2$ in $\{x_0, \ldots, x_{2n-1}\}$ in order to have fresh eigenvariables and again context $n$. Indeed, assume for $i=1,2$ that $\Pi_i$ has $f_i$ free variables and $e_i$ eigenvariables and $f_i + e_i \le n$ and all the variables among $\{x_0, \ldots, x_{2n-1}\}$. From $f_i + e_i \le n$ we deduce $e_i \le n-f_i$, then $(e_1 + e_2) \le 2n - (f_1 + f_2)$ summing the two inequations, so there are enough variables in $\{x_0, \ldots, x_{2n-1}\}$ different from the free variables of $\Pi_1, \Pi_2$ for a renaming making the eigenvariables of the two proofs disjoint each other and from the free variables. The renaming produces $\le n$ variables in in $\{x_0, \ldots, x_{2n-1}\}$, therefore context $n$. Thanks to the freshness condition, during the computation of $\nu$ we do not have to rename variables. 
%
%We do not include in $\nu$ the case of inductive predicates for the moment. 

Since we have a rule for elimination of equality, whenever we have formulas $A[a,b]$ and $a=b$ in a cut-free proof we may also have $A[b,a]$.



\section{A cut-elimination algorithm for $\infty$-proofs}
In this section we define a normalization procedure for $\infty$-proofs, which may destroy regularity.
%20:26 03/10/2018

We consider an $\infty$-proof $\Pi$ having the generalized cut-free $\nu(\Pi_1, \Pi_2)$ instead of the cut-rule. We write $\Pi: \Gamma \vdash \Delta$ for: $\Pi$ has conclusion $\Gamma \vdash \Delta$ 

We say that $\Pi$ is in head-normal form when the root $r$ at the root of $\Pi$ is $\not = \nu$. 

We define a cut-elimination procedure whose goal is reducing $\Pi: \Gamma \vdash \Delta$ to some $\Pi_0: \Gamma \vdash \Delta$ in head-normal form. Then the procedure is recursively applied to the premises of $r$. For a generic infinite proof-tree the procedure may fail to terminate for some nodes. We prove that in the case $\Pi$ has the global trace condition then $\Pi: \Gamma \vdash \Delta$ always reduces to a normal form $\Pi_0: \Gamma \vdash \Delta$, and $\Pi_0$ has the global trace condition. However, $\Pi$ may be regular and $\Pi_0$ be not regular.
  
%We define the cut-elimination procedure on a generalized cut $\nu(\Pi_1, \Pi_2)$, taking $\infty$-proofs $\Pi_1, \Pi_2$ and returning a cut-free $\infty$-proof $\Pi$.  The procedure has state a pair $\nu_1 \in \Pi_1$, $\nu_2 \in \Pi_2$ of nodes and a tree of cuts between main formulas of rules with conclusion some $\mu_1 \le \nu_1$, $\mu_2 \le \nu_2$. At each step we replace either $\nu_1$ with some ancestor $\nu'_1$ or $\nu_2$ with some ancestor $\nu'_2$ or both. Either we increase the tree of cuts or we move some non-cut rule to the root. We use as main index of the procedure the maximum distance between $\nu_1$ and some ancestor which is a progress point. A progress point is in the left-hand side, the cut formula of $\nu_1$ is in the right-hand side, therefore the progress point infers some non-cut formula. If we reach a progress point then we move some non-cut rule from $\nu_1$ to the root. If we do not reach a progress point in the first premise, then after some steps $\nu_1$ stops moving up and $\nu_2$ moves. Each move of $\nu_2$ either produces some non-cut rule to the root, or we move up some $\mu_1 \le \nu_1$. Let $k$ be the height of $\nu_1$. Then after $k$ finitely many steps $\mu_1$ stops moving up, when it reaches $\nu_1$. Therefore a cut formula in $\nu_2$ may move up at most $k$ times. We cannot reach $k+1$ progress point in the same trace, which nevertheless exists by the global trace condition. Therefore after finitely many steps we reach a progress point, then we move some non-cut rule from $\nu_2$ to the root.

%20:42 03/10/2018

When we cut $\Gamma \vdash \Delta, Su$ with $t=Sx, \Nat x, \Gamma \vdash \Delta$ we perform a substitution $[u/x]$ on the eigenvariable $x$. There are finitely many terms $u$ which are witnesses of a $\Nat R$-rule in a regular proof, but in general substitutions may nest and in this case they may produce infinitely many terms. In the particular case we remove global cuts, then for any two substitutions $[u/x], [v/y]$ the terms $u,v$ include no eigenvariables, therefore we $x,y \not \in \FV(u)$ and our substitution cannot nest.
%Due to the fresh variable rule, the result is not completely cut-free in the usual sense. The fresh-variable rule hydes a cut with $\exists x.x=u$. This means that in the ``normal'' proof we may have any number of statements $x=u$ which are not part of the conclusion. 


We may use our procedure to remove an ordinary cut: from $\Gamma_1 \vdash A, \Delta_1$ and $A, \Gamma_2 \vdash \Delta_2$ deduce  $\Gamma_1, \Gamma_2 \vdash \Delta_1, \Delta_2$. Since $A \sim A $, we first deduce deduce  $A=A, \Gamma_1, \Gamma_2 \vdash \Delta_1, \Delta_2$, we remove cuts, then we cut all identities in $A=A$ and we obtain a cut-free cyclic proof of $\Gamma_1, \Gamma_2 \vdash \Delta_1, \Delta_2$.

Assume the last rule $r$ of $\Pi$ is $\not = \nu$. Then we recursively reduce all premises of $r$ (if any). If there is no assumption then the procedure termine. Thus, whenever we obtain a finite cut-free proof the procedure terminate. 

Assume $\Pi = \nu(\Pi_1, \Pi_2)$ with $\Pi_1:\Gamma_1 \vdash \Delta_1, C$ and $\Pi_2:C,\Gamma_2 \vdash \Delta_2$ and $c$ is the first connective of $C$. 
%10:00 06/10/2018

%If there is some redundant reduction on $\Pi$ (this is always the case if $\Pi$ has height $n+1$) we apply it. Otherwise no $\Pi$ has height $\le n$: in this case we define a reduction which expands $\Pi$ to some height $\le n+1$. 

Assume that for $i=1,2$ we have $\Pi_i = r_i(\{\Pi_{i,j} | j <  a_i\})$ for some $r_i$ of arity $a_i$. If $r_1=\nu$ we reduce on $\Pi_1$. If $r_1 = \sub$ we apply substitution on the unique assumption of the rule. Assume $r_1 \not = \nu, \sub$. If $r_1 = \id$ then $\Pi \reduces \Weak(\Pi_2)$.  Assume $r_1 \not = \nu, \sub, \id$.  If $r_1$ does not infer $C$ we move $r_1$ outside $\nu$: then $\Pi \reduces r_1(\{\nu(\Pi_{i,j},\Pi_2) | j <  a_i\})$. Assume $r_1$ infers $C$. If $r_1=\Weak$ then $\Pi \reduces \Weak(r_{1})$.  Assume $r_1$ infers $C$ and $r_1 \not = \nu, \sub, \id, \Weak$.  

If $r_2=\nu$ we reduce on $\Pi_2$.  If $r_2 = \sub$ we apply substitution on the unique assumption of the rule. If $r_2 = \id$ then $\Pi \reduces \Weak(\Pi_1)$.  Assume $r_2 \not = \nu, \sub, \id$.   Assume $r_2$ does not infer $C$. We move $r_2$ outside $\nu$: then $\Pi \reduces r_2(\{ \nu(\Pi_1, \Pi_{2,j}) | j \le a_i\})$. Assume $r_2$ infers $C$. If $r_2=\Weak$ then $\Pi \reduces r_{2}$.   Assume $r_2$ infers $C$ and $r_2 \not = \nu, \sub, \id, \Weak$. In this case we perform a mix reduction, according to the value of the outermost connective $c$. 


%18:37 02/10/2018

Assume $c \equiv \vee, \wedge,\rightarrow$. We consider as example the case $c \equiv \vee$. In this case $C = C_1 \vee C_2$ and $\Pi_1$ has a unique assumption $\Pi_{1,1}$ inferring $\Gamma \vdash \Delta, C_j$ for some $j$, while $\Pi_2$ has two assumption, with each $\Pi_{2,j}$ inferring $C_j, \Gamma \vdash \Delta$ for $y=1,2$. Then $\Pi \reduces \nu(\Pi'_1, \Pi'_2)$, with  $\Pi'_1 = \nu(\Pi_{1,1},\Pi_{2})$ and $\Pi'_2 = \nu(\Pi_{1},\Pi_{2,j})$. 

%If $c \not \equiv \rightarrow$ or $i \not \equiv 1$ then there is no sign change in the assumption $i$ and  $S \reduces \nu(S'_1, S'_2)$, otherwise there is a sign change and $S \reduces \nu(S'_2, S'_1)$. 

%16:21 23/06/2018
Assume $c$ is some atomic predicate. We distinguish three sets of atomic predicates: logical ($p(\vec{t})$), equality ($=$), and inductively defined predicates. 

Assume $c$ is a logical predicate: then $A = p(\vec{t})$ and $B = p(\vec{u})$. This is the only case in which we need the assumption $A=B$. 
Then $\Pi_1, \Pi_2 \reduces \id$ and $\Pi \reduces \Pi_0$, a proof of $A=B, A, \Gamma_1, \Gamma_2 \vdash \Delta_1, \Delta_2, B$.

Assume $c$ is the equality predicate. Then $\Pi_1 \reduces \eqR : \Gamma \vdash \Delta. t=t$ and $\Pi_2 \reduces \eqL(\Pi'_2)$. In this case, the rule $\eqL$ deduces $u=v,\Gamma_2[u,v] \vdash \Delta_2[u,v]$ from $\Gamma_2[v,u] \vdash \Delta_2[v,u]$. 
%We set $\Pi \reduces \eqL^4(\nu(\Pi_1, \Pi'_2))$, a degenerate $\nu$-rule, without the right-cut formula, proving $(t=t) = (u=v), \Gamma_1, \Gamma_2[v,u] \vdash \Delta_1, \Delta_2[v,u]$. We extend it to a proof of $(t=t) = (u=v), \Gamma_1, \Gamma_2[u,v] \vdash \Delta_1, \Delta_2[u,v]$. 
We set $\Pi \reduces \eqR(\eqR(\Pi_2)): t=u, t=v,\Gamma_2[u,v] \vdash \Delta_2[u,v]$. 

Assume $c$ is an inductively defined predicate: the example is $c \equiv \Nat$. 
Assume  $\Pi_1 = \Nat R(\Pi_{1,1}) :\Gamma_1 \vdash \Delta_1, \Nat Su$ and $\Pi_2 = \Nat L(\Pi_{2,1}, \Pi_{2,2}):\Nat t, \Gamma_2 \vdash \Delta_2$, with $\Pi_{1,1} : \Gamma_1 \vdash \Delta_1, \Nat u$ and $\Pi_{2,1}: t=0, \Gamma_2 \vdash \Delta_2$  and $\Pi_{2,2}: t=Sx, \Nat x, \Gamma_2 \vdash \Delta_2$. Then $\Pi \reduces \eqL(\nu(\Pi_{1,1}, [u/x]\Pi_{2,2}))$, where:
\begin{enumerate}
\item
$[u/x]\Pi_{2,2}: t=Su, \Nat u, \Gamma_2 \vdash \Delta_2$
\item
$\nu(\Pi_{1,1}, [u/x]\Pi_{2,2}) : t=Su, \Gamma_1, \Gamma_2 \vdash \Delta_1,  \Delta_2 $ 
\item
$\eqL(\nu(\Pi_{1,1}, [u/x]\Pi_{2,2})): Su=t, \Gamma_1, \Gamma_2 \vdash \Delta_1,  \Delta_2 $
is the required conclusion.
\end{enumerate}


%13:49 06/10/2018

\begin{theorem}[Cut-elimination for $\infty$-proofs] 
If $\Pi$ is an $\infty$-proof then $\Pi$ eventually reduces to some cut-free $\infty$-proof $\Pi_0$. 
\end{theorem}

\begin{proof}
Let $\Pi_0$ be the limit of the reduction process, defined as the set of nodes which eventually stop changing. 
\begin{enumerate}
\item
We prove that $\Pi_0$ is well-defined and cutfree. If all nodes of $\Pi_0$ are defined and are cut-free we are done. Assume this is not the case. If some node is undefined then some cut reduces to itself infinitely many times, and there is some node equal to a cut. Thus, some node of $\Pi_0$ is a cut. Take any cut-node $\alpha$ and the leftmost path $\pi_0$ in $\Pi_0$ from $\alpha$. Either $\pi_0$ is infinite, or $\pi_0$ stops. In the first case we set $\pi=\pi_0$. In the second case, $\pi_0$ stops in a cut $\nu(\Pi_1, \Pi_2)$ reducing to itself infinitely many times. Thus, there is some trace $\tau_1$ of the cut-formula of $\Pi_1$ passing infinitely many times through a main formula. Let $\pi_1$ be the path including $\tau$. We set $\pi = \pi_1@\pi_2$. $\pi$ is infinite, by the global trace condition we have an infinitely progressing trace $\tau$ in $\pi$. 

We claim that this cannot be. Let $p(\vec{t})$ be the inductive formula which is the element of $\tau$. There are two possibilities: $p(\vec{t})$ is an ancestor of a formula in $\alpha$ or $p(\vec{t})$ is an ancestor of a formula in some cut above $\alpha$.
\begin{enumerate}
\item
Assume that $p(\vec{t})$ is an ancestor of a formula in $\alpha$. Then either the cut in $\alpha$ would be replaced by the introduction rule of $p(\vec{t})$, or all progress points of $\tau$ would be removed while removing some cut. This is impossible because removing a cut may remove at most a finite prefix of $\tau$
\item
Assume $p(\vec{t})$ is an ancestor of a cut formula above $\alpha$. $p(\vec{t})$ is a formula in the left-hand-side, therefore $p(\vec{t})$ would be an ancestor of the cut-formula in the second premise of a cut. If the first premise reduces infinitely many times and $\pi$ should take the first premise of the cut, while $\tau$ is in the second premise, contradiction. If the first premise stabilizes to a cut, then $\pi$ should take the first premise of the cut, while $\tau$ is in the second premise, contradiction. If the first premise stabilizes to an introduction rule for the cut-formula, then the cut would remove a finite premix of $\tau$ and then stabilize, contradiction because $\tau$ is infinite. 
\end{enumerate}
Thus, $\Pi_0$ is cut-free. 
\item
We prove that $\Pi_0$ satisfies the global trace condition. Take any infinite path $\pi_0$ in $\Pi_0$. Then $\pi_0$ is the subsequence of some infinite path $\pi$ in $\Pi$. $\pi$ has some infinitely progressing trace $\tau$. Removing a cut may remove at most a finite prefix of $\tau$, therefore there is some infinitely progressing trace $\tau_1$ in $\pi_1$.
\end{enumerate}
\end{proof}

%18:41 02/10/2018
%13:32 11/09/2018

%Here is the case $(c,i) = (\rightarrow, 1)$. 
%Assume that from $\Pi_{2,1}:\Gamma_2, A_2 \rightarrow B_2 \vdash A_2, \Delta_2$ and $\Pi_{2,2}:\Gamma_2, A_2 \rightarrow B_2, B_2 \vdash \Delta_2$ we deduce $\Pi_{2}: \Gamma_2, A_2\rightarrow B_2 \vdash \Delta_2$.  
%Assume that from $\Pi_{1,1}:\Gamma_1, A_1 \vdash A_1 \rightarrow B_1, \Delta_1$ we deduce $\Pi_1: \Gamma_1 \vdash A_1 \rightarrow B_1, \Delta_1$. 
%We want to check that $\Pi' = \cut( \Pi'_2,  \Pi'_1 ): A_1 \rightarrow B_1 = A_2 \rightarrow B_2, \Gamma_1, \Gamma_2 \vdash \Delta_1, \Delta_2$. 
%
%%18:00 12/06/2018
%We cut $\Gamma_1, A_1 \rightarrow B_1 \vdash A_1, \Delta_1$ and $\Gamma_2 \vdash A_2 \rightarrow B_2, \Delta_2$ and we obtain $A_1 \rightarrow B_1 = A_2 \rightarrow B_2, \Gamma_1, \Gamma_2 \vdash A_1, \Delta_1, \Delta_2$. We cut $\Gamma_1, A_1 \rightarrow B_1 \vdash A_1, \Delta_1$ and $\Gamma_2 \vdash A_2 \rightarrow B_2, \Delta_2$ and we obtain $A_1 \rightarrow B_1 = A_2 \rightarrow B_2, \Gamma_1, \Gamma_2 \vdash A_1, \Delta_1, \Delta_2$. We cut $\Gamma_1, A_1 \rightarrow B_1, B_1 \vdash \Delta_1$ and $\Gamma_2 \vdash A_2 \rightarrow B_2, \Delta_2$ and we deduce $A_1 \rightarrow B_1 = A_2 \rightarrow B_2, B_1, \Gamma_1, \Gamma_2 \vdash \Delta_1, \Delta_2$. We cut $\Gamma_2, A_2 \vdash B_2, A_2 \rightarrow B_2, \Delta_2$ with $\Gamma_1, A_1\rightarrow B_2 \vdash \Delta_1$ and we deduce $A_1 \rightarrow B_1 = A_2 \rightarrow B_2, \Gamma_1,\Gamma_2, A_2 \vdash B_2, \Delta_1, \Delta_2$. With two more cuts we obtain $B_1=A_1, A_2 = B_2, A_1 \rightarrow B_1 = A_2 \rightarrow B_2, \Gamma_1, \Gamma_2 \vdash \Delta_1, \Delta_2$, which is $A_1 \rightarrow B_1 = A_2 \rightarrow B_2, \Gamma_1, \Gamma_2 \vdash \Delta_1, \Delta_2$.

%Let $\Pi$ be a cyclic proof with $n$ sequents. We define an order on $\Pi$-states. A state which is a leaf is smaller than a state which is not a leaf. A state $\nu(S_1, S_2)$ with $S_1, S2$ both missing the cut formula is smaller than a state with one cut-formula missing, which is smaller than a state with zero cut-formulas missing. We associate to a state $\nu(S_1, S_2)$ an index. Assume $\alpha$ is a leaf of $S_1$ (hence a subtree $\Pi_\alpha$ of $\Pi$) with distance $d_\alpha$ from the root of $S_1$. Assume $p_\alpha$ is the maximum lenght of a branch of $\Pi_\alpha$ before a progress point. We define the progress index of $S$ as the sum of all $p_\alpha * 2^{n-d_\alpha}$ for $\alpha$ leaf of $S_1$. We compare states of the form $S = \nu(S_1, S_2)$ with zero cut-formulas missing according to the progress index of $S$, if they are the same we compare first $S_1$, then $s_2$. We prove that if $S \reduces S'$, then either $S'$ starts with a non-cut rule or $S > S'$. Since states are finitely many this implies that eventually we have $S \reduces^+ T$ for some $T$ starting with a non-cut rule. This proves both that cut-elimination procedure is well-defined and that the resulting tree is regular.

%In order to prove the global trace condition for the cut-free form $\Pi_0$ of $S$, we remark that each infinite branch $\alpha$ of $\Pi_0$ is obtained cutting some rules from some interleaving $\beta$ of the branches of a subtree $T$ of $\Pi$. From $\alpha$ infinite we deduce that $\beta$ is infinite, hence $T$ is infinite. $T$ is binary, hence $T$ has a leftmost infinite branch $\gamma$, consisting of all leftmost nodes with infinitely many descendants. $\gamma$ contains some infinitely progressing trace $\tau$. 
%Let us choose the infinitely progressing trace $\tau_1$ starting from the lowest possble node of $\gamma_1$. 
%The nodes of $\tau$ which are cutted in $\alpha$ form a prefix of $\tau$. If the prefix is a finite prefix then $\alpha$ has some infinitely progressing trace. Assume the entire $\tau$ is cutted: in this case we derive a contradiction. Then $T$ contains another infinite branch $\delta$ of left-occurrences of inductive predicates which is cutted with the formulas of $\tau$ in $\gamma$. $\delta$ is infinite, is in $T$ and it is to the left of $\gamma$, contradiction.

%Our interpretation of this result is the following one. 
%Cyclic proofs (with all inductive definitions and) quantifier-free cuts interpret proofs of Arithmetic with $\Sigma^0_2$-induction. 
%We may remove quantifier-free cuts only from cyclic proofs interpreting Polynomial Arithmetic with $\Ind R$-induction, and these proofs form a proper subset of cyclic proofs quantifier-free cuts. 

We claim that in the case $\Pi$ is a cyclic proof and all cuts of $\Pi$ are global then the normal form of $\Pi$ is a cyclic proof. We prove that the tree of cuts generated by the reduction process has only finitely many possible values. Indeed, this tree has leaves which are subformulas of $\Pi$ with applied some substitution from a finite set, and this tree has an upper bound to its heigh. By regularity, there is a maximum number $k$ of steps before a progress point from any point of $\Pi$. There is a maximum number $h$ of steps before finding a trace progressing $k+1$ times. When a branch goes definitely to the left it may have at most $k$ more steps. Any branch has at most $h$ steps. Thus the tree is finite. Each subtree of the normal form is obtained from one state of the reduction process, hence the normal form has finitely many subtrees. 

\begin{thebibliography}{1}

\bibitem[1]{Anupam}
On the logical complexity of cyclic arithmetic, Anupam Das, unpublished draft. {\tt http://www.anupamdas.com/wp/log-comp-cyc-arith/  }

\bibitem[2]{Bellantoni-Cook}Bellantoni, S. and Cook, S. 
A new recursion-theoretic characterization 
of the polytime functions,
in: Computational Complexity (1992) 2: 97. 
{\tt https://doi.org/10.1007/BF01201998}

\bibitem[3]{FOSSACS2017} Berardi, Stefano and Tatsuta, Makoto. Classical system of Martin-Löf's inductive definitions is not equivalent to cyclic proof system. FOSSACS 2017,  p. 301-317.

\bibitem[4]{LICS2017} Berardi, Stefano and Tatsuta, Makoto. Equivalence of inductive definitions and cyclic proofs under arithmetic. LICS 2017,  p. 1-12.

\bibitem[5]{Brotherston2006} Brotherston, James (2006). Sequent calculus proof systems for inductive definitions, phd. thesis, Laboratory for Foundations of Computer Science, School of Informatics, University of Edinburgh.
\bibitem[6]{Buss}

Buss, Samuel R. Handbook of proof theory, volume 137. Elsevier, 1998

\bibitem[7]{Simpson2017} Simpson, Alex. Cyclic Arithmetic is equivalent to Peano Arithmetic. FOSSACS 2017, p. 283-300.

\end{thebibliography}



\section{Appendix 1. The strength of the quantifier-free fragment of cyclic proofs}
%19:43 17/06/2018
In this section we describe which formulas may be interpreted by inductive predicates and quantifier-free formulas in cyclic proofs. 

Definable inductive predicates in Martin-L\"{o}f theory inductive definitions are exactly the $\Sigma^0_1$-predicates. Indeed, it is immediate that all inductive predicates are $\Sigma^0_1$-predicates. We prove that every $\Sigma^0_1$-predicate is inductively definable. If $f$ is any primitive recursive $n$-ary map, then $f(\vec{x}) = y$ is inductively definable as some predicate of arity $n+1$, which we denote as $f(\vec{x]},y)$. Indeed, $f(\vec{x})=0$ is inductively definable by the clause $f(\vec{x]},0)$, and $f(\vec{x})=x_i$ by the clause $f(\vec{x]},x_i)$, and $f(\vec{x})=Sx_i$ by the clause $f(\vec{x},Sx_i)$. If $f(\vec{x}) = y$ and $g(\vec{x},y) = z$ are inductive then $h(\vec{x}) = g(\vec{x},f(\vec{x}))$ is definable by the clause: if $f(\vec{x},y)$ and $g(\vec{x},y,z) $ then $h(\vec{x},z)$. If $f(\vec{x}) = y$ and $g(\vec{x},y,n) = z$ are inductive then we may define $h(\vec{x},0) = f(\vec{x})$, $h(\vec{x},n+1) = g(\vec{x},h(\vec{x},n),n)$. $h$ is definable by the clauses: if$f(\vec{x},y)$ then $h(\vec{x},0,y)$ and if $h(\vec{x},n,y)$ and $g(\vec{x},y,n,z)$ then $h(\vec{x},n+1,z)$. Any $\Sigma^0_1$-predicate $p(\vec{x})$ is definable by $\exists x.f(\vec{x},x)=0$, therefore by the clause: if $f(\vec{x},x,0)$ then $p(\vec{x})$.

All totality statements $t(\vec{x})$ for primitive recursive maps $f(\vec{x})$ are expressible by the clause: if $f(\vec{x},y)$ then $t(\vec{x})$. For any primitive recursive map $f(\vec{x})$ we may prove $\Nat(x_1), \ldots, \Nat(x_n) \vdash t(\vec{x})$ by induction. If we assume that successor is injective and never zero, we may prove all axioms of Robinson arithmetic. Thus, the quantifier-free fragment of Martin-L\"{o}f contains $I\hy \Sigma^0_1$-arithmetic.

%%%%%%%%%%%%%%%%%%%%%%%%%%%
% REMARKS ABOUT  SIZE INCREMENT DUE TO NORMALIZATION ARE NOT TRUE
% 17:33 14/07/2018
%%%%%%%%%%%%%%%%%%%%%%%%%%%
%The representation of $\Sigma^0_1$-predicates by inductive predicates allows much shorter cut-free proofs. Cut-free intuitionistic arithmetical proofs of existential statements have the witness property: if $\vdash \exists x.p(x)$ then the proof ends with a proof of $p(t)$ for some term $t$ of the language. Nothing similar holds for cyclic proofs and for inductive predicates: here is an example. Let $\exp(x,y)$ be the inductive predicate for $2^{2^x} = y$ and $t(x)$ be the totality predicate representing $\exists y. \exp(x,y)$. Then $t(10)$, hence there is some $y$ with value $2^{2^10} \equiv 2^{1024}$. If we represent $t(x)$ with an existential quantifier, then any intuitionitistic normal proof of $t(1024)$ would include the proof that $2^{1024}$ is an integer, obtained from $\Nat 0$ by applying $2^{1024}$ times the successor rule: if $\Nat x$ then $\Nat Sx$. This proof would have size $2^{1024}$, which is more than $10^{300}$, hence would not fit in the known universe. Instead a cut-free cyclic proof of $t(10)$, with $t$ as inductive predicate, fits in few pages. The reason is that removing cuts from a cyclic proof does not expand a proof of $p(n)$ into $n$ cuts of $p(x)$ with $p(x) \rightarrow p(x+1)$, one for each $x<n$, as it happens when we reduce the induction rule. Instead, the graph representation of a cut-free proof of $t(x)$ and of $t(n)$ differ at most by a factor $2$, no matter what is $n$.
%
%When we remove cuts from a proof in $I\hy \Sigma^0_1$ we expand the induction rule and therefore we expand the size by the ackermann function (at least). Instead, we expect that if remove cuts from an arithmetical or cyclic quantifier-free proof then we expand the size in an elementary way (some finite and fixed number of exponentials),  due to the use of back-links for reducing the size of the proof-graph. Rather, we expect the proof complexity be hidden in proving the global trace condition for the normal form. Possibly this cut-elimination proof requires K\"{o}nig's Lemma for $\Sigma^0_1$-trees and the I-$\Sigma^0_2$ fragment of arithmetic. 
%%%%%%%%%%%%%%%%%%%%%%%%%%%%%
%%%%%%%%%%%%%%%%%%%%%%%%%%%%%

%12:34 29/06/2018
\section{Appendix 2. Robinson's arithmetic in cyclic proofs}
In this section we explain how to interpret Robinson's axioms thorugh inductive definitions. As a consequence, connective-free cyclic proofs are far more expressive than we expect, they prove the same theorems without alternating quantifiers as Arithmetic with $\Sigma^0_2$-induction. In particular, they prove the consistency of arithmetic with $\Sigma^0_2$-induction, and this theorem has no cut-free proof.  

%We may provide an  example of quantifier-free (in fact, connective-free) cyclic proof without cut-free form and simpler than induction. Connective-free cyclic proofs in the languague $0,S$ prove the totality of the Ackermann–Peter function, and no cut-free proof of the same result exists (in the full system). A similar result holds for the doubling map $f(x) = 2x$. 

We consider as inductive definitions exactly all \emph{``productions'' defined in Brotherston thesis}, Def. 2.2.1, p. 27. We define an inductive predicate $d(x,y)$ for ``$x,y$ are different integers'', as follows: $d(0,Sy)$ and $d(Sx,0)$ and if $d(x,y)$ then $d(Sx,Sy)$. Using $d(x,y)$ we may define the \emph{negation} of the $0,S$-axioms and of the axiom for $+,*$. We define: if $0=Sx$ then $Z$, and if $Sx=Sy$ and $d(x,y)$ then $S$. Then $\neg Z$ is equivalent to $\forall x.0\not=Sx$ and $\neg S$ is equivalent to $\forall x. \neg(Sx =Sy \wedge d(x,y))$. We claim that this is a way of stating injectivity of successor. 

By $\Nat L$ we may prove $\vdash t=u, d(t,u)$ from $\vdash d(0,Sy)$ and $\vdash d(Sx,0)$ and $\vdash 0=0$ and $\vdash x=y, d(x,y)$ for any $x, y \not \in \FV(t,u)$. The first three sequents are axioms and in the forth one both $\Nat t$ and $\Nat u$ progress. This is a cyclic proof. 

Using rule $d L$ twice and weakening, we prove that $x=y, d(x,y) \vdash Z,S$ is a consequence of $Sx=0 \vdash Z$ and $0=Sy \vdash Z$ and $Sx=Sy, d(x,y) \vdash S$. All these sequents are axioms.

Besides, we may prove $\vdash 0 \not = Sx, Z$ and $Sx = Sy \rightarrow x=y, S$. Thus, we derive $\vdash (x \not y) \leftrightarrow d(x,y), Z,S$. 

If we interpret $\Gamma \vdash \Delta$ as $\Gamma \vdash \Delta,Z,S$ then zero and successor axioms are derivable. We may complete the interpretation of Robinson's arithmetic of the base and the inductive case of the definition of $+$ and $*$. We express their negations. The original axioms are: $x+0=x$ and $x+Sy=S(x+y)$ and $x*=x$ and $x*Sy=(x*y)+y$. Their negations are expressed as follow: if $d(x+0,x)$ then $Z+$, if $d(x+Sy, S(x+y) )$ then $S+$, if $d(x*0,0)$ then $Z*$, if $d(x*Sy, x*y+x)$ then $S*$.  Let $\Xi = Z,S,Z+,S+,Z*,S*$: this is the list of negations of $6$ Robinson's axiom. The only Robinson's axioms missing is expressed by the rule $\Nat L$.  Thus we interpret $\Gamma \vdash \Delta$ as $\Gamma \vdash \Delta,\Xi$: in this way all Robinson's axioms are derivable.

We conjecture:
\begin{enumerate}
\item
the quantifier-free fragment of $\LKID$ with Robinson's axioms is equivalent to $I \hy \Sigma^0_2$-arithmetic. 
\item
the cut-free and quantifier-free fragment of $\CLKID^\omega$ with Robinson's axioms is equivalent to polynomial arithmetic.
\end{enumerate}
The reason is the same in both cases: we may express $\Sigma^0_1$-formulas of arithmetic (with the inductive predicate $\Nat$ and $0,S,+,*$) through quantifier-free formulas of $\LKID$ (with all inductive definitions). 


%
%
%%21:26 05/06/2018
%\section{A cut-elimination procedure for quantifier-free cuts}.
%
%We define a cut-elimination procedure for $A=B$-cuts. We consider cyclic-proofs without substitution rule, with the fresh variable rule. 
%
%Due to the fresh variable rule, the result is not completely cut-free in the usual sense. The fresh-variable rule hydes a cut with $\exists x.x=u$. This means that in the ``normal'' proof we may have any number of statements $x=u$ which are not part of the conclusion. Since we have a rule for elimination of equality, whenever we have formulas $A[a,b]$ and $a=b$ in the proof we may also have $A[b,a]$.
%
%
%
%%We will define a cut-elimination between to proofs $\Pi_1, \Pi_2$ having two finite pairwise disjoint sets of eigenvariables. The eigenvariables will not change through the cut-elimination.
%
%We will prove that there is a maximum number of variables in a sequent during reduction, therefore only finitely many indexes of fresh variables are required. 
%
%%We start from a cut between $\Pi_1$ and $\Pi_2$, with bound variables in $\Pi_1$ disjoint from free or bound variables from $\Pi_2$ and conversely. We always assume that the conclusion of the cut is the union of one sequent $\Theta_1 \vdash \Xi_1$ from $\Pi_1$ and one sequent $\Theta_2 \vdash \Xi_2$  from $\Pi_2$, and if the rule $r$ binding the fresh variable $x$ is from $\Pi_1$ then it has conclusion $\Theta_1 \vdash \Xi_1$, if it is from $\Pi_2$ then it has conclusion $\Theta_2 \vdash \Xi_2$. If $r$ is from $\Pi_1$ then $x$ is fresh w.r.t.  $\Theta_1 \vdash \Xi_1$ by definition of $r$, and is fresh w.r.t. $\Theta_2 \vdash \Xi_2$ by the initial choice. If $r$ is from $\Pi_2$ we reverse the roles of $\Pi_1$ and $\Pi_2$ and we conclude again that $x$ is fresh.
%
%We use the fresh variable rule while removing a cut of $\pi_1:\Gamma_1 \vdash \Delta_1, \Nat Su$ and $\pi_2:\Nat t, \Gamma_2 \vdash \Delta_2$, when $\pi_1$ introduces $\Nat Su$ from $\Gamma_1 \vdash \Delta_1, \Nat u$ and $\pi_2$ introduces $t$ from $\pi_2: t=0, \Gamma_2 \vdash \Delta_2$  and $t=Sx, \Nat x, \Gamma_2 \vdash \Delta_2$. The cut reduces to a cut of $\Gamma_1 \vdash \Delta_1, \Nat u$ with $t=Sx, \Nat x, \Gamma_2 \vdash \Delta_2$. By cut we obtain $x = u, t = Sx, \Gamma_1, \Gamma_2 \vdash \Delta_1,  \Delta_2 $, then by $\eqL$ we get $u=x, t= Su, \Gamma_1, \Gamma_2 \vdash \Delta_1,  \Delta_2 $. Eventually by the fresh variable rule we conclude $t= Su, \Gamma_1, \Gamma_2 \vdash \Delta_1,  \Delta_2$.
%
%Let $\Pi$ be any infinitary binary proof. A \emph{reduction state} for $\Pi$, a state for short, is any finite application of cuts starting from subtrees of $\Pi$. An height for a state is inductively defined as follows: $0$ is an height for a subtrees of $\Pi$, if $\Pi_1$, $\Pi_2$ have height $h_1, h_2$ then any cut of $\Pi_1$, $\Pi_2$ has height $h = \max(h_1, h_2)$. A state has bound $b$ if it has some height $\le b$. For instance, if $\cut(\Pi_1, \Pi_2)$ is a subtree of $\Pi$, then it is a state, it has height $0$, it has height $1$ (being the cut of two subtrees), and it has bound $0$.
%
%Assume $\Pi$ is regular, with $k$ subtrees.
%\begin{enumerate}
%\item
%For any $n$, $\Pi$ has are finitely many states with bound $n$. For $0$ this is true by assumption. 
%Assume there at most $a$ states with bound $n$. They have a maximum number of formulas $b$ for the left- and right-hand side of their conclusion. 
%Then the number of cuts is at most $a*b*b*a$: for a cut we have to choose a first premise $\Pi_1$ among $a$, one of the at most $b$ formulas in the right-hand-side of the conclusion of $\Pi_1$, a second premise $\Pi_2$ among $a$, one of the at most $b$ formulas in the left-hand-side of the conclusion of $\Pi_2$. 
%Not all choices produce a correct cut, but there are are at most $a*b*b*a$ choices, hence at most $k + a^2b^2$ states with bound $n+1$, all subtrees union all cuts of two states with bound $n$. 
%
%%%%%%%%%%%%%%%%%%%%%%
%% 14:31 11/06/2018
%%%%%%%%%%%%%%%%%%%%%%
%% CORREGGERE: LE h FORMULE TAGLIATE 
%% VENGONO DIVISE TRA LE DUE PREMESSE DEL
%% TAGLIO, PER CUI SONO MENO DI h
%%%%%%%%%%%%%%%%%%%%%%
%
%\item
%Assume $m$ is the maximum number of formulas in a sequent of $\Pi$: then all states of $\Pi$ have bound $m$. Indeed, let $\Pi_0$ be a state of height $h$: we prove that $h \le m$. If we move upward from the conclusion of $\Pi_0$, the children have one formula in the sequent, which is cut, the grand-children have a second one, different form the first, and so forth. If the height is $h$ then the leaves of the state have conclusion a sequent with $h$ pairwise distinct formulas which are cut, hence al least $h$ formulas. Thus, $h \le m$. 
%
%\item
%$\Pi$ has finitely many states. By point $2$ there is a common bound $m$ for all states. By point $1$ the states of bound $m$ are finitely many. 
%\end{enumerate}
%
%A leaf of a state is paused if is the left premise of a cut and introduces the cut formula (in the right-hand-side). In all subtrees, the rightmost leaf is never paused. The active leaf of a state is the leftmost non-paused leaf: there is some non-paused leaf, hence there is the leftmost non-paused leaf.
%
%A $k$-labeled state is a state with some $h \le k$ associate to each cut: we associate no integer to the subtrees of $\Pi$, one to each cut we add. If the states are finitely many then the $k$-labeled states are finitely many. We order the labeled states. The subtrees of $\Pi$ are minimal, we compare two cuts by the label, if the label is equal we compare the first, then the second premise of the cut. By induction on the definition of states we prove that this is an order relation.
%
%We define a reduction procedure for an infinitary $\Pi$ generating some infinitary $\Pi_0$. The initial state of the procedure is $\Pi$, and each step the procedure either move to a smaller state or generate the conclusion of $\Pi_0$ and produce one or two premises for $\Pi_0$ of out of one or two states of $\Pi$. Thus, each subtree of $\Pi_0$ is generated from some state of $\Pi$. Hence if $\Pi$ is regular then $\Pi_0$ is. 
%
%
%
%%08:42 11/06/2018
%We will prove that the global trace condition, if any, is preserved. Indeed, each reduction step preserves the first $a$ leaves of the state and replaces the leaf number $a$ with an immediate subtree. If we continue in this way some leaf is changed enough times to provide a progress point. This progress point may be used to reduce the index of some cut, but if we continue enough times, we obtain a progress point which is not used to reduce the index of some cut, but it is moved to the root.
%
%Let $A$ be an occurrence of formula in the conclusion of $\Pi$. The $A$-index of $\Pi$ is the number of inferences of $A$ we have in the branches of $\Pi$ before the first progress point. Every branch of $\Pi$ has some progress point,  by K\"{o}nig's lemma if we prune the branches of $\Pi$ to the first progress point we obtain a finite tree, hence the $A$-index is finite.
%
%Assume $\Pi$ is regular. Let $p$ be the maximum length of a path from some subtree of $\Pi$ without a progress. For each subtree there is a maximum length, otherwise by K\"{o}nig's lemma there would be an infinite path without progress. If $\Pi$ is regular there is a maximum length common to all subtrees, hence $p$ is finite. Any formula $A$ has $A$-index $\le p$ in any subtree of $\Pi$. Assume $\Pi$ is cyclic (it has the global trace condition and is regular) and the cuts in $\Pi$ are all quantifier-free and $q$ is the maximum degree of a cut in $\Pi$. Let $k=p+q$. We associated $k$-indexes to cuts of states of $\Pi$. Assume we have a cut on $A$ and $A$ has degree $h \le q$. If $A$ is not an inductive predicate the index of the cut is $p+h$. 
%If $A$ is an inductive predicate the index of the cut is defined as the $A$-index in the active of first premise of the cut. By assumption it is $\le p$. If $A$ is a proper subformula of $B$, then a cut on $A$ has index smaller than the cut on $B$.
%
%
%\begin{theorem}[Normalization Theorem]. 
%If a proof-tree has finitely many cuts (say, one cut at the root) and the global trace condition then it has a normal form with the global trace condition.
%\end{theorem}
%
%\begin{proof}[Sketch]
%We define a reduction process, such that if $\cut(\Pi_1,\Pi_2)$ reduces to $\Pi_3$ then any branch of $\Pi_3$ is the interleaving of a branch of $\Pi_1$ and of a branch of $\Pi_2$, after removing all inferences of some ancestor of a removed cut-formula. As a consequence, any reduction step preserves the global trace condition. If we have the global trace condition, for each atomic formula $A$ on the right-hand-side, we cannot have an infinite sequence $\ldots, A_2, A_1, A_0 \equiv A$ such that $A_{i+1}$ is the side formula of an inference with conclusion $A_i$. Indeed, $A_0$ is coming from some $\Pi_i$, and these rules in the right-hand side, together with any rules in the right-hand side we removed, form some infinite branch of $\Pi_i$ with no progressing point : progressing points are for rules in the left-hand side. We should have some infinite progressing trace in any infinite branch of $\Pi_i$, contradiction. We may to prove that in finite time a cut on $A, B$ at the root is replaced by cut on some hereditary premises $A', B'$ of $A, B$, or by a rule different from cut. The first case cannot go on forever, because it first transforms $A, B$ into a cut $A',B'$ between atomic formulas, then it cannot continue forever on the formula in the right-hand side. Thus, after a finite time the root stabilizes to some cut, some child of the root to some cut and so forth. We would have an infinite branch made of cuts, corresponding to some infinite branch in $\Pi_1$ and $\Pi_2$. After some time the branch has two progress point. The first could become the assumption right formula of a cut $A_i, B_i$, but the second one would make thecut $A_i, B_i$ progress, contradiction. We conclude that  after finite time the cut at the root is replaced by a rule different from the cut. By repeating this operation on the children of the root, on the grand children and so forth, we push cuts into the infinity. 
%$\Box$ \end{proof}

%%%%%%%%%%%%%%%%%%%%%%%%%%%%%
%%22:36 07/06/2018 - parti precedenti rimosse
%%%%%%%%%%%%%%%%%%%%%%%%%%%%%

%%%%%%%%%%%%%%%%%%%%%%%%%%%%%%
%\section{Appendix. A test: proving cut-elimination for equations in sequent calculus}
%Cut-elimination for first order logic usually does not consider introduction and elimination rules for equality, while cyclic proofs include introduction and elimination for equality. As a test, we prove cut-elimination for a sequent calculus extended with introduction and elimination for equality. This result is much simpler that cut-elimination for equality and cyclic proof, but we proved it as a test, before trying to prove the stronger result. 
%
%We assume having a first-order language, for sake of simplicity zero and successor: $0,S$. Equality rules for sequent calculus are $(\eqR)$: $\Gamma \vdash \Delta,t = t$ and $(\eqL)$: if $\Gamma[a,b] \vdash \Delta[a,b]$ then $a=b,\Gamma[b,a] \vdash \Delta[b,a]$. We want to prove cut-elimination for cuts on $a=b$ and for sequent calculus.
%
%\begin{lemma}
%\begin{enumerate}
%\item
%Assume $\pi_1:\Gamma_1 \vdash \Delta_1, a=b$ and $\pi_2:\Gamma_2[a,b] \vdash \Delta_2[a,b]$ are cut-free. Then there is some cut-free $\pi_3:  \Gamma_1, \Gamma_2[b,a] \vdash \Delta_1, \Delta_2[b,a]$. 
%\item
%Assume $\pi_1:\Gamma_1 \vdash \Delta_1, a=b$ and $\pi_2:a=b, \Gamma_2 \vdash \Delta_2$ are cut-free. Then there is some cut-free $\pi_3: \Gamma_1, \Gamma_2 \vdash \Delta_1, \Delta_2$.
%\end{enumerate}
%\end{lemma}
%
%\begin{proof}
%\begin{enumerate}
%
%\item
%By induction on $\pi_1$. Let $r$ be the last rule of $\pi_1$. We argue by cases on $r$. 
%\begin{enumerate}
%\item
%Assume $r \not \equiv (\eqL)$. Then the assumption number $i$ of $\pi_1$ is $\pi_{1,i} : \Gamma_{1,i} \vdash \Delta_{1,i}, a=b$, with the same $a,b$ (we add $a=b$ everywhere in $\pi_{1,i}$  if it is not there). By induction hypothesis on $\pi_{1,i}$ there is some cut-free $\pi_{3,i}:  \Gamma_{1,i}, \Gamma_2[b,a] \vdash \Delta_{1,i}, \Delta_2[b,a]$. We apply $r$ and we define some cut-free $\pi_3:  \Gamma_1, \Gamma_2[b,a] \vdash \Delta_1, \Delta_2[b,a]$. 
%\item
%Assume $r \equiv (\eqL)$. Then for some $t=u \in \Gamma_1$ we have $\Gamma_1 \vdash \Delta_1, a= b 
%\equiv t=u,\Gamma'_1[u,t] \vdash \Delta'_1[u,t], a'[u,t]=b'[u,t]$, in particular $a \equiv a'[u,t]$ and $b \equiv b'[u,t]$). The premise of $r$ is $\pi_{1,1}: \Gamma'_1[t,u] \vdash \Delta'_1[t,u], a'[t,u]=b'[t,u]$. 
%By assumption, $\pi_2:\Gamma_2[a'[u,t],b'[u,t]] \vdash \Delta_2[a'[u,t],b'[u,t]]$. By $(\eqL)$ we obtain a cut-free proof of $t=u,\Gamma_2[a'[t,u],b'[t,u]] \vdash \Delta_2[a'[t,u],b'[t,u]]$.
%We apply the induction hypothesis on $\pi_{1,1}: \Gamma'_1[t,u] \vdash \Delta'_1[t,u], a'[t,u]=b'[t,u]$ and we obtain a cut-free proof of $t=u,\Gamma'_1[t,u],\Gamma_2[b'[t,u], a'[t,u]] \vdash \Delta'_1[t,u], \Delta_2[b'[t,u], a'[t,u]]$. By $(\eqL)$ we obtain a cut-free proof of $t=u,\Gamma'_1[u,t],\Gamma_2[b'[u,t], a'[u,t]] \vdash \Delta'_1[u,t], \Delta_2[b'[u,t], a'[u,t]]$. This latter is $\Gamma_1,\Gamma_2[b,a] \vdash \Delta_1, \Delta_2[b,a]$, as wished.  
%\end{enumerate}
%
%\item
%By induction on $\pi_2$. Let $r$ be the last rule of $\pi_2$. We argue by cases $r$. 
%\begin{enumerate}
%
%\item
%Assume $r \not \equiv (\eqL)$. Then the assumption $\pi_{2,i}$ having number $i$ in $\pi_2$ is $a=b,\Gamma_{2,i} \vdash \Delta_{2,i}$, with the same $a,b$  (we add $a=b$ everywhere in $\pi_{2,i}$  if it is not there). By induction hypothesis on $\pi_{2,i}$ there is some cut-free $\pi_{3,i}: \Gamma_{1}, \Gamma_{2,i} \vdash \Delta_{1}, \Gamma_{2,i}$. We apply $r$ and we define some cut-free $\pi_3:  \Gamma_1, \Gamma_2 \vdash \Delta_1, \Delta_2$. 
%
%\item
%Assume $r = \equiv (\eqL)$, and $r$ has assumption $a=b$.  Then $a=b, \Gamma_2, \Delta_2 \equiv a=b,\Gamma'_2[b,a] \vdash \Delta'_2[b,a]$ and $\pi_2$ has premise $\pi_{2,1} :  (a=b),\Gamma'_2[a,b] \vdash \Delta'_2[a,b]$ (we add $(a=b)$ everywhere in $\pi_{2,1}$ if $(a=b)$ is not already there). By induction hypothesis on $\pi_{2,1}$ we have a cut-free proof of $\Gamma_1,\Gamma'_2[a,b] \vdash \Delta_1,\Delta'_2[a,b]$.  By point $1$ above we obtain a cut-free proof of $\Gamma_1,\Gamma'_2[b,a] \vdash \Delta_1,\Delta'_2[b,a]$, which is $\Gamma1,\Gamma_2 \vdash \Delta_1, \Delta_2$. 
%
%\item
%Assume $r = \equiv (\eqL)$, and $r$ has assumption $t=u \in \Gamma_2$ different from $(a=b)$.  Then $a=b,\Gamma_2, \Delta_2 \equiv a'[u,t]=b'[u,t],t=u,\Gamma'_2[u,t] \vdash \Delta'_2[u,t]$ and $\pi_2$ has premise $\pi_{2,1} : a'[t,u]=b'[t,u], \Gamma'_2[t,u] \vdash \Delta'_2[t,u]$. There is a cut-free proof of $t=u,\Gamma1 \vdash (a'[t,u]=b'[t,u]),\Delta_1$. By induction hypothesis on $\pi_{2,1}$ we have a cut-free proof of $ \Gamma_1, \Gamma'_2[t,u] \vdash \Delta_1, \Delta'_2[t,u]$, then of $t=u, \Gamma_1, \Gamma'_2[u,t] \vdash \Delta_1,\Delta'_2[u,t]$, which is $\Gamma_1, \Gamma_2 \vdash \Delta_1, \Delta_2$. 
%\end{enumerate}
%\end{enumerate}
%
%$\Box$ \end{proof}
%
%%12:00 05/06/2018
%
%\section{Appendix 2: normalization without regularity}
%We may define a normalization algorithm on infinitary binary substitution-free proof trees. The cut we may remove is: from $E, A \vdash B$ and $E, B \vdash A$ only using $\id, \eqL$, and $\Gamma_1 \vdash \Delta_1, A$ and $B, \Gamma_2 \vdash \Delta_2$ then $E, \Gamma_1, \Gamma_2 \vdash \Delta_1, \Delta_2$. The algorithm uses substitution on bound proof variables. The algorithm may diverge in some case, but if the tree has the global trace condition then it converges to some proof-tree with global trace condition. However, even if the original tree was regular the normal tree often is not regular, because as we said cyclic proofs have no cut-elimination in general. 
%
%\begin{theorem}[Normalization Theorem]. 
%If a proof-tree has finitely many cuts (say, one cut at the root) and the global trace condition then it has a normal form with the global trace condition.
%\end{theorem}
%
%\begin{proof}[Sketch]
%We define a reduction process, such that if $\cut(\Pi_1,\Pi_2)$ reduces to $\Pi_3$ then any branch of $\Pi_3$ is the interleaving of a branch of $\Pi_1$ and of a branch of $\Pi_2$, after removing all inferences of some ancestor of a removed cut-formula. As a consequence, any reduction step preserves the global trace condition. If we have the global trace condition, for each atomic formula $A$ on the right-hand-side, we cannot have an infinite sequence $\ldots, A_2, A_1, A_0 \equiv A$ such that $A_{i+1}$ is the side formula of an inference with conclusion $A_i$. Indeed, $A_0$ is coming from some $\Pi_i$, and these rules in the right-hand side, together with any rules in the right-hand side we removed, form some infinite branch of $\Pi_i$ with no progressing point : progressing points are for rules in the left-hand side. We should have some infinite progressing trace in any infinite branch of $\Pi_i$, contradiction. We may to prove that in finite time a cut on $A, B$ at the root is replaced by cut on some hereditary premises $A', B'$ of $A, B$, or by a rule different from cut. The first case cannot go on forever, because it first transforms $A, B$ into a cut $A',B'$ between atomic formulas, then it cannot continue forever on the formula in the right-hand side. Thus, after finite time the cut at the root is replaced by a rule different from the cut. By repeating this operation on the children of the root, on the grand children and so forth, we push cuts into the infinity. 
%$\Box$ \end{proof}
%%%%%%%%%%%%%%%%%%%%%%%%%%%%%

%%%%%%%%%%%%%%%%%%%%%%%%%%%%%%%%%%
%NORMALIZATION WITH RENAMING INCLUDED - removed 07/06/2018
%%%%%%%%%%%%%%%%%%%%%%%%%%%%%%%%%%
%A cut with substitution or an s-cut consists in proofs $\pi_1:\Gamma_1 \vdash \Delta_1, A_1$ and  $\pi_2:A_2, \Gamma_2 \vdash \Delta_2$ and substitutions $\theta_1, \theta_2$ such that $\theta_1 \Gamma_1, \theta_2 \Gamma_2 \vdash \theta_1(A_1) = \theta_2(A_2)$. The conclusion is $\theta_1 \Gamma_1, \theta_2 \Gamma_2 \vdash \theta_1 \Delta_1, \theta_2 \Delta_2 $. Equality of formulas is short for having equal terms in equal positions. We add a fresh variable rule: if $x=t,\Gamma \vdash \Delta$ and $x \not \in \FV(t,\Gamma,\Delta)$ then $\Gamma \vdash \Delta$. The rule is implicitly a cut with a proof of $\vdash \exists x \in \Nat.x=t$. We use the rule while removing an s-cut of $\pi_1:\Gamma_1 \vdash \Delta_1, \Nat Su$ and $\pi_2:\Nat t, \Gamma_2 \vdash \Delta_2$, when $\pi_1$ introduces $\Nat Su$ from $\Gamma_1 \vdash \Delta_1, \Nat u$ and $\pi_2$ introduces $t$ from $\pi_2: t=0, \Gamma_2 \vdash \Delta_2$  and $t=Sx, \Nat x, \Gamma_2 \vdash \Delta_2$. The cut reduces to a cut of $x=u, \Gamma_1 \vdash \Delta_1, \Nat u$ with $t=Sx, \Nat x, \Gamma_2 \vdash \Delta_2$. This is an s-cut because $x=u, \Gamma_1, \Gamma_2 \vdash x=u$ and proves $x=\theta_1 u, theta_2 t=Sx, \theta_1 \Gamma_1, \theta_2 \Gamma_2 \vdash \theta_1 \Delta_1, \theta_2 \Delta_2 $. From this sequent we deduce $x=\theta_1 u, theta_2 t=S \theta_1u, \theta_1 \Gamma_1, \theta_2 \Gamma_2 \vdash \theta_1 \Delta_1, \theta_2 \Delta_2 $. By a cut with $\theta_1 \Gamma_1, \theta_2 \Gamma_2 \vdash \theta_2(t) = S \theta_1(u)$ we obtain $x=\theta_1 u, \theta_1 \Gamma_1, \theta_2 \Gamma_2 \vdash \theta_1 \Delta_1, \theta_2 \Delta_2 $. Eventually by the fresh variable rule we obtain $x=\theta_1 u, \theta_1 \Gamma_1, \theta_2 \Gamma_2 \vdash \theta_1 \Delta_1, \theta_2 \Delta_2 $.



\section{Appendix 3. Proofs with cyclic context}
Assume $L$ is a language without function symbols: $L$ consists of logical predicates, constants and inductive predicates defined out of them. We are looking for a set of $\infty$-proofs including cyclic proofs for the language $L$. This set should have cut-elimination and should be described by finite automata. Cyclic proofs for $L$ have no cut-elimination: we checked that there are algorithm requiring non-constant space  which may be proved terminating by a cyclic proof with cuts, and no cyclic proof without cuts.

Cyclic proofs have finitely many context. We define \emph{proofs with cyclic context}, in which contexts are inductively defined using two operations: concatenation of contexts and creation of fresh variables. We inductively define at the same time the lists of pairwise distinct variables of the context. We call these proofs: proofs with cyclic context. They are defined by finite graphs. A graph node may be the conclusion of one rule for each possible shape of the context decorating it. The unfolding of the graph may give a pseudu-proof with infinite contexts, but correct proofs 

As example, we consider the canonical cut-free $\infty$-proof of $\Backward(x,y) \vdash \Forward(x,y)$. We prove  $\Backward(x,y) \vdash \Forward(x,y)$ by $\Backward L$ with first premise $R(x,y) \vdash \Forward(x,y)$ and second premise $\Backward(x,a), R(a,y) \vdash \Forward(x,y)$ ($a$ fresh). The first premise is proved from $R(x,y) \vdash R(x,y)$, which is proved by $\id$. We prove the second premise $\Backward(x,a), R(a,y) \vdash \Forward(x,y)$ by $\Backward L$ with first premise $R(x,a), R(a,y) \vdash \Forward(x,y)$ and second premise  $\Backward(x,b), R(b,a), R(a,y) \vdash \Forward(x,y)$ ($b$ fresh). The first premise is proved from $R(x,a), R(a,y) \vdash R(x,a)$ (which is an $\id$) and $R(x,a), R(a,y) \vdash \Forward(a,y)$, which is proved from $R(x,a), R(a,y) \vdash R(a,y)$ (which is an $\id$). We prove a general version of the second premise: $\Backward(x,b), R(b,a), \Gamma[a,\vec{a},y] \vdash \Forward(x,y)$. 
Here $\Gamma[a,\vec{a},y]$ denotes a context we will inductively define. The base case is $\Gamma[a,\vec{a},y] = R(a,y)$ and $\vec{a}=\nil$. The base case covers $\Backward(x,b), R(b,a), R(a,y) \vdash \Forward(x,y)$. 
%21:32 25/09/2018
We prove $\Backward(x,b), R(b,a), \Gamma[a,\vec{a},y]\vdash \Forward(x,y)$ by $\Backward L$ with the following premises.

\begin{enumerate}
\item
\emph{First premise:}
$R(x,b), R(b,a), \Gamma[a,\vec{a},y] \vdash \Forward(x,y)$. We prove it by $\Forward R$ with first premise $R(x,b), R(b,a), \Gamma[a,\vec{a},y] \vdash R(x,b)$ (which is an $\id$-rule) and $R(x,b), R(b,a), \Gamma[a,\vec{a},y] \vdash \Forward(b,y)$, which we prove by $\Weak$ and $R(b,a), \Gamma[a,\vec{a},y] \vdash \Forward(b,y)$. Here we loop. If $R(b,a), \Gamma[a,\vec{a},y] = R(b,a), R(a,y)$, then we already considered the case. If $R(b,a), \Gamma[a,\vec{a},y] = R(b,a), R(a,c), \Gamma[\vec{a},y]$, then $R(b,a), R(a,c), \Gamma[\vec{a},y] \vdash \Forward(b,y)$ is a renaming by $[b/x,a/b,c/a]$ of $R(x,b), R(b,a), \Gamma[\vec{a},y] \vdash \Forward(x,y)$, which we already considered.  Fresh variables are supposed to be pairwise distinct, hence  $\Gamma[\vec{a},y][b/x,a/b,c/a] \equiv \Gamma[\vec{a},y]$.
\item
\emph{Second premise:} $\Backward(x,c), R(c,b), R(b,a), \Gamma[a,\vec{a},y] \vdash \Forward(x,y)$. This is an instance  of $\Backward(x,c), R(c,b), \Gamma[b,a,\vec{a},y] \vdash \Forward(x,y)$, which is an instance of $\Backward(x,c), R(c,b), \Gamma[b,\vec{b},y] \vdash \Forward(x,y)$, hence a renaming of the root $\Backward(x,b), R(b,a), \Gamma[a,\vec{a},y]\vdash \Forward(x,y)$. We loop.
\end{enumerate}

We have a notion of global trace condition for proofs with cyclic contexts. We ask that in each infinite path $\pi$, either there is a trace of the inductive definition of some context which is infinitely progressing, or there is a trace of some inductively defined predicate which is infinitely progressing.

We prove the global trace condition for the unfolding of the proof-tree $\Pi$ we just defined. Take any path $\pi$ in $\Pi$. Either $\pi$ chooses the first premise of at least one $\Backward L$-rule, or $\pi$ always chooses the second premise. In the first case the left-hand side of the sequent in $\pi$ reduces by one at each step in $\pi$, therefore $\pi$ is finite. In the second case $\pi$ is the rightmost path in the proof-tree. At each step the unique assumption $\Backward(x,y)$ progresses, therefore if $\pi$ is infinite there is some infinitely progressing trace in $\pi$. 

The same argument saying that cut-free cyclic proofs may only derive the convergence of polynomial-time algorithms works for proofs with cyclic contexts. Instead, proofs with cyclic contexts are not limited to constant space algorithms, because for such proofs the contexts may be of any size. 

In order to add function symbols to the language $L$, we should consider cyclic contexts and terms. We have inductive definitions of the proof-tree, of the contexts, and of the terms. This means that contexts and terms could be infinite, but are not infinite in correct proofs. Also in this case, $2$-Ramsey theorem may be used to prove soundness for this notion of proof. Assume we forbid inductive definitions discarding variables in $L$. Then we conjecture: the set of algorithms we may prove convergent with cyclic contexts and term and cut-free proofs is the set of primitive recursive algorithms.


\end{document}
